% !TEX program = xelatex
\documentclass[11pt,a4paper]{article}

\usepackage{xeCJK}
% Fandol：非可变字体，xdvipdfmx 可正确内嵌
\setCJKmainfont{FandolSong-Regular}[
  BoldFont=FandolHei-Regular, ItalicFont=FandolKai-Regular,
  BoldItalicFont=FandolHei-Regular, Extension=.otf]
\setCJKsansfont{FandolHei-Regular}[Extension=.otf]
\setCJKmonofont{FandolFang-Regular}[Extension=.otf]
\setmainfont{TeX Gyre Pagella}
\setmonofont{DejaVu Sans Mono}[Scale=0.85]

\usepackage{amsmath,amssymb,amsthm}
\usepackage{geometry}
\geometry{margin=2.6cm}
\usepackage{booktabs}
\usepackage{longtable}
\usepackage{xcolor}
\usepackage[colorlinks=true,linkcolor=blue!55!black,citecolor=blue!55!black,
            urlcolor=blue!55!black]{hyperref}
\usepackage{enumitem}
\usepackage{mdframed}
\usepackage[normalem]{ulem}

\linespread{1.15}
\setlength{\emergencystretch}{2em}

\renewcommand{\contentsname}{目录}
\renewcommand{\abstractname}{摘要}
\renewcommand{\figurename}{图}
\renewcommand{\tablename}{表}

\definecolor{axcol}{RGB}{20,80,140}
\definecolor{lncol}{RGB}{100,100,100}

\newtheoremstyle{sudef}{10pt}{10pt}{}{}{\bfseries}{.}{.5em}{}
\theoremstyle{sudef}
\newtheorem{theorem}{定理}[section]
\newtheorem{lemma}[theorem]{引理}
\newtheorem{proposition}[theorem]{命题}
\newtheorem{corollary}[theorem]{推论}
\newtheorem{definition}[theorem]{定义}
\newtheorem{remark}[theorem]{注}

\newenvironment{axiom}[1]
  {\begin{mdframed}[linecolor=axcol,linewidth=0.9pt,backgroundcolor=axcol!4,
    skipabove=10pt,skipbelow=10pt]\noindent\textbf{\color{axcol}公理 #1.}\ }
  {\end{mdframed}}

% Lean 出处标记
\makeatletter
\newcommand{\breakunder}{\discretionary{\char`\_}{}{\char`\_}}
\makeatother
\newcommand{\lean}[1]{\hfill{\small\color{lncol}\begingroup\let\_\breakunder\texttt{#1}\endgroup}}
\newcommand{\leanline}[1]{\par\vspace{2pt}\noindent{\small\color{lncol}\sloppy\let\_\breakunder%
  \setlength{\emergencystretch}{4em}\sloppy Lean 出处：\texttt{#1}\par}}

\newcommand{\dd}{\partial}
\newcommand{\R}{\mathbb{R}}
\newcommand{\N}{\mathbb{N}}
\newcommand{\kron}{\delta}
\newcommand{\Defm}{D}
\renewcommand{\d}{\mathrm{d}}
\newcommand{\ad}{\operatorname{ad}}
\newcommand{\covD}{\nabla}

\title{\textbf{SCD —— 标度耦合动力学}\\[6pt]
\large 内部交流报告：一个无量纲公理体系的当前状态\\[2pt]
\normalsize 全部结论在 Lean 4 中机器验证，无一处 \texttt{sorry}}
\author{形式化验证于 Lean 4 (v4.32.2) + Mathlib}
\date{}

\begin{document}
\maketitle

\begin{abstract}
\noindent
\textbf{这是一份内部交流报告，不是论文。}它记录一个仍在推进的公理体系的当前状态，
包括已被推翻和已被更正的部分——那些与结论同样重要，集中列在第 \ref{sec:audit} 节。

体系的基本命题是：\textbf{时空不是预先给定的几何背景，标度才是基本量}。
由 A3（$\varepsilon\cdot s=1$）局域能动量\emph{就是}局域标度；标度之差表现为曲率；
而“存在什么”是标度轴上的一个截面，不是世界的属性。

\medskip
\noindent\textbf{被导出而非假设的：}标度／转动的拆分、Lorentz 号差、记账形式、
\textbf{标度与转动的耦合}（两次缩放的括号必为转动，由词义强制）、
\textbf{实秩一}从而 Lorentz 代数族、$n=k+1$、
守恒律出自 Bianchi、以及从度规到 $1.7515''$ 与 $42.99''/$世纪的完整引力链条。

\medskip
\noindent\textbf{理论的分叉点}是各向同性／异性二分：
各向同性轨迹上无序参量、无荷、无多重态、无转动；
各向异性轨迹上它们\emph{一并}出现。
而花样是否各向同性是标度相关的事实——这就是“相互作用随能标下降而分离”的导出。

\medskip
\noindent\textbf{当前的记分牌：}一条活张力（DESI 对 $w$ 不演化，$2.8$--$4.2\sigma$，
但 BAO 单独自洽且各联合互相矛盾），一个活的开放问题（$f_0(500)$ 极点落在界的哪一侧），
其余通过。\textbf{光子无色散}这条与所有预设最小长度的纲领分道，且正在赢。

\medskip
\noindent\textbf{边界声明：}Lean 验证的是\emph{从公理到结论的蕴含}，不是公理为真。
第 \ref{sec:scope} 节逐条区分定理与外部输入。

\medskip
\noindent\textbf{形式化规模：}637 条定理、13001 行 Lean 代码、64 个文件；
\textbf{全部 637 条定理}经 \texttt{\#print axioms} 审计
（清单由源码生成而非人工挑选，故无一条能悄悄漏掉），
仅依赖 \texttt{propext}、\texttt{Classical.choice}、\texttt{Quot.sound}，
无 \texttt{sorryAx}。
\end{abstract}

\section*{如何读这份报告}
\addcontentsline{toc}{section}{如何读这份报告}

九个部分按\textbf{理论最终定下的逻辑顺序}排列，不是开发顺序：

\begin{itemize}[leftmargin=1.6em,itemsep=2pt]
\item \textbf{第一部分}——七条公理，全在一处。
\item \textbf{第二部分}——公理逼出什么。单线：地基拆分 $\to$ 号差 $\to$ 记账形式
      $\to$ 方向 $\to$ 曲率即对易子 $\to$ \textbf{耦合（铰链）} $\to$ 实秩一
      $\to$ 维数 $\to$ 轴数。
\item \textbf{第三部分}——\textbf{分叉点}。此后两条路互不相同。
\item \textbf{第四／五部分}——各向同性扇区（几何与引力）／各向异性扇区（标签与荷）。
\item \textbf{第六／七部分}——量子扇区／标度流与内容。
\item \textbf{第八部分}——预言，以及与 DESI／PDG／Fermi-LAT 的对质。
\item \textbf{第九部分}——\textbf{登记册}：撤回、更正、引用未证、登记假设。
\end{itemize}

\begin{mdframed}[linecolor=red!60!black,linewidth=1pt,backgroundcolor=red!3]
\textbf{在采信任何结论之前，请先读第 \ref{sec:audit} 节。}
那里列出被撤回的 4 条、被限缩适用域的 6 条、引用未证的 2 条、登记假设的 2 条，
以及\textbf{同一个“标量代替方向性”错误出现的 5 次}——其中一次挺过了一次形式化和一版报告。
\textbf{没有撤回记录的开发是没有被审计的开发。}
\end{mdframed}

\tableofcontents

\section{导言：把“尺度”而不是“时空”当作基本量}

标准图景中，时空流形与度量先验存在，物质在其上运动，Einstein 方程把二者耦合起来。
这一图景在量子场论面前是尴尬的：QFT 的核心操作——重整化——是关于\emph{能标}的，
而广义相对论的核心对象是\emph{长度}。二者在自然单位制下互为倒数，却被安放在
理论结构的两个完全不同的层次上。

本报告采取相反的次序。设定如下：

\begin{itemize}[leftmargin=1.6em]
\item 基底只有一套\emph{平直的欧氏记账结构}（一张坐标纸），它不携带任何物理；
\item 真正的物理量是每一点上的\emph{局域度量单位}，记其对数为 $\sigma$（稍后将看到它是矢量而非标量）；
\item 长度不是原始概念，而是“沿路径累加局域单位”的结果，于是
      在各向同性情形下 $g=e^{2\sigma}\kron$（第 \ref{sec:frame} 节证明这个特例太强）；
\item 能量标度是长度标度的倒数：$\varepsilon\cdot s=1$，$s=e^{\sigma}$。
\end{itemize}

在这个设定下，“不同点的能动量不同”直接表现为“不同点的标度不同”，
而标度的空间不均匀性\emph{就是}曲率。第 \ref{sec:geom} 节把这句话变成一条精确的恒等式。

\medskip
\noindent\textbf{但那条恒等式不是本理论的核心，它是各向同性扇区的结果。}
开发过程中最重要的一步发生在别处：标度必须是\emph{方向性}的（标量版本被证明太强，
它禁止 Schwarzschild），而一旦标度有方向，\textbf{标度与转动的耦合就被词义强制了}——
缩放是对称算子、转动是反称算子，而两个对称算子的括号必然反称。
沿两个方向各缩放一次，两种次序必然相差一次转动。
这一步（第 \ref{sec:coupling} 节）是全部下游结果的铰链：
它给出实秩一、从而 Lorentz 代数族、从而 $n=k+1$，
也给出量子引力修正的全部内容。

\medskip
\noindent\textbf{理论的分叉点}在第 \ref{sec:locus} 节：
标度花样各向同性时，序参量、荷、多重态、转动\emph{一个都没有}；
各向异性时它们\emph{一并}出现。而花样是否各向同性是标度相关的事实。

\medskip
全部论证在 Lean 4 中完成。之所以坚持零 \texttt{sorry}，是因为一个公理体系一旦
在地基处留下未证之处，其上层的一切“回归”都失去意义。
而之所以坚持记录\emph{撤回}（第 \ref{sec:audit} 节），
是因为一个没有撤回记录的开发，是一个没有被审计的开发。

\section{物理图像：这个世界看起来是什么样}\label{sec:picture}

在写下任何代数之前，先把图像说清楚。\textbf{代数本身不是理论}——
一个只有交换子和 Cartan 分解的陈述，抽象程度并不低于它试图取代的东西。
以下十条是这套公理所描述的世界，用可以想象的语言。

\subsection{没有空间，只有一个比值}

每一"处"有一个数，而这个数是\textbf{比值}：局域度量单位对某个基准的比。
基准的选择不可观测（A5），所以\textbf{只有差有意义}。

\begin{mdframed}[linecolor=axcol,linewidth=1pt,backgroundcolor=axcol!5]
把宇宙想成一张处处标着"这里的一米有多长"的图，而\textbf{图上没有画格子}。
格子是从那些数字之间的\emph{不一致}里长出来的，不是先有格子再往里填数。
\end{mdframed}

\subsection{"两个地方"不是原语，而是被构造出来的}

如果两处的标度花样\textbf{平行}，就没有任何观测量能把它们分开——
不是"看起来一样的两个地方"，而是\textbf{根本没有第二个地方}。

\textbf{完全均匀的区域没有内部结构，连广延都没有。}
距离是"非平行度"累积出来的量。而楔积是反称的，所以
\textbf{这个理论里的分离是一个有向面积，不是一个长度}（第 \ref{sec:native} 节）。

\subsection{标度是有方向的，因为能动量是张量}

每一处的标度不是一个数，是一个\textbf{矢量}。理由很物理：
流动的流体沿流向与横向携带的能量密度不同，
\textbf{张量源不可能只设定一个标量标度}。第 \ref{sec:frame} 节证明标量版本
\emph{禁止} Schwarzschild——它太强，不是太弱。

\subsection{引力是尺子在变，不是空间在弯}

矢量的\textbf{量值}随位置变化，就是引力。曲率恰是标度的二阶非均匀性。

\begin{mdframed}[linecolor=axcol,linewidth=1pt,backgroundcolor=axcol!5]
\textbf{没有任何东西被掰弯。}光走"直线"，但你用来判断"直"的尺子在沿途改变，
于是轨迹看起来是弯的。潮汐力是尺子的\emph{二阶}变化——
一阶（整体变大变小）不可观测，因为那正是 A5 说的基准平移。
\end{mdframed}

\subsection{转动是方向在变——而且它由缩放\emph{生成}}

矢量的\textbf{方向}随位置变化，就是转动。而这里有本理论最硬的一步：

\begin{center}
\textbf{沿两个不同方向各做一次缩放，两种次序的结果必然相差一次转动。}
\end{center}

这不是选来的耦合，是\textbf{词义强制}的：缩放是纯拉伸故对称，转动不改变长度故反称，
而两个对称算子的括号必然反称。物理上这就是 Thomas--Wigner 转动，
而 Terrell--Penrose 效应是它的视觉形式——
高速物体被\emph{看成}转动而非收缩，因为标度改变正在被读作转动。

\textbf{一个方向：平直。两个或更多：曲率，而且是被缩放自己生成的。}

\subsection{$\hbar$ 是标度轴上的一个步长}

不是一个神秘常数。\textbf{沿标度轴平移一步，与"在哪里"这件事不对易}，
而失配的大小就是 $\hbar$（第 \ref{sec:nc} 节，精确非交换，非一阶近似）。

\begin{mdframed}[linecolor=axcol,linewidth=1pt,backgroundcolor=axcol!5]
不确定性不是"测量扰动"，是\textbf{"标度锐利"与"内容丰富"互相排斥}：
完全锐利的标度是一个点，而标度轴上的一个点\textbf{不携带任何内容}。
要分辨 $N$ 个结构，至少得花 $N/\rho$ 宽的窗口。
\end{mdframed}

\subsection{"存在什么"是一个截面，不是世界的属性}

换一个分辨率，你\textbf{不是}对同一批粒子知道得更多，
而是\textbf{在看另外一批}。所以"基本粒子清单"这个问题是不良定的——
第 \ref{sec:emergence} 节证明任何标度之上都还有未分辨的内容，没有最后一层。

\textbf{而标准模型标签之多，度量的是它覆盖的对数范围，不是物理的深度}：
用一套形式覆盖三十个 e-fold，就必须为其中任何地方出现过的东西备好标签，
而绝大多数在任一给定标度上并未实现——这正是壳外内容与规范冗余的形状。

\subsection{粒子是标度矢量场里梳不开的结}

不是场的激发，是\textbf{位形}：矢量绕着某处转而无法被梳平。

\begin{itemize}[leftmargin=1.6em,itemsep=2pt]
\item 它的\textbf{整数荷}是矢量绕球面的圈数；
\item 它的\textbf{质量}不是它"带着"的属性，而是\textbf{它在标度轴上出现的位置}；
\item 它的\textbf{转动标签}是 $v\wedge\nabla v$——所以\textbf{带荷必然带转动}，
      框架禁止"带电而无自旋的基本粒子"（观测上从未见过；带电 $\pi$ 是复合的）。
\end{itemize}

\textbf{共振峰就是阈值}：标度轴上内容改变的那个位置，而峰宽是那个位置有多模糊。

\subsection{宇宙不在膨胀，是标度在跑}

星系并没有在往外飞。\textbf{是尺子在缩}，而由 A3 能量标度是长度标度的倒数，
所以能量标度在耗散——这就是 A6。

\begin{mdframed}[linecolor=axcol,linewidth=1pt,backgroundcolor=axcol!5]
\textbf{没有第一时刻，因为没有零标度。}标度是一个单位，永不为零。
所以"宇宙年龄"这个量不存在，存在的是\textbf{流逝的比值}——有限、可测，但没有起点。
问"大爆炸初值是多少"，等于问一个绝对标度，而 A5 说那不可表达。
\end{mdframed}

\subsection{黑洞任意黑，但不全黑}

视界是 $g_{tt}=0$ 之处，而 $g_{tt}=-s_t^2$，$s_t$ 是单位——\textbf{永不为零}。

\textbf{没有奇点、没有视界、没有无穷红移，三者出自同一个事实：比值可逆。}
黑洞是标度比值\emph{极端但有限}的区域。外部几何原封未动，
所以阴影、ISCO、铃宕全部与广义相对论一致；
差别在于\textbf{没有任何东西被因果隔断，也没有信息被销毁}。

\subsection{两个时间箭头其实是一个序}

一个别处的长期谜题：为什么熵增的方向与宇宙膨胀的方向一致？
通常的回答是"靠一个特殊的初条件对齐"。\textbf{在这里它们不是两件恰好一致的事。}

三条事实，各自因各自的理由成立，彼此互不引用：
第 \ref{sec:signature} 节使\textbf{时间就是标度漂移的方向}；
第 \ref{sec:entropy} 节的 H 定理\textbf{只用有限分辨率}，其证明中\emph{完全没有出现标度}；
第 \ref{sec:horizon} 节使区域的熵为 $\rho\ln R$。

\begin{theorem}[按熵排序与按标度比值排序是同一个序]\label{thm:onearrow}
$\rho\ln R_1 \le \rho\ln R_2 \iff R_1 \le R_2$，严格版本亦然。
\textbf{这是等价，不是单向蕴含}——所以两个箭头是\emph{一个}箭头，
而不是两个恰好同向的箭头。
\leanline{Native.entropy\_order\_is\_scale\_order,
entropy\_strict\_iff\_scale\_strict, arrows\_are\_one\_order}
\end{theorem}

\begin{mdframed}[linecolor=axcol,linewidth=1pt,backgroundcolor=axcol!5]
\textbf{热力学箭头、宇宙学箭头、时间方向，是同一个序被读了三遍。}

\textbf{内容在哪、不在哪：}H 定理确实独立于标度扇区——这是它不循环的原因。
把二者接起来的是"熵 $=\rho\ln R$"，即把熵认作\emph{可分辨阈值的计数}。
\textbf{那个认定是一个建模选择}，由第 \ref{sec:spectrum} 节的阈值密度所激发但并非被迫。
故结论是：\emph{在那个熵读法之下}，两个箭头不是靠边界条件对齐的，它们是同一个关系。
\textbf{削弱那个读法，论证随之削弱。}
\end{mdframed}

\subsection{把十条串起来}

\begin{mdframed}[linecolor=axcol,linewidth=1.2pt,backgroundcolor=axcol!6]
一个数（比值）$\to$ 一个矢量（因为源是张量）$\to$
量值变化给出引力、方向变化给出转动 $\to$
两者的耦合被词义强制 $\to$ 耦合的对易子给出 $\hbar$ 与量子修正 $\to$
矢量梳不开的地方是粒子 $\to$ 矢量整体在跑就是宇宙学。

\textbf{全部物理是同一个矢量场的不同读数。}
"统一"在这里不是把几种力塞进一个群，而是\textbf{它们本来就是一件事的不同侧面}。
\end{mdframed}

\begin{remark}[这张图像的代价]
它买不到的东西必须同样明说：没有耦合常数、没有质量数值、没有代结构、
没有散射振幅。\textbf{这张图像说的是"是什么"，不是"多少"}——
定量的部分只在引力扇区闭合（第四部分），
而粒子扇区给出的是\emph{标签结构}与\emph{界}，不是数值。
\end{remark}

\clearpage

\clearpage
\part{第一部分\quad 公 理}

\section{基底与公理}\label{sec:axioms}

\subsection{微分基底}

为使一切陈述有限、可机器检验，我们不使用解析意义上的光滑函数空间，
而使用它的代数影子。

\begin{definition}[标度代数]
一个 \emph{$n$ 维标度代数} 是一个交换环 $A$，配备 $n$ 个两两交换的导子
$\dd_0,\dots,\dd_{n-1}:A\to A$，满足
\[
\dd_i(a+b)=\dd_i a+\dd_i b,\qquad
\dd_i(ab)=(\dd_i a)b+a(\dd_i b),\qquad
\dd_i\dd_j a=\dd_j\dd_i a .
\]
\leanline{SCD.ScaleAlgebra}
\end{definition}

由 Leibniz 律立刻得到 $\dd_i 1=0$、$\dd_i 0=0$，因而所有整数常量被导子零化
（\texttt{d\_one}、\texttt{d\_zero}、\texttt{d\_natCast}）。
Kronecker 记号 $\kron_{ij}$ 取值于 $A$，是纯记账量，满足 $\dd_k\kron_{ij}=0$。

\begin{remark}
选择代数化而非解析化，使得后文所有曲率恒等式成为\emph{纯代数恒等式}，
对任意交换环成立。这既消除了对分析条件的依赖，也使机器验证不留缝隙。
\end{remark}

\subsection{七条公理}

\begin{axiom}{A1（基底）}
世界的基底是一个环 $A$，配备 $n$ 个满足 Leibniz 律且两两交换的导子 $\dd_i$，
连同平直的 Kronecker 记账 $\kron$。基底本身不携带物理。

\textbf{$A$ 不要求交换。} 交换性是\emph{经典扇区}的性质，不是基底的公理；
第 \ref{sec:nc} 节证明上述微分结构在任意（含非交换）环中由对易子自动实现，
而量子力学正是非交换扇区。第 \ref{sec:geom} 节起的几何推导在交换扇区中进行，
这是一个\emph{物理限制}（引力是经典的），而非数学便利。
\end{axiom}

\begin{axiom}{A2（标度场）}
存在一个无量纲标量 $\sigma\in A$（\emph{对数标度}），及其乘法代表元
$s\in A^{\times}$（$s=e^{\sigma}$），满足指数场的定义性微分方程
\[
\dd_i s = s\,\dd_i\sigma .
\]
$s$ 取值于单位群，故处处非零。
\leanline{SCD.ScaleField}
\end{axiom}

\begin{axiom}{A3（标度—能动量对偶）}
局域能量标度是局域长度标度的倒数：
\[
\varepsilon\cdot s = 1,\qquad \varepsilon := s^{-1}.
\]
在自然单位 $\hbar=c=1$ 下这不是建模选择，而是能量的定义本身。
\leanline{ScaleField.en, ScaleField.en\_mul\_scale}
\end{axiom}

\begin{axiom}{A4（测量，方向性）}
局域度量单位是\emph{每个方向一个}标度 $s_a\in A^{\times}$。被测线元为
\[
g_{ab} = \eta_a\,s_a^{2}\,\kron_{ab}.
\]
\textbf{标度必须有方向}，因为能动量是张量：流动的流体沿流向与横向携带的能量密度
不同，张量源不可能设定一个标量标度。

\emph{各向同性}特例（所有 $s_a$ 相等）给出 $g=e^{2\sigma}\kron$，
这是本报告早期版本所用的公理；第 \ref{sec:frame} 节证明它\emph{太强}
（它禁止 Schwarzschild），并证明它恰是本公理的一个子轨迹，
故一切在其上证明的结果不失。
\leanline{Frame.DirScale, Unify.toMetric\_of\_isotropic}
\end{axiom}

\begin{axiom}{A5（基准协变性）}
$\sigma$ 只在相差一个全局常数的意义下有定义：基准单位的选择不可观测。
形式上，对任意满足 $\dd_i c=0$ 的 $c$，替换 $\sigma\mapsto\sigma+c$ 是理论的对称性。
\leanline{geometry\_fiducial\_invariant}
\end{axiom}

\begin{axiom}{A6（开放性）}
宇宙不是封闭系统：其总能量标度单调耗散。
\leanline{Cosmology.dissipation\_solution}
\end{axiom}

\begin{axiom}{A7（可观测性）}
耦合生成的每一个结构，都携带公理所提供的标签：
$\Lambda^2\mathfrak p$ 与 A4 所标记的 $\mathfrak p$ 同维。
\leanline{Postulates.Observability, Locus.Observability}
\end{axiom}

\begin{mdframed}[linecolor=red!60!black,linewidth=0.9pt,backgroundcolor=red!3]
\textbf{A7 是晚期扶正的，而这是降级不是升级。}
它曾以“判断”之名被携带数版；第 \ref{sec:dimension} 节证明这种形状的条件\emph{不止一条}
且答案不同，故它是一个关于\textbf{什么算可观测}的选择，不是算术。
给定 A7，$k=3$ 是定理，$n=4$ 随之。\textbf{任何人可以拒绝 A7 并得到别的 $k$。}
\end{mdframed}

A5 正是“完全无量纲化”的精确含义，我们将在定理 \ref{thm:fiducial} 中证明它
对全部几何量严格成立。而 A5 的承重远超其外观：
它同时强制耗散速率恒定（第 \ref{sec:cosmos} 节），
故\textbf{DESI 的张力是与 A5 的张力，不是与某个拟合参数的张力}。

\clearpage

\clearpage
\part{第二部分\quad 公理逼出什么}

\section{地基重建：拆分不是一个选择}\label{sec:foundation}

自我审查（第 \ref{sec:audit} 节）指出 A1 白送了太多东西：一组方向、一个维数 $n$、
以及一个随后在干实活的平直度规 $\kron$。更糟的是有两处本该被\emph{导出}的结论被\emph{假设}了：
输运可拆为"标度部分"与"转动部分"；且标度部分是可交换、可积的，而转动部分承载全部曲率。

本节把这两处假设去掉。方法是问：\textbf{在没有坐标纸垫底的情况下，一次标度比较究竟是什么？}

把一切剥掉，剩下的是：局域单位的比较可以复合、复合有结合律、每次比较都可撤销。
那就是一个\textbf{群} $\Gamma$——再多的东西没有，再少又不够。

于是整个拆分由一个观察落地。\textbf{标度是一种相乘的量。}
先按 $a$ 比再按 $b$ 比，必须与先 $b$ 后 $a$ 给出相同的\emph{标度}，因为标度是比值而比值可交换。
所以一次比较所携带的标度，恰是它在 $\Gamma$ 的\emph{最大交换商}——即交换化——中的像。

\begin{theorem}[对易子的标度恒为平凡]\label{thm:commtrivial}
$\mathrm{scale}(aba^{-1}b^{-1})=1$，且反之标度为 $1$ 者恰是对易子子群的元素。
故\textbf{"标度部分"与"转动部分"的拆分不是一个选择}：
转动恰是标度看不见的那部分，而这是被\emph{定义}逼出来的，不是被假设的。
\leanline{Foundation.scale\_commutator, Foundation.scale\_eq\_one\_of\_mem\_commutator,
Foundation.mem\_commutator\_of\_scale\_eq\_one, Foundation.scale\_comm}
\end{theorem}
\textbf{这里没有选择}：交换化是唯一的泛性交换商。

\begin{theorem}[对易子不携带标度]\label{thm:scalecomm}
$\mathrm{scale}(ghg^{-1}h^{-1}) = 1$，无条件成立。
\textbf{不可交换性对标度是隐形的。}
\leanline{Foundation.scale\_commutator}
\end{theorem}

\begin{theorem}[曲率必然对标度隐形]\label{thm:curvinvis}
由第 \ref{sec:conn} 节，曲率\emph{就是}两次输运的对易子；
与定理 \ref{thm:scalecomm} 合起来，曲率的标度必然平凡。

第 \ref{sec:conn} 节的定理 \ref{thm:rotpart} 只在\emph{假设}了拆分之后才证明这一点。
这里拆分不是假设——它是交换化——结论被\textbf{逼出来}。
\leanline{Foundation.curvature\_is\_scale\_invisible}
\end{theorem}

\begin{theorem}[拆分是一个划分，不是一种分解方式]\label{thm:splitcanon}
两次比较携带相同标度，当且仅当它们相差一个纯转动（对易子子群中的元素）。
故"标度部分"与"转动部分"不是按某种约定挑出的两个分量，
而是一个典范商的两侧。
\leanline{Foundation.scale\_eq\_iff\_differ\_by\_rotation}
\end{theorem}

物理读法正是本纲领一直想要、此前却无法辩护的那一句：
\textbf{引力恰好住在标度看不见的地方。}

\section{时间就是标度正在跑的那个方向}\label{sec:signature}

审查把 Lorentz 号差列为无解释的输入。\emph{把它当作输入本身就是惯性}——
在度规为原始量的理论里它才是输入，而这里度规不是原始量。

换个问法：在一个内容只有标度的世界里，\emph{什么可能挑出一个方向}？
候选只有一个。A6 说宇宙耗散：标度在\emph{漂移}。而漂移有方向。

\begin{mdframed}[linecolor=axcol,linewidth=0.9pt,backgroundcolor=axcol!4]
\textbf{时间是标度确实在变的那个方向；空间是标度不变的那些方向。}
\end{mdframed}

\begin{theorem}[号差形状 $1+(n-1)$]\label{thm:sig}
非零漂移的值域是一维的，故由秩—零化度定理其核的余维恰为一。
\textbf{恰好一个方向被挑出，其余全部不被挑出}——这正是 Lorentz 形状，
由耗散得到而非公设。
\leanline{Signature.codim\_ker\_eq\_one}
\end{theorem}

\begin{theorem}[没有耗散就没有时间]\label{thm:notime}
若标度不漂移，则每个方向都是空间的：没有任何方向被挑出，也就根本没有时间方向。
\textbf{时间不是理论被赠予的一个维度——它就是宇宙开放这件事。}
\leanline{Signature.no\_time\_of\_no\_drift,
Signature.spatial\_eq\_top\_of\_no\_drift, Signature.spatial\_ne\_top\_of\_drift}
\end{theorem}

\begin{theorem}[时空拆分与基准无关]\label{thm:sigfid}
把漂移整体重标（这是基准的改变，由 A5 不可观测）不改变空间子空间。
故时空的划分是基准无关的，正如任何物理陈述都必须如此。
\leanline{Signature.spatial\_smul, Signature.drift\_direction\_unique}
\end{theorem}

号差因此不是钦定的 $(-,+,+,+)$。它是 $1+(n-1)$，
因为漂移是\emph{一个}线性泛函；而它之所以存在，只是因为宇宙是开放的。
\textbf{A5 与 A6 联手做了原先由公设 $\eta$ 做的事。}

\begin{remark}[仍未导出]
维数 $n$ 本身，以及使那个被挑出的方向为\emph{类时}（而非仅仅"被挑出"）的符号约定，
仍是输入。
\end{remark}

\section{记账度规也不是公设}\label{sec:invariant}

审查的第一条最要命：A1 宣称 Kronecker $\kron$"不携带物理"，
然后 \texttt{Conformal.lean} 用它升降指标、定义给出 Ricci 的缩并。
一个平直度规被白送进来，随后在干实活。

第 \ref{sec:foundation} 节已说明拆分不是选择。把同一个问题对 $\kron$ 问一遍：
\textbf{比较代数上有没有什么典范的东西可以充当记账形式？}

有，而且它需要的比度规少得多。设代数带一个\emph{迹}——一个满足
$\tau(xy)=\tau(yx)$ 的线性泛函。这比内积弱得多：无正定性、无非退化性、
无基的选择；它在任何矩阵代数上都是典范的，而且这就够了。

\begin{theorem}[记账形式由迹给出且对称]\label{thm:formsymm}
$\kappa(x,y):=\tau(xy)$ 仅凭迹性质即为对称。
\leanline{Invariant.Trace.form\_symm}
\end{theorem}

\begin{theorem}[不变性——$\kron$ 真正在提供的那条性质]\label{thm:forminv}
\[
\kappa([z,x],y)+\kappa(x,[z,y])=0 .
\]
那些其对易子\emph{就是}曲率的输运（第 \ref{sec:conn} 节）保持这个形式。
这正是被公设的平直度规在 \texttt{Conformal.lean} 中所扮演的角色，
而在这里它是关于任意迹的一条\emph{定理}。
\leanline{Invariant.Trace.form\_invariant, Invariant.Trace.form\_ad\_skew}
\end{theorem}

\begin{mdframed}[linecolor=axcol,linewidth=0.9pt,backgroundcolor=axcol!4]
三个各自独立的输入——一组方向、一个平直度规 $\kron$、一个号差 $\eta$——
现在压缩成一个：\textbf{一个带迹的比较代数}。
号差出自耗散（第 \ref{sec:signature} 节），拆分出自交换化（第 \ref{sec:foundation} 节），
形式出自迹（本节）。
\end{mdframed}

\begin{remark}[压缩不是推导]
\emph{维数}仍然是那个代数的维数，仍是输入。本节不假装相反。
\end{remark}

\section{方向不是给定的：方向就是代数}\label{sec:direction}

最后一个未解决的地基输入是维数 $n$，以及 A1 附赠的指标集 $\mathrm{Fin}\,n$。
本节把指标集去掉。而在过程中发现，A1 夹带的\emph{不止}一个维数。

看 A1 到底说了什么：有 $n$ 个导子 $\dd_0,\dots,\dd_{n-1}$，它们\textbf{两两交换}。
这里被假设的是两件事，不是一件。指标集是显眼的那件，另一件是\emph{交换性}。

\textbf{更正（见第 \ref{sec:dynamics} 节）。}我原先把那个交换性读作一个隐藏的
\emph{平直性}假设——"相互交换的方向就是平直的坐标方向"。
\textbf{那个解读是错的，在此撤回。}相互交换的坐标方向只是选了一张图，
局部永远可选，与曲率无关；
定理 \ref{thm:noflatcomm} 把导子全部杀死——它们能有的最强交换性——曲率原封不动。
下面的定理 \ref{thm:expose} 本身成立（它带着平直前提），
但从它得出的那个解读不成立。A1 假设的是一张图和一个指标集，
这是真实的输入，但比我当时宣称的要小。

把指标集丢掉。方向不是标签，而是比较代数的元素，输运对它是线性的。于是：

\begin{theorem}[曲率即输运不尊重括号的程度]\label{thm:dircurv}
\[
F(x,y) = [D_x,D_y] - D_{[x,y]} .
\]
没有任何外部的东西参与定义——无指标集、无度规、无号差。
（对比第 \ref{sec:conn} 节的 $F$，那里需要一个联络一形式和一个平直 $\kron$ 来升降指标。）
\leanline{Direction.DirTransport.curv, curv\_antisymm}
\end{theorem}

\begin{theorem}[平直 $\iff$ 输运是 Lie 映射]\label{thm:flatlie}
输运无曲率，当且仅当它把方向的括号送到输运的对易子。
曲率是这件事的\emph{唯一}障碍。
\leanline{Direction.DirTransport.flat\_iff\_lie\_hom}
\end{theorem}

\begin{theorem}[旧理论是交换情形]\label{thm:abcase}
当方向代数交换时，曲率退化为纯对易子 $[D_x,D_y]$——
即第 \ref{sec:conn} 节所计算的对象。
\leanline{Direction.DirTransport.curv\_of\_abelian}
\end{theorem}

\begin{mdframed}[linecolor=gray!60!black,linewidth=0.9pt,backgroundcolor=gray!4]
\begin{theorem}[交换性与括号的关系]\label{thm:expose}
对平直输运，"诸方向相互交换"（A1 的要求）\emph{等价于}"括号平凡地作用"。

注意前提：这是关于\emph{平直}输运的陈述，
\textbf{不是}"相互交换的方向必平直"的推导。见第 \ref{sec:dynamics} 节的更正。
\leanline{Direction.DirTransport.commute\_iff\_bracket\_acts\_trivially}
\end{theorem}
\end{mdframed}

于是 $n$ 不再被手工挑选。选基是便利而非公理
（\texttt{transport\_determined\_by\_generators}）。

\begin{mdframed}[linecolor=red!60!black,linewidth=0.9pt,backgroundcolor=red!3]
\textbf{更正（见第 \ref{sec:dimension} 节）。}我曾写"方向\emph{就是}代数，
其个数就是代数的维数"。代数被定名之后这句话可以核验，而它是\textbf{假的}：
$\dim\mathfrak{so}(3,1)=6$，而时空有 $4$ 个方向。
方向是代数所\textbf{作用}的空间——它的定义表示——不是代数本身。
但什么也没丢：$\mathfrak{so}(k,1)$ 带有一个典范的 $k+1$ 维表示，
故定名代数仍然定住了 $n$，即 $k+1$。
\end{mdframed}

\begin{remark}[诚实的界限]
本节导出的是"方向\emph{是}代数元素"以及"其个数\emph{是}代数的维数"。
它\emph{没有}导出是\emph{哪个}代数，因而不预言 $n=4$。
输入从"$n$ 个交换标签 $+$ 平直度规 $+$ 号差"变成"一个代数"；
这是压缩，不是维数的推导。
\end{remark}

\section{曲率就是对易子}\label{sec:conn}

\subsection{避开一次惯性}

补各向异性曲率的"显然"办法，是写下 Cartan 结构方程然后硬算。
那等于把广义相对论的整套机器原样搬过来——正是最该避开的路径依赖。
有一条更本原的路，而且它把引力认成了本报告早已造好的东西。

问：输运局域单位\emph{究竟}是什么？沿方向 $i$ 移动，就是把一切改用邻近的单位表述，
那是一个算子 $D_i$。先 $i$ 后 $j$ 未必等于先 $j$ 后 $i$。这个差额就是曲率的全部：
\[
F_{ij} = [D_i,\;D_j].
\]

这与 \texttt{Quantum.ad} 是\emph{同一个构造}。那里的输运是相空间平移，
它们的不对易就是 $\hbar$；这里的输运是标架输运，它们的不对易就是引力。

\begin{mdframed}[linecolor=axcol,linewidth=0.9pt,backgroundcolor=axcol!4]
\textbf{曲率与 $\hbar$ 是同一个构造在两族输运上的取值。}
这比"两个理论之间存在类比"强得多——它们是一件事。
\end{mdframed}

\subsection{Yang--Mills 曲率不是被假设的，是被算出来的}

\begin{theorem}[曲率即输运的对易子]\label{thm:commcov}
以 $D_i m = \dd_i m + [A_i,m]$ 为协变输运，则
\[
[D_i,D_j]\,m \;=\; [F_{ij},\,m],
\qquad
F_{ij} = \dd_iA_j-\dd_jA_i+[A_i,A_j].
\]
两件事同时发生：其一，Yang--Mills 曲率不是被写下的，而是左边算出来的结果；
其二，结果对 $m$ 是\emph{代数}作用——$m$ 的导数全部消失——
这正是"曲率是张量"的精确含义。
\leanline{Connection.comm\_covD}
\end{theorem}

\begin{theorem}[Bianchi 恒等式]\label{thm:bianchi}
$\;\sum_{\text{cyc}} D_iF_{jk}=0$。它由 \texttt{ad} 的 Jacobi 恒等式
（定理 \ref{thm:adjac}）与混合偏导可交换直接给出。
\textbf{引力场的守恒律就是输运的结合律}，别无其他内容。
\leanline{Connection.bianchi}
\end{theorem}

\subsection{究竟是哪一部分在弯曲}

\begin{theorem}[梯度联络是平的]\label{thm:gradflat}
若联络是标量的梯度且自对易，其曲率恒为零。
这是定理 \ref{thm:weyl} 在非交换情形下的重述：
输运的\emph{标度}部分是可积的，故无论几何多弯，第二时钟效应都不存在。
\leanline{Connection.curv\_gradient\_eq\_zero}
\end{theorem}

\begin{theorem}[全部曲率都住在不对易的部分]\label{thm:rotpart}
把输运拆为标度部分（由 A5 是梯度，故对易、可积）加标架转动部分：
\[
F\bigl(\dd\sigma+\omega\bigr)_{ij} \;=\; F(\omega)_{ij}.
\]
\textbf{标度部分对曲率的贡献恰好为零。}
\leanline{Connection.curv\_eq\_rotation\_part, curv\_split, curv\_pure\_scale\_eq\_zero}
\end{theorem}

于是引力不是"标度场把空间掰弯了"。引力是\textbf{标架转动的不对易}，
而标度本身可积地输运。这既解释了可积 Weyl 几何为何能承载引力却不产生第二时钟效应，
也正是标量公理 A4 表达不了的那个各向异性曲率——而且全程没有出现一个 Christoffel 记号。

推论：纯标量输运的曲率恒为零，这就是旧公理禁止 Schwarzschild 的结构性原因，
与定理 \ref{thm:frameflat} 从度量一侧给出的结论一致。

\begin{remark}[仍未闭合的一环]
本节给出的是\emph{给定联络}的曲率。由标架\emph{定出}联络
（Cartan 第一结构方程，即无挠条件 $de^a+\omega^a{}_b\wedge e^b=0$）尚未形式化。
因此各向异性扇区现在有了曲率机器与拆分定理，但还没有"从 $s_a$ 算到 $F$"的完整链条，
光线偏折与近日点进动仍未导出。缺口比上一版窄，但没有消失。
\end{remark}

\section{耦合一直都在：两次缩放生成一次转动}\label{sec:coupling}

第 \ref{sec:expressive} 节把框架的弱点诊断为"标度路与转动路按构造独立"，
并断言进展需要一个公理不提供的\emph{通路间耦合}。

那个诊断是通过分析 \texttt{Foundation.scale} 得到的，
而它把标度取为交换化，理由是"标度是比值，比值可交换"。
\textbf{那个理由说的是标量标度。} A4 给的不是标量标度，是\emph{每方向一个}标度。
\textbf{而沿不同方向的缩放并不对易。}

这是同一个替换错误第\textbf{三}次出现——在公理是方向性的地方用标量推理——
前两次是第 \ref{sec:frame} 节（几何）与第 \ref{sec:axis} 节（缺陷）。
这次它在地基上，而改正它\emph{恰好}供出了那个被说成缺失的耦合。

\subsection{耦合由词义强制}

缩放是纯拉伸：它改变沿其轴的长度而不转动任何东西，故作为算子是\textbf{对称}的。
转动不改变任何长度，故是\textbf{反对称}的。于是：

\begin{theorem}[两次缩放的括号是转动]\label{thm:ss}
若 $S,T$ 对称，则 $[S,T]$ \textbf{反对称}。
即：沿一个方向缩放后再沿另一个方向缩放，与反序执行的结果\textbf{相差一次转动}。
没有任何东西是被选的；生成元的对称性就逼出了它。
\leanline{Coupling.bracket\_scaling\_scaling}
\end{theorem}

\begin{theorem}[完整结构]\label{thm:cartan}
$[S,S]\subseteq A$、$[A,S]\subseteq S$、$[A,A]\subseteq A$。
\textbf{标度扇区不是子代数}，而它不闭合的那部分\emph{恰好落在转动扇区}。
\leanline{Coupling.bracket\_rotation\_scaling, bracket\_rotation\_rotation,
coupling\_structure}
\end{theorem}

\begin{theorem}[具体地，在三维中]\label{thm:concrete}
$0$--$1$ 面与 $0$--$2$ 面上的两个方向性缩放，其括号\emph{恰好}是
$1$--$2$ 面上的转动：$[K_1,K_2]=J$，逐分量算出，不诉诸任何群论。
故这两个缩放\textbf{不对易}——标度扇区一旦具有方向性就不再交换。
\leanline{Coupling.boost\_bracket\_eq\_rotation, scalings\_do\_not\_commute}
\end{theorem}

\subsection{一个方向什么也生成不了}

\begin{theorem}[单方向 $\Rightarrow$ 平直]\label{thm:onedir}
单个缩放与自身的括号为零，故不生成转动。
\leanline{Coupling.bracket\_self, no\_rotation\_of\_commuting}
\end{theorem}

这与已证的一切自洽：定理 \ref{thm:rotpart} 说纯标度联络是平的；
而\emph{单个}缩放方向正是没有括号可用、因而不生成转动的情形。
定理 \ref{thm:frameflat} 是同一事实的度量侧表述。

\begin{mdframed}[linecolor=axcol,linewidth=0.9pt,backgroundcolor=axcol!4]
\textbf{一个方向：平直。两个或更多：曲率，而且是被缩放自己生成的。}
\end{mdframed}

\subsection{这改变了什么}

\begin{enumerate}[leftmargin=1.8em]
\item \texttt{Foundation.scale} 作为交换化，描述的只是\emph{标量}情形。
      方向性标度下，缩放不构成交换商——\textbf{它们生成转动}。
\item 定理 \ref{thm:indepchan}（两路独立）因此是关于\emph{标量扇区}的陈述，
      不是关于本框架的。\textbf{两路是耦合的}，而且是最强的耦合：
      转动扇区被标度扇区\emph{生成}。
\item 第 \ref{sec:crosscheck} 节找不到跨扇区检验，原因也随之改写：
      不是"按构造不可能"，而是我当时分析的是退化情形。
\end{enumerate}

\begin{remark}[物理上这是什么]
这就是 Thomas--Wigner 转动：两个不同方向的推促复合出一个推促\emph{加一次转动}。
Terrell--Penrose 效应是它的视觉形式——高速运动的物体被看成\emph{转动}而非收缩，
因为标度改变正被读成转动。\textbf{两者都不是从外面引进来的}；
它们就是两个对称生成元的括号。
\end{remark}

\section{到底是哪个代数：公理逼出实秩一}\label{sec:algebra}

第 \ref{sec:direction} 节把输入从"$n$ 个交换标签 $+$ 平直度规 $+$ 号差"压到"一个代数"，
然后停在那里：它没说是\emph{哪个}代数。这是此后最尖锐的未决项，
而下游全在等它——交叉积、$\beta$ 函数的结构、量子扇区的任何定量陈述。

结果是：\textbf{两条早已证过、且当初并非为此而证的定理，合起来把候选压进了一份有限的分类清单。}

\subsection{分裂是典范的，不是表示的选择}

第 \ref{sec:coupling} 节把缩放定义为对称算子、转动定义为反称算子。
转置是相对某个形式取的，所以这看起来像偷偷塞进了一个度规。它没有：
第 \ref{sec:invariant} 节已从一个迹导出形式 $\kappa(x,y)=\tau(xy)$。

\begin{theorem}[两扇区在导出的形式下正交]\label{thm:orth}
$S$ 对称、$A$ 反称 $\Rightarrow$ $\operatorname{tr}(SA)=0$，
仅用 $\operatorname{tr}(X^{\mathsf T})=\operatorname{tr}X$ 与
$\operatorname{tr}(XY)=\operatorname{tr}(YX)$。
故标度／转动分解在框架\emph{自己拥有}的形式下正交——
连同定理 \ref{thm:cartan} 的括号关系，这恰是一个\textbf{Cartan 分解}
$\mathfrak g = \mathfrak k \oplus \mathfrak p$。
\leanline{Algebra.scaling\_rotation\_orthogonal}
\end{theorem}

\subsection{公理逼出实秩一}

两条分别为别的目的证过的结论：

\begin{itemize}[leftmargin=1.6em]
\item 定理 \ref{thm:onedir}：相互对易的缩放不生成转动。
      故一族两两对易的缩放，就是一组可以\textbf{同时指定而不产生曲率}的标度方向；
      这样的极大族即 $\mathfrak p$ 的极大交换子代数，其维数按定义就是\textbf{实秩}。
\item 定理 \ref{thm:sig}：漂移是\emph{一个}线性泛函。标度只沿\textbf{一个}方向跑，
      故只有\textbf{一个}标度坐标。
\end{itemize}

\begin{mdframed}[linecolor=axcol,linewidth=1pt,backgroundcolor=axcol!5]
\[
\boxed{\ \text{实秩} = 1\ }
\]
这是框架\emph{对自己}施加的约束，不是为了让什么对上而选的。
\end{mdframed}

\begin{theorem}[两个推促的括号是空间转动]\label{thm:bbr}
$[K_v,K_w]$ 在空间块上是 $v_iw_j-v_jw_i$，其余分量为零。
这就是定理 \ref{thm:ss} 在 Lorentz 族上的显式形式，也就是 Thomas--Wigner 转动。
\leanline{Algebra.boost\_bracket\_spatial}
\end{theorem}

\begin{theorem}[两个推促对易 $\iff$ 方向平行]\label{thm:rank1}
故\textbf{不存在二维的相互对易缩放族}：缩放扇区的极大交换子代数是一条\emph{直线}。
实秩一是算出来的，不是断言的。
\leanline{Algebra.boosts\_commute\_iff\_parallel,
no\_two\_dimensional\_commuting\_family, commuting\_boosts\_lie\_on\_a\_line}
\end{theorem}

\begin{theorem}[两个"秩一"是同一句话]\label{thm:samefact}
一族对易缩放所诱导的漂移两两成比例。
故第 \ref{sec:signature} 节的"一个漂移泛函"与耦合的"一个对易缩放"，
是 Cartan 分解两侧对同一事实的两次读数。
\leanline{Algebra.commuting\_drifts\_proportional, boost\_drift}
\end{theorem}

\subsection{这剩下什么}

实秩一的实单李代数是被完全分类的，清单很短：
\[
\mathfrak{so}(n,1),\quad \mathfrak{su}(n,1),\quad \mathfrak{sp}(n,1),\quad \mathfrak{f}_{4(-20)} .
\]
它们的对称空间分别是实、复、四元数、八元数双曲空间。

\begin{mdframed}[linecolor=red!60!black,linewidth=0.9pt,backgroundcolor=red!3]
\textbf{这条分类是李理论的定理，不是本报告的推导，此处\emph{未被证明}。}
它是引用，并已登记进第 \ref{sec:audit} 节的污染表作为外部输入。
本报告证明的是\emph{筛掉}这份清单的那条约束，以及 $\mathfrak{so}(n,1)$ 满足它。
\end{mdframed}

四者之中，$\mathfrak{so}(n,1)$ 是标度为\textbf{实}比值的那一个；
其余三者要求标度带复、四元数或八元数结构，而 A3 不提供——
那里的标度是量值之比，且第 \ref{sec:axis} 节已发现读数是无向的实数。
\textbf{于是框架指向 Lorentz 族，而它是从耗散加耦合走到那里的。}

\subsection{三个等着的项目随之落地}

\begin{enumerate}[leftmargin=1.8em]
\item \textbf{交叉积。}秩一意味着 Iwasawa 分解的交换部分是一维的，
      故在交叉积中作用的群是 $\mathbb R$ 而非 $\mathbb R^k$。
      \textbf{这正是第 \ref{sec:exact} 小节单独一个平移 $U$ 就够用的原因}——
      那本来看起来只是建模上的省事。
      \leanline{Algebra.commuting\_family\_one\_parameter, boost\_add, boost\_smul}
\item \textbf{$\beta$ 函数。}秩一根系至多有两个正根 $\alpha,2\alpha$，
      故至多\textbf{两个}标度权重。（随分类一并引用。）
\item \textbf{双曲几何，未被塞入。}秩一对称空间就是双曲空间，
      其 Lorentz 形式即反德西特，标度坐标是径向方向，重整化流沿它跑。
      \textbf{这不是放进去的，这就是秩一本身。}
\end{enumerate}

\begin{remark}[诚实的界限]
这导出的是\emph{族}，不是维数：$n=4$ 仍未被预言，这里没有任何东西挑出它。
分类是引用来的。而从"一个漂移泛函"到"一维极大交换子代数"这一步是一个
\emph{物理认定}——认定独立的对易缩放就是独立的标度坐标——不是形式推演；
它被明写在外面，没有藏进任何证明里（在 Lorentz 模型\emph{内部}它是定理 \ref{thm:samefact}）。
\end{remark}

\section{最后两项输入究竟是什么}\label{sec:dimension}

剩下两项：$n=4$ 靠一个判断，标度量值自由。
去追它们之间的关系时，翻出了夹在中间、且一直被我混淆的第三件事，故先说它。

\subsection{方向是表示，不是代数}

第 \ref{sec:direction} 节去掉指标集时宣称：\textbf{方向就是代数，其个数就是代数的维数}。
代数定名之后这句话可核验，而它\textbf{失败了}。

\begin{theorem}[代数严格大于方向集]\label{thm:algbig}
$\dim\mathfrak{so}(3,1) = 6 > 4 = n$。
故方向\textbf{不是}代数，而是代数所作用的空间——定义表示。
\leanline{Dimension.algebra\_bigger\_than\_directions}
\end{theorem}

更正不损失什么，反而有所得。$\mathfrak{so}(k,1)$ 带一个典范的 $k+1$ 维定义表示。

\begin{mdframed}[linecolor=axcol,linewidth=1pt,backgroundcolor=axcol!5]
\textbf{故代数一旦定名，$n$ 就不再是一项额外输入——它是 $k+1$，
剩下要定的只有 $k$。}\textbf{两项输入合并为一项}，
这是"两个未决项是否相关"这个问题的前一半答案。
\leanline{Dimension.rep\_dim\_determined}
\end{mdframed}

\subsection{两条同形的自匹配条件，答案不同}

第 \ref{sec:slice} 节用的条件是：耦合生成的转动与 A4 所标记的方向同样多，
\[
\dim\Lambda^2\mathfrak p = \dim\mathfrak p \iff \frac{k(k-1)}{2}=k \iff k=3 .
\]
把方向与代数区分开之后，\emph{第二条}同形条件浮现出来——要求代数与它作用的对象一样大：
\[
\dim\mathfrak{so}(k,1) = k+1 \iff \frac{k(k+1)}{2}=k+1 \iff k=2 .
\]

\begin{theorem}[两条条件给出不同答案]\label{thm:twocond}
$k=3$ 满足前者而不满足后者；$k=2$ 满足后者而不满足前者。
\leanline{Dimension.rot\_matches\_scale\_iff, alg\_matches\_rep\_iff,
two\_matching\_conditions\_differ}
\end{theorem}

\textbf{所以"要求结构与自身匹配"在说清\emph{哪个}指标集须匹配\emph{哪个}之前是不良定的。}
把这一点明说出来是本小节的要点：数值巧合恰恰藏在这种缝隙里。

框架确实能在两者间决断：A4 按\textbf{方向}指派标度，而那些方向是缩放扇区 $\mathfrak p$；
第 \ref{sec:gauge} 节使一切可观测标签都是它们之间的花样。
故可用的标签是 $\mathfrak p$ 的，而需要标签的是耦合生成的 $\Lambda^2\mathfrak p$。
第二条比较的是代数与其表示，而\textbf{框架并不要求这两者相等}——
Lorentz 代数是六维而时空是四维，这不构成任何缺陷。

于是 $k=3$ 成立，$n=4$ 随之，而它所依赖的判断现在被尖锐地陈述为：
\emph{A4 提供的标签必须足以标记耦合生成的转动}。
这是\textbf{一条关于可观测性的假设}，不是数值偏好。

\subsection{自由的量值不是缺陷}

标度花样是一个矢量；方向是转动规范，量值未定。这看起来像第二项松动的输入。它不是同一类东西。

\begin{theorem}[一切都对花样齐次，故比值不变]\label{thm:ratio}
$\text{wedge}(cv, cw) = c^2\,\text{wedge}(v,w)$，
故一切转动标签的\emph{比值}在重标度下\textbf{不变}。
\leanline{Dimension.wedge\_smul\_both, wedge\_ratio\_invariant,
only\_ratios\_are\_fixed}
\end{theorem}

而比值正是框架所预言的东西——第 \ref{sec:spectrum} 节的换算表、
第 \ref{sec:particle} 节的离散度、引力那些数——没有一个会动。

\begin{mdframed}[linecolor=axcol,linewidth=1pt,backgroundcolor=axcol!5]
\textbf{量值不是一项缺失的预言：它是预言被表达于其中的\emph{单位}。}
一个无量纲理论必须恰好携带一个这样的数，而本框架恰好携带一个——\textbf{这是可能的最少}。
理论把一切定到一个总体因子，而定不了那个因子，是因为它\emph{不可能}定——
再没有别的东西可以用来表达它。
\end{mdframed}

\subsection{两项自由在何处相遇}

$k$ 离散、量值连续，故它们不是同一项自由。但它们由一条换算相连，
而这正是"$n$ 应当与标度有关"这个直觉成立的意义。

秩一对称空间在径向坐标下体积按 $e^{(k-1)r}$ 增长，而径向坐标\emph{就是}标度坐标
（第 \ref{sec:algebra} 节）。故在该体积中均匀分布的群体，其阈值密度是每单位径向 $k-1$，
而实测密度是每单位对数能量；两者差一个根的归一化 $\alpha$，而那正是自由的量值：
\[
\rho \cdot \alpha = k - 1 .
\]

\begin{theorem}[密度--归一化换算]\label{thm:dennorm}
$k=3$ 时 $\rho\alpha = 2$；代回即返回输入。
\leanline{Dimension.density\_normalization\_relation, normalization\_returns\_input}
\end{theorem}

\begin{mdframed}[linecolor=red!60!black,linewidth=0.9pt,backgroundcolor=red!3]
\textbf{这是换算，不是预言}——一个方程一个未知数，正是第 \ref{sec:crosscheck} 节
记录过的陷阱。它\emph{确实}确立的是：\textbf{两项剩余自由并不独立}——
定住 $k$ 就把 $\rho$ 的测量变成对归一化的确定，反之亦然。

它还额外依赖"缺陷在对称空间体积中均匀分布"这一条，而本报告没有任何定理提供它。
\textbf{已登记为假设，未被夹带。}
\end{mdframed}

\begin{remark}[输入清单现状]
代数族——已导出（实秩一，加一条引用的分类）；
$n$——不再独立，代数定名后即 $k+1$；
$k=3$——一条判断，已在上文尖锐陈述；
量值——不是同种意义的输入，它是单位，而无量纲理论恰需一个。
\end{remark}

\section{两难消解，而且更多轴不可得}\label{sec:axes}

第 \ref{sec:slice} 节把一个两难称作框架最尖锐的未决问题。
\textbf{那个两难建立在一处误读上，而误读是我的。}

定理 \ref{thm:uniax} 证的是 $\text{wedge}\,v\,v=0$——一个矢量与\emph{自身}的楔积。
我把它读成"单轴花样不携带转动"。但\emph{单轴}指标度只取一个不同的\textbf{值}，
不是只有一个\textbf{方向}；而定理讲的是一个\emph{矢量}，不是一根轴。
$\text{wedge}\,v\,v=0$ 实际说的是：\textbf{均匀}花样不生成转动。
那是关于恒定性的陈述，不是关于轴数的。

\textbf{这是"标量代替方向性"的第五次}——继第 \ref{sec:frame}、\ref{sec:axis}、
\ref{sec:coupling} 节与第 \ref{sec:dynamics} 节撤回的那处解读之后。
值得记下的是：这次它\textbf{挺过了一次形式化和一版论文}——定理本身是真的，
错的只是它的名字和我的注解。

\subsection{转动是花样与其自身变化的楔积}

定理 \ref{thm:binj} 说标度花样是\emph{单个}矢量 $v$。
故在任一处 $v\wedge v = 0$，转动标签只能来自花样在\textbf{不同位置取不同值}。

\begin{theorem}[转动需要变化，且只需要变化]\label{thm:rotvar}
非零转动标签强制花样的两个值不平行；反之，值处处平行的花样处处不生成转动。
\leanline{Axes.rotation\_needs\_variation, combable\_no\_rotation,
rotation\_iff\_variation}
\end{theorem}

\begin{mdframed}[linecolor=axcol,linewidth=1pt,backgroundcolor=axcol!5]
\textbf{角动量在这里是 $v\wedge\nabla v$：标度矢量与其自身变化的楔积。}
它\textbf{不是}一个位形在某点所具有的属性。
\end{mdframed}

而这恰好\emph{有利于}框架地消解了两难，因为缺陷\emph{正是}花样并非处处平行的位形——
这就是绕数的含义：

\begin{theorem}[绕数缺陷必然携带转动]\label{thm:wound}
取两个不平行值的花样有非零转动标签；
反之，无转动的位形可梳，故拓扑平凡。
\leanline{Axes.wound\_carries\_rotation, rotationless\_is\_trivial}
\end{theorem}

\textbf{框架禁止带拓扑荷而转动标签为零的缺陷。}
观测上，从未见过基本的带电自旋零粒子——带电 $\pi$ 介子是复合的。
作为吻合陈述，并附带保留：此处"基本"本身就是标度相关的（第 \ref{sec:emergence} 节）。

\subsection{更多轴不是不受青睐，而是不可得}

要求是：没有充分论证不得排除双轴乃至更多轴。论证如下，而且很短。

\textbf{"单轴／双轴"是\emph{张量}序参量的词汇}——
向列相有一个对称张量，其不同本征值的个数就是轴数。
\textbf{本框架的序参量不是张量。}实秩一强制标度花样是单个\textbf{矢量}
（定理 \ref{thm:rank1}、\ref{thm:binj}），
而矢量只有一个量值和一个方向，别无其他。
\textbf{没有第二个本征值可以与第一个不同。}

\begin{theorem}[任何第二个缩放要么是同一根轴，要么是一次变化]\label{thm:noax2}
给定花样 $v$，任何另一个缩放 $w$ 要么是 $v$ 的倍数（同一根轴），
要么与之有非零楔积——那它就不是第二根轴，而是一次\emph{变化}，表现为转动。
\textbf{没有第三种可能。}故"有几根轴"的答案是\textbf{一根}，且非选择所致。
\leanline{Axes.no\_second\_axis, axis\_or\_variation}
\end{theorem}

\begin{theorem}[花样的稳定化子是正维数的]\label{thm:stabpos}
一个非平凡转动生成元湮灭该花样。这正是\emph{单轴}序参量的标志。
\leanline{Axes.axis\_stabilizer\_nontrivial}
\end{theorem}

而这一点让整件事扣上了。\textbf{连续}稳定化子正是使序参量空间有非平凡 $\pi_2$、
从而有点缺陷的东西；双轴或完全各向异性的花样有\textbf{离散}稳定化子，
而李群模离散子群的 $\pi_2=0$，故\textbf{根本没有点缺陷，也没有整数荷}。

\begin{mdframed}[linecolor=axcol,linewidth=1pt,backgroundcolor=axcol!5]
\begin{center}
单轴 $\longrightarrow$ 连续稳定化子 $\longrightarrow$ 点缺陷、整数荷\\[2pt]
双轴及以上 $\longrightarrow$ 离散稳定化子 $\longrightarrow$ 无点缺陷、无 $\mathbb Z$
\end{center}
\textbf{第 \ref{sec:particle} 节的整数荷只在单轴情形存在，而实秩一使单轴成为唯一情形。}
两件事互相印证，这正是那组标签成立的原因。
\end{mdframed}

（$\pi_2(S^n)$ 与"李群模离散子群的 $\pi_2$ 为零"是标准拓扑，
与秩一分类同等地属于引用。此处证明的是代数那一半：
花样是单个矢量，且其稳定化子正维数。）

\subsection{这在别处的代价}

第 \ref{sec:dynamics} 节把定理注解为"多重态内部平直"。
它的 \texttt{Realizes} 前提\textbf{按方向}指派缩放算子，而实秩一不提供这个——
它提供一个矢量、一个缩放。定理对满足其前提者仍成立，
但\textbf{那条物理解读撤回}。取代它的是本节：转动由单一花样的\emph{变化}携带。
\leanline{Axes.multiplet\_flatness\_needs\_per\_direction\_scalings}

\clearpage

\clearpage
\part{第三部分\quad 分 叉 点}

\section{A7，以及立场转变所迫的一次审计}\label{sec:locus}

两件事在此了结。第一件是一条一直以"判断"之名在干活的公设，此处把它扶正并命名。
第二件是一次审计——若不做，一个反驳我便无法回答。

\subsection{A7（可观测性），作为公设}

第 \ref{sec:dimension} 节把最后一项数值输入压到一步：
\emph{A4 提供的标签必须足以标记耦合生成的转动}。
它同时表明这种形状的条件\emph{不止一条}且彼此不合，
所以这一步是关于\textbf{什么算可观测}的\emph{选择}，不是算术。

\textbf{这是一项哲学承诺，而且没有从其余七条公理推导出它的前景}——
它就是每个物理理论都必须在某处对"什么算一次观测"所做的那种承诺。故予以扶正：

\begin{mdframed}[linecolor=axcol,linewidth=1.2pt,backgroundcolor=axcol!6]
\textbf{A7（可观测性）。}耦合生成的每一个结构，都携带公理所提供的标签。
\end{mdframed}

\begin{theorem}[给定 A7，$k=3$ 是定理]\label{thm:a7}
且 $n=4$ 随之，因为定理 \ref{thm:algbig} 已经使 $n$ 不再是独立输入。
而 A7 是一条真实的限制：它在其他任何维数上都不成立，故非平凡。
\leanline{Locus.three\_from\_observability, observability\_fails\_elsewhere,
observability\_iff\_three}
\end{theorem}

\textbf{这是对一项主张的降级，不是对一项结果的升级}：
原先被呈现为"几乎已导出"的东西，现在公开地是第七条公理。
框架的输入是 A1--A7 加一个单位，这是诚实的记账。

\subsection{反驳}

本纲领的立场发生过转变：它从试图复现每一个粒子开始，
终于第 \ref{sec:emergence} 节"物种清单是某一截面的性质"。反驳是：
先前那些面向相互作用的结果是在旧立场下写的，可能已与新立场不一致——

\begin{quote}
假如有任何基本相互作用是标度不变的，那就等于承认存在标度不变的粒子及其性质；
而且这与"相互作用只在能标下降后才分离"的图像矛盾。
\end{quote}

\textbf{这个反驳在形式上是对的}，而审计有一个伤亡和一项实质收获。

\textbf{存活的部分。}开发中每一项面向相互作用的结果，
都是关于某个量\emph{如何随标度变化}的陈述，而非关于某个标度不变的相互作用：
第 \ref{sec:running} 节的跑动、第 \ref{sec:qcd} 节的渐近自由与量纲嬗变、
第 \ref{sec:particles} 节的动量依赖寿命、暗物质的探测失效。
\textbf{没有一条断言相互作用在标度上是任何固定的东西。}

\begin{theorem}[没有标度不变的相互作用可供承认]\label{thm:noinv}
耦合在标度上恒定，当且仅当其缺陷密度为零——即无物可分辨的退化情形，
而不是一种基本力。
\leanline{Locus.no\_scale\_invariant\_interaction,
constant\_coupling\_iff\_no\_content}
\end{theorem}

\textbf{伤亡。}第 \ref{sec:charges}、\ref{sec:particle} 节给缺陷指派拓扑标签，
而同伦类\emph{就是}标度不变的。这与第 \ref{sec:emergence} 节存在真实张力，
也正是该反驳的真正靶心。

\subsection{解药一直在文件里}

第 \ref{sec:axis} 节：完全简并的标度花样什么也不区分，\textbf{完全不支持任何缺陷}。
第 \ref{sec:gauge} 节：其稳定化子是全部，故无任何多重态可分辨。
而下面这条：两个各向同性花样恒平行，故也不生成任何转动。

\begin{theorem}[各向同性轨迹上没有任何标签]\label{thm:iso}
各向同性 $\Rightarrow$ 无转动标签、无区分方向；
连同上述两条，\textbf{各向同性轨迹不携带任何种类的标签}。
而非零转动标签\textbf{强制}花样各向异性。
\leanline{Locus.isotropic\_no\_rotation, labels\_require\_anisotropy,
everything\_switches\_on\_together, particle\_locus\_is\_anisotropic}
\end{theorem}

\begin{mdframed}[linecolor=axcol,linewidth=1pt,backgroundcolor=axcol!5]
\begin{center}
各向同性花样 $\longrightarrow$ 无序参量、无荷、无多重态、无转动：什么也不被区分\\[3pt]
各向异性花样 $\longrightarrow$ 序参量出现，荷、多重态与曲率随之一并出现
\end{center}
\textbf{而花样是否各向同性，是一个标度相关的事实。}
\end{mdframed}

于是那些拓扑标签\textbf{不是}世界的标度不变性质。
它们是\emph{各向异性轨迹}的性质，在各向同性破缺的那个标度上诞生——
而且是\textbf{一并}诞生的，因为它们同源。

\begin{mdframed}[linecolor=axcol,linewidth=1pt,backgroundcolor=axcol!5]
\textbf{这就是"相互作用随能标下降而分离"，是导出的而非假设的}；
而第 \ref{sec:emergence} 节没有被违反：某标度上存在什么，仍然取决于该标度。
\end{mdframed}

\textbf{适用域更正备案。}第 \ref{sec:charges}、\ref{sec:particle} 节应被读作描述
\emph{各向异性轨迹}。定理不变，其定义域现在被写明。
第 \ref{sec:frame} 节早已把各向同性情形降格为度规扇区的一个轨迹——
同一降格现在也适用于荷，这是我此前没注意到的一处自洽。

\section{耦合让标度场做什么}\label{sec:dynamics}

第 \ref{sec:coupling} 节给出了耦合。本节问它让理论\emph{说}什么——
关于地基、关于演化、关于谱。三者是同一个问题，故在同一处回答。

\subsection{一处更正，以及交换环真正的代价}

\begin{theorem}[杀死导子并不能压平曲率]\label{thm:noflatcomm}
令所有导子为零——方向能有的最强交换性——曲率仍是 $\ad(A_i)(A_j)$，完全自由。
而第 \ref{sec:coupling} 节的两个方向性缩放给出
$\ad K_1 K_2 = J \neq 0$，是一个现成的反例。
\textbf{平直性住在联络对不对易里，不在方向对不对易里。}
\leanline{Dynamics.no\_flatness\_from\_commuting\_derivations, nonflat\_witness}
\end{theorem}

于是交换的\emph{函数}代数并不是对量子非对易性的限制。
空间上的函数本就交换，那是经典且正确的；非对易性属于比较算子，
而 \texttt{Foundation}、\texttt{Connection}、\texttt{Crossed}、\texttt{Quantum}、
\texttt{Coupling} 全都保持它们非交换。

\textbf{但交换扇区确实有代价，而这里把它精确定住。}

\begin{theorem}[交换扇区的输运恒对易]\label{thm:commcost}
在交换环上 $\ad$ 恒为零，故 $[\covD_i,\covD_j]=0$：
无论联络是什么，该扇区\textbf{完全没有输运意义下的曲率}。
\leanline{Dynamics.commutative\_sector\_transports\_commute,
commutative\_curvature\_undetectable}
\end{theorem}

这正是第 \ref{sec:geom} 节非公设 $\kron$ 不可的原因：
它的曲率只能从共形因子进来，进不来输运。
\textbf{代价就是这个，而且被限制在经典各向同性几何里，不是对量子扇区的限制。}

\subsection{演化方程就是 Bianchi，其内容是守恒}

\begin{theorem}[Bianchi 强制源守恒]\label{thm:conserv}
若标度场的曲率正比于一个源——任何场方程的形状——
则该源的循环协变散度\emph{恒等于零}。
\leanline{Dynamics.source\_conserved\_of\_field\_equation}
\end{theorem}

\begin{theorem}[尖锐形式：不守恒的源没有场方程]\label{thm:nofe}
若循环散度在某处不为零，则该源\textbf{不可能}是任何联络的曲率之源。
本框架不是偏好守恒的源，而是\emph{无法表述}不守恒的源。
\leanline{Dynamics.no\_field\_equation\_of\_nonconserved}
\end{theorem}

\begin{mdframed}[linecolor=axcol,linewidth=0.9pt,backgroundcolor=axcol!4]
\textbf{能动量守恒在本框架中不是额外公设，而是场方程可被写下的条件。}
它经由 Bianchi 来自 Jacobi 恒等式——来自输运的结合律，别无其他。
\end{mdframed}

再加上耦合的贡献：$F$ 反称，而由定理 \ref{thm:ss} 它取值于转动。
故\textbf{标度场的源是一个转动密度}。
\leanline{Dynamics.source\_antisymm, field\_equation\_content}

\subsection{多重态内部是平的}

这里是耦合真正付账的地方。第 \ref{sec:gauge} 节说多重态是方向性标度的简并块
（$a\sim b \iff s_a = s_b$）；第 \ref{sec:coupling} 节说缩放与自身的括号为零。

\begin{theorem}[同一多重态内不生成转动]\label{thm:multflat}
同一多重态中的两个方向携带同一单位，故有同一缩放，故括号为零。
\textbf{标度的简并块内部是平的；曲率只在多重态\emph{之间}生成。}
\leanline{Dynamics.same\_multiplet\_no\_rotation}
\end{theorem}

\begin{theorem}[逆命题：转动是简并破缺的见证]\label{thm:rotwit}
两方向之间有非零转动，就证明它们属于\textbf{不同}的多重态。
\leanline{Dynamics.rotation\_implies\_different\_multiplet}
\end{theorem}

在耦合之前，规范多重态与曲率是标度花样的两个互不相干的读数；
\textbf{现在前者控制后者}。而全简并（各向同性）的极端情形给出：
所有括号为零，处处不生成转动——\textbf{最大对称的位形就是平的}。
这从规范一侧第三次独立地回收了定理 \ref{thm:frameflat} 与 \ref{thm:rotpart}。
\leanline{Dynamics.larger\_multiplet\_fewer\_generators}

\begin{theorem}[多重态只能在阈值处改变]\label{thm:multthresh}
内部对称性完全由标度的重合花样决定，故花样不变则多重态不变。
\textbf{新的多重态不会在两个阈值之间出现}，
与第 \ref{sec:emergence} 节的内容变化发生在同一批位置。
\leanline{Dynamics.multiplets\_fixed\_by\_pattern,
multiplet\_change\_needs\_pattern\_change}
\end{theorem}

\begin{remark}[诚实的界限]
以上都没有预言\emph{哪些}多重态。它说的是它们可以在何处改变、以及它们内部是平的。
标度花样仍然是输入。
\end{remark}

\clearpage

\clearpage
\part{第四部分\quad 各向同性扇区：几何与引力}

\begin{mdframed}[linecolor=axcol,linewidth=1.2pt,backgroundcolor=axcol!6]
\textbf{先解一个表面矛盾。}第三部分说各向同性轨迹上\textbf{无转动、无荷、无多重态}，
而本部分整部在各向同性扇区\emph{计算曲率}。这不矛盾，但必须说清：

\begin{center}
\textbf{“各向同性”指标度矢量的\emph{方向}处处一致，不是它的\emph{量值}不变。}
\end{center}

方向一致 $\Rightarrow$ 楔积恒为零 $\Rightarrow$ 无转动标签、无荷、无多重态。
但量值仍可随位置变化 $\Rightarrow$ 共形因子非常数 $\Rightarrow$
\textbf{几何仍然弯曲}。这正是第 \ref{sec:cosmos} 节据以导出暗物质的那一条：
\emph{只引力，无其他}。

\medskip
\noindent
所以本部分的定量结果——$1.7515''$、$42.99''/$世纪、$\gamma=1$——
是\textbf{纯引力扇区}的结果：那里没有标签可算，也不需要。
而第 \ref{sec:frame} 节证明把各向同性当\emph{公理}太强（它禁止 Schwarzschild），
故它是一个\emph{轨迹}，不是世界。
\end{mdframed}

\section{几何的涌现：曲率即标度形变}\label{sec:geom}

\begin{mdframed}[linecolor=axcol,linewidth=1pt,backgroundcolor=axcol!5]
\textbf{本节全部结果属于\emph{各向同性扇区}。}
A4 给的是\emph{每方向一个}标度；本节所用的 $g=e^{2\sigma}\kron$ 是它的
各向同性特例（所有 $s_a$ 相等），第 \ref{sec:frame} 节证明把这个特例当作\emph{公理}
太强——\textbf{它禁止 Schwarzschild}。特例被降格为一个\emph{轨迹}，
故本节的定理不失，但它们的适用域是那个轨迹。
\end{mdframed}

\subsection{Levi-Civita 数据}

在各向同性轨迹上，度量 $g=e^{2\sigma}\kron$ 的 Christoffel 记号为
\begin{equation}
\Gamma^{a}{}_{bc}
= \kron_{ab}\,\sigma_c + \kron_{ac}\,\sigma_b - \kron_{bc}\,\sigma_a,
\qquad \sigma_i:=\dd_i\sigma .
\end{equation}
\leanline{SCD.Chr}

值得注意：共形因子 $e^{2\sigma}$ \emph{完全消失}。联络只看得见标度的\emph{变化}，
看不见标度本身——这已经预示了 A5。

\begin{definition}[标度形变张量]
\begin{equation}
\Defm_{ij} := 2\sigma_{ij} - 2\sigma_i\sigma_j + \kron_{ij}\,|\nabla\sigma|^2,
\end{equation}
其中 $\sigma_{ij}=\dd_i\dd_j\sigma$，$|\nabla\sigma|^2=\sum_k\sigma_k\sigma_k$。
它度量局域单位\emph{偏离仿射变化}的程度。（因子 $2$ 是为避免除法，
使一切陈述在任意交换环上成立。）
\leanline{SCD.Defm}
\end{definition}

\subsection{主定理}

\begin{theorem}[曲率即标度形变]\label{thm:master}
在任意标度代数中，对任意 $\sigma\in A$ 与任意指标 $a,b,c,e$，
\begin{equation}\boxed{\;
2\,R^{a}{}_{bce}
= \kron_{ae}\Defm_{bc} - \kron_{ac}\Defm_{be}
+ \kron_{bc}\Defm_{ae} - \kron_{be}\Defm_{ac}\;}
\end{equation}
即 $g=e^{2\sigma}\kron$ 的完整 Riemann 张量恰是欧氏记账 $\kron$ 与标度形变 $\Defm$
的 Kulkarni–Nomizu 乘积。
\leanline{SCD.two\_Rm\_eq}
\end{theorem}

这是\emph{各向同性扇区}的核心恒等式（不是全部纲领的核心；那是第 \ref{sec:coupling}
节的耦合）。它把“欧氏空间在不同点按不同比例缩放，等价于黎曼几何”
从一句口号变成一条恒等式：\textbf{几何 $=$ 记账 $\times$ 标度形变，没有别的成分}。
曲率没有任何独立于标度场的自由度。

\begin{proof}[证明要点]
Riemann 张量的导数部分给出
\[
\dd_c\Gamma^a{}_{be}-\dd_e\Gamma^a{}_{bc}
=\kron_{ae}\sigma_{bc}-\kron_{be}\sigma_{ac}-\kron_{ac}\sigma_{be}+\kron_{bc}\sigma_{ae},
\]
其中 $\kron_{ab}$ 项因混合偏导可交换而相消（\texttt{d\_Chr\_diff}）。
二次部分由一条主收缩恒等式给出：
\[
\sum_m \Gamma^{a}{}_{cm}\Gamma^{m}{}_{be}
= 2\kron_{ac}\sigma_b\sigma_e-\kron_{ac}\kron_{be}|\nabla\sigma|^2
+\kron_{ab}\sigma_c\sigma_e+\kron_{ae}\sigma_b\sigma_c
-\kron_{bc}\sigma_a\sigma_e-\kron_{ce}\sigma_a\sigma_b
\]
（\texttt{sum\_Chr\_Chr\_full}；Lean 中通过把九个乘积项整理成可收缩形式，
再用 $\sum_a\kron_{ab}f(a)=f(b)$ 逐项化简）。两部分相加，$|\nabla\sigma|^2$ 项
恰好成对相消，余下即所断言的形式。全部步骤为环上恒等式，由 \texttt{ring} 封闭。
\end{proof}

\subsection{迹与标量曲率}

\begin{theorem}[Ricci 曲率]\label{thm:ric}
\begin{equation}
R_{be} = -(n-2)\bigl(\sigma_{be}-\sigma_b\sigma_e\bigr)
- \kron_{be}\bigl(\Delta\sigma + (n-2)|\nabla\sigma|^2\bigr).
\end{equation}
\leanline{SCD.Ric\_eq}
\end{theorem}

\begin{theorem}[标量曲率]\label{thm:rsc}
\begin{equation}
e^{2\sigma}R = -2(n-1)\,\Delta\sigma - (n-1)(n-2)\,|\nabla\sigma|^2 .
\end{equation}
\leanline{SCD.RscBare\_eq}
\end{theorem}

（我们把共形因子显式保留在左侧，从而无需任何可逆性假设。）

\begin{corollary}[真空即标度的仿射变化]\label{cor:flat}
若 $2$ 在 $A$ 中可逆（对任何 $\R$-代数模型自动成立），则
\[
\Defm_{ij}\equiv 0 \;\Longrightarrow\; R^{a}{}_{bce}\equiv 0 .
\]
即：形变自由的标度场给出\emph{完全平直}的几何，而非仅仅 Ricci 平坦。
\leanline{Rm\_eq\_zero\_of\_Defm\_eq\_zero}
\end{corollary}

真空因此不是“没有标度”，而是“标度均匀地变化”。这与本纲领的直觉一致：
标度处处存在，引力只是它的不均匀性。

\subsection{可积 Weyl 几何：没有第二时钟效应}

Weyl 于 1918 年提出：单位长度可以逐点变化，由一个独立的一形式 $A_i$ 描述。
Einstein 的反驳是致命的：若 $A$ 不可积，原子谱线将依赖于世界线历史
（“第二时钟效应”），与观测矛盾。

本理论中 $A_i$ \emph{不是}独立场，而是全局标量的梯度 $A_i=\dd_i\sigma$。

\begin{theorem}[标度联络平坦]\label{thm:weyl}
\[
F_{ij} := \dd_i A_j - \dd_j A_i = 0 \quad\text{恒成立}.
\]
\leanline{SCD.scaleCurv\_eq\_zero}
\end{theorem}

证明只用到混合偏导可交换。因此沿任意闭合回路输运单位后单位不变，
第二时钟效应\emph{按构造}不存在。这正是可积 Weyl 几何成功、
而 Weyl 原版失败的原因——本纲领所需的恰是前者。

\subsection{无量纲性是一条定理}

\begin{theorem}[基准协变性]\label{thm:fiducial}
设 $\dd_i c=0$。则对 $\sigma\mapsto\sigma+c$，下列量逐一不变：
\[
\sigma_i,\quad \sigma_{ij},\quad \Delta\sigma,\quad |\nabla\sigma|^2,\quad
\Gamma^a{}_{bc},\quad R^a{}_{bce},\quad R_{be},\quad \Defm_{ij},\quad e^{2\sigma}R .
\]
\leanline{SCD.geometry\_fiducial\_invariant}
\end{theorem}

这就是 A5 的精确内容，也是“全部物理量无量纲化”的严格实现：
没有任何可观测量依赖于 $\sigma$ 的零点位置，只有标度\emph{差}是物理的。

\begin{remark}[重整化群的起源]
定理 \ref{thm:fiducial} 说的是\emph{几何}对基准平移不变。但\emph{物质的描述}
未必不变：为在基准移动时保持预言不变，耦合常数必须相应流动。
这一补偿运动就是重整化群（第 \ref{sec:rg} 节）。
换言之，RG 不是 QFT 的技术手续，而是 A5 在物质部门的影子。
\end{remark}

\section{A4 太强：标度必须是方向性的}\label{sec:frame}

\subsection{诊断}

第 \ref{sec:geom} 节的一切都建立在 $g=e^{2\sigma}\kron$ —— \emph{一个}标量标度，
各方向相同。本报告早期版本把这当作 A4 本身。它无法承载 Schwarzschild，
而原因是物理的，不是技术的。

本纲领说局域能动量决定局域标度。但\textbf{能动量是张量}：
流动的流体沿流向与横向携带的能量密度不同。张量源不可能设定一个标量标度。
\textbf{标度必须有方向。}

这就是第 \ref{sec:axioms} 节把 A4 直接写成\emph{方向性}的理由：
局域度量单位是每个方向一个标度 $s_a\in A^{\times}$，被测线元为
$g_{ab}=\eta_a s_a^2\kron_{ab}$。各向同性（所有 $s_a$ 相同）是它的特例。

\subsection{旧公理不是没能给出 Schwarzschild，而是禁止了它}

Schwarzschild（及 Reissner–Nordström）族的定义性关系是 $g_{tt}g_{rr}=-1$。
在标度语言中它读作

\begin{mdframed}[linecolor=axcol,linewidth=0.9pt,backgroundcolor=axcol!4]
\[
s_t\cdot s_r=1 .
\]
\textbf{时间标度与径向标度互为倒数}：钟被压缩的因子，恰好等于径向尺子被拉长的因子。
这是一句关于\emph{标度}的陈述，标量理论连表述它的能力都没有。
\end{mdframed}

\begin{theorem}[各向同性 $+$ 倒数关系 $\Rightarrow$ 平直]\label{thm:frameflat}
若所有方向共用一个标度，且 $s_ts_r=1$，则 $s_t^2=1$，
于是每个度量分量都等于裸记账值：几何是平直的。
\leanline{Frame.DirScale.isotropic\_reciprocal\_forces\_flat,
isotropic\_reciprocal\_metric\_bare}
\end{theorem}

这就是诊断的精确形式：\textbf{$g=e^{2\sigma}\kron$ 并非未能产生 Schwarzschild，
而是排除了它}。公理确实太强。

\begin{theorem}[方向性 A4 下该关系有非平凡解]\label{thm:frameexists}
对任意满足 $u^2\ne1$ 的单位 $u$，存在方向性标度场使 $s_t=u$、$s_r=u^{-1}$，
倒数关系成立且\emph{非}各向同性。
\leanline{Frame.exists\_reciprocal\_nonisotropic}
\end{theorem}

\subsection{一套公理，而不是两套}

若把方向性 A4 与旧的标量 A4 并列，体系就同时挂着两条互斥的公理——
\texttt{Conformal.lean} 等文件用标量 $\sigma$，而 \texttt{Frame.lean} 说标度有方向。
这必须消除，消除的办法是\textbf{降格旧公理}：方向性 A4 是公理，标量形式是关于其
各向同性轨迹的定理。

\begin{theorem}[旧公理是新公理的各向同性轨迹]\label{thm:unify}
若所有方向共用一个标度 $s$，则方向性 A4 的度量\emph{恰好}等于共形度量
$\eta_a s^2\kron_{ab}$——无近似、无附加假设。
\leanline{Unify.toMetric\_of\_isotropic}
\end{theorem}

\begin{theorem}[各向同性轨迹恰是标量理论的像]\label{thm:unifyiff}
一个方向性标度场是各向同性的，当且仅当它形如 $\mathrm{ofScalar}(u)$。
故旧公理不只是被新公理蕴含，而是新公理限制在一个子轨迹上的\emph{全部}内容，
没有剩余。
\leanline{Unify.isotropic\_iff\_ofScalar, Unify.sector\_coherence}
\end{theorem}

于是此前一切关于 $\sigma$ 的定理原封不动保留，只是获得了明确的适用域：
它们是关于\emph{各向同性扇区}的陈述。公理表由两条变回一条。

\begin{remark}[相容但尚不完备]
本报告\emph{未}完成\emph{各向异性}扇区的曲率计算（Cartan 结构方程、
Schwarzschild 曲率张量的逐分量验证、光线偏折与近日点进动）。
公理体系现在是\textbf{相容}的；它还不\textbf{完备}——
方向性 A4 的几何推论目前只在标度恰为各向同性处被算出。
这是当前最优先的下一步。
\end{remark}

\section{引力与规范场是同一个东西}\label{sec:transport}

从标架定出联络，是否需要无挠条件？\textbf{这样问本身就是惯性。}
“挠率”来自一个把标架与联络当作\emph{独立场}的理论，所以才需要一个条件把二者拴住。
在这里它们并不独立，也就没有什么可强加的。

从 $x$ 到 $y$ 的输运只有一件事要做：把 $x$ 处每个局域单位送到 $y$ 处对应的单位。
两端的标度花样就是全部，于是输运就是"使单位对上"的那个重标号。
\textbf{没有做出任何选择，因而不需要任何条件来消除选择。}

\begin{theorem}[非简并 $\Rightarrow$ 输运唯一]\label{thm:tunique}
若目标处各局域单位\emph{互不相同}，则输运\textbf{唯一}。
没有无挠条件、没有度规相容性公设、没有变分原理：
联络是标度场的一个\emph{函数}，它从来就不是一个待约束的独立场。
\leanline{Transport.transport\_unique\_of\_nondegenerate}
\end{theorem}

\begin{theorem}[简并 $\Rightarrow$ 自由度恰是内禀对称性]\label{thm:tcoset}
两个输运同一花样的重标号相差一个 $\mathrm{stabilizer}$ 元素，且仅此而已；
反之乘上任一稳定子元素仍是输运。故输运集合恰是稳定子的一个陪集。
\leanline{Transport.transport\_coset, Transport.transports\_of\_stabilizer}
\end{theorem}

\begin{theorem}[唯一性 $\iff$ 无规范自由]\label{thm:tiff}
输运唯一，当且仅当稳定子平凡。
\leanline{Transport.transport\_unique\_iff\_stabilizer\_trivial,
Transport.nondegenerate\_no\_gauge\_freedom}
\end{theorem}

\begin{mdframed}[linecolor=axcol,linewidth=0.9pt,backgroundcolor=axcol!4]
\textbf{输运是一个对象；标度花样的简并把它劈成"被定死的部分"与"自由的部分"。}
被定死的部分逐点变化并弯曲——那是\emph{引力}；
自由的部分是块内转动——那是\emph{规范场}。
它们不是两个有待统一的场，而是一个场，
被"局域单位是否恰好重合"这一件事分开。
\end{mdframed}

\section{引力红移与 Newton 极限}\label{sec:newton}

\subsection{引力红移}

\begin{theorem}[红移]\label{thm:redshift}
对 A2、A3 的标度场，
\[
\dd_i \varepsilon = -\,\varepsilon\,\dd_i\sigma .
\]
\leanline{SCD.ScaleField.d\_en}
\end{theorem}

证明只是对 $\varepsilon s=1$ 求导。物理读法：标度\emph{升高}处能量标度\emph{降低}。
一个固定激发从标度大的区域向外传播时能量下降——引力红移是定理，不是附加假设。

\subsection{标度响应律}

要得到引力，还需要一句话说明\emph{什么在源生标度}：

\begin{mdframed}[linecolor=red!60!black,linewidth=1pt,backgroundcolor=red!3]
\textbf{A6$'$（标度响应）不是本体系的公理。}公理只有 A1--A7（第 \ref{sec:axioms} 节）。
下式是一条\textbf{外部物理输入}，本报告\emph{未}从更基本的东西导出它，
已登记在第 \ref{sec:scope} 节：
\[
e^{2\sigma}R = \kappa\rho,
\]
其中 $\rho$ 是局域无量纲能量密度；$n=3$ 时这正是共形平直切片上的 Hamilton 约束。
\textbf{凡依赖它的结论（本节的 Newton 极限）都继承这一输入。}
\end{mdframed}

\subsection{一阶区域}

\begin{theorem}[线性化是精确的]\label{thm:linear}
若 $|\nabla\sigma|^2=0$，则
\[
e^{2\sigma}R = -2(n-1)\Delta\sigma .
\]
\leanline{SCD.RscBare\_linear}
\end{theorem}

\begin{theorem}[Poisson 方程]\label{thm:poisson}
在一阶区域，令 $\Phi:=-\sigma$（\emph{Newton 势就是负的对数标度}），则
\[
2(n-1)\,\Delta\Phi = \kappa\rho .
\]
$n=3$ 时为 $4\Delta\Phi=\kappa\rho$；取 $\kappa=16\pi G$ 即得
\[
\Delta\Phi = 4\pi G\rho .
\]
\leanline{SCD.poisson, SCD.poisson\_three}
\end{theorem}

符号自洽性值得强调：$\sigma=-\Phi$ 意味着引力势阱深处 $\sigma>0$，
即\emph{标度被拉大}、尺子被拉长——这与 Schwarzschild 各向同性坐标下
径向固有距离大于坐标距离一致；同时由定理 \ref{thm:redshift}，
$\varepsilon=e^{-\sigma}=e^{\Phi}$，正是引力红移因子。两个符号都不是调出来的。

\subsection{一阶区域不是“扔掉高阶项”}

条件 $|\nabla\sigma|^2=0$ 看似是把困难项抹去。我们证明它是\emph{精确可实现}的：
把基底扩张为对偶数环 $A[\varepsilon]/(\varepsilon^2)$，并取无穷小标度场
$\sigma=\varepsilon\varphi$。

\begin{theorem}[一阶区域的实现]\label{thm:dual}
对偶数环上的导子逐分量延拓后仍是标度代数；且对任意 $\varphi\in A$，
\[
|\nabla(\varepsilon\varphi)|^2 = 0 \quad\text{恒等成立}.
\]
因而定理 \ref{thm:linear}、\ref{thm:poisson} 在其上\emph{精确}成立，
没有任何被丢弃的项。
\leanline{SCD.Dual.ScaleAlgebra, Dual.gradsq\_inr, Dual.poisson\_exact}
\end{theorem}

一阶微扰论由此被安置在公理体系\emph{内部}，而不是从外部强加。

\section{使倒数关系得以成立的那次相消}\label{sec:cancel}

第 \ref{sec:vacuum} 节由两件事导出 $s_t\cdot s_r=1$：真空使标度和为常数，
以及渐近平直性定住那个常数。后者已证；前者被标注为输入——
"标准的 Ricci 组合"——它是引力链条上最后一个未证的环节。
本节证明其中\emph{带内容}的那部分。

把度规写成 A4 的方向性形式，只有时间与径向单位随 $r$ 变，
$A=e^{2\sigma_t}$、$B=e^{2\sigma_r}$。相关的两个 Ricci 分量各自都很乱：
二阶导数、二次项、看不出结构。但取那个特定组合 $R_{tt}/A+R_{rr}/B$，
\textbf{所有难看的东西全部相消}：

\begin{theorem}[相消]\label{thm:cancel}
\[
\frac{R_{tt}}{A}+\frac{R_{rr}}{B} \;=\; \frac{1}{rB}\Bigl(\frac{A'}{A}+\frac{B'}{B}\Bigr).
\]
二阶导数彼此抵消，二次项也彼此抵消——对\emph{一切} $A,B$ 精确成立。
剩下的是一个全对数导数。
\leanline{Schwarzschild.ricci\_combination}
\end{theorem}

这就是真空何以对 $AB$ 有话可说的原因，也是倒数关系存在的全部理由。

\begin{theorem}[真空 $\Rightarrow$ 标度和驻定]\label{thm:vacstat}
$R_{tt}=R_{rr}=0$ 强制 $A'/A+B'/B=0$（因前因子 $1/rB$ 永不为零），
在标度变量下即 $\sigma_t'+\sigma_r'=0$——
这恰是第 \ref{sec:vacuum} 节所需的那条假设。
\leanline{Schwarzschild.vacuum\_log\_derivative,
Schwarzschild.vacuum\_iff\_scale\_sum\_stationary}
\end{theorem}

\begin{mdframed}[linecolor=axcol,linewidth=0.9pt,backgroundcolor=axcol!4]
相消的那个组合恰是按 $t$--$r$ 块的逆度规加权的组合。
所以真空方程限制在该块上只说了一件事：\textbf{那里的面元是驻定的}；
它对时间标度与径向标度\emph{各自}如何分布只字未提。
这正是"真空中引力只是拿一个换另一个"的精确含义。
\leanline{Schwarzschild.area\_element\_stationary}
\end{mdframed}

\paragraph{链条现状。} 引力扇区现在跑通为：
Ricci 分量（输入）$\to$ 相消（\emph{本节已证}）$\to$ 标度和为常数（已证）
$\to$ 渐近平直性（已证，第 \ref{sec:vacuum} 节）$\to$ $s_ts_r=1$
$\to$ $\gamma=1$（已证，第 \ref{sec:ppn} 节）$\to$ 偏折 $1.7515''$ 与进动 $42.99''$（已证）。

仍为输入的只剩两个 Ricci 表达式本身，它们来自对角拟设的 Levi-Civita 联络；
那一步是机械的，而\emph{本节证明的这一步才是有非平凡事情发生的地方}。

\section{引力扇区闭合}\label{sec:riccidiag}

第 \ref{sec:cancel} 节证明了相消，但把两个 Ricci 分量当作给定。本节把它们导出来，
引力扇区从此从度规一路跑到实测数字，中间不留假设。

对角度规 $-A(r)dt^2+B(r)dr^2+r^2d\Omega^2$ 的联络系数由 Levi-Civita 公式给出，
对角情形下每个符号都是一行计算。由它们按
$R_{ab}=\dd_c\Gamma^c_{ab}-\dd_a\Gamma^c_{cb}+\Gamma^c_{cd}\Gamma^d_{ab}-\Gamma^c_{ad}\Gamma^d_{cb}$
求和：

\begin{theorem}[两个 Ricci 分量，已导出]\label{thm:ricciderived}
求和的结果\emph{恰好}是第 \ref{sec:cancel} 节所假设的两个表达式。
\leanline{RicciDiag.ricciTT\_eq, RicciDiag.ricciRR\_eq}
\end{theorem}

\begin{mdframed}[linecolor=axcol,linewidth=0.9pt,backgroundcolor=axcol!4]
\textbf{引力链条端到端。} 度规 $\to$ 联络 $\to$ Ricci（\emph{已证}）$\to$ 相消（已证）
$\to$ 标度和驻定（已证）$\to$ 渐近平直性（已证）$\to$ $s_ts_r=1$
$\to$ $\gamma=1$（已证）$\to$ 偏折 $1.7515''$ 与进动 $42.99''$（已证）。

\textbf{从度规到那两个数字之间，已无任何一步是假设。}
\leanline{RicciDiag.gravity\_chain}
\end{mdframed}

扇区内仍为输入的只剩：标度响应律 A6$'$，以及静态球对称拟设的选取。引力波未涉及。

\section{倒数关系的那个常数从哪来}\label{sec:vacuum}

第 \ref{sec:ppn} 节让倒数条件 $s_ts_r=1$ 补上了偏折的另一半，
但把倒数性本身留成了"被描述的解的性质"。本节关掉这个缺口中\emph{能}关掉的那半。

真空条件并不直接给出 $s_ts_r=1$。它给出更弱的东西：
时间与径向对数标度之\emph{和}的导数为零，故 $s_ts_r$ 是\emph{常数}。
从"常数"到"一"是另一步，而那一步才有物理内容：它是\textbf{渐近平直性}——
远离源处什么也没有，故局域单位必须与基准一致。

\begin{theorem}[渐近平直性定住倒数关系]\label{thm:recip}
若 $\sigma_t+\sigma_r$ 导数为零且在远处为零，则它处处为零，即 $s_t\cdot s_r=1$。
那个常数不是拟合的：它由公理允许的唯一边界条件定住——
由 A5，"远处"恰恰就是基准的意思。
\leanline{Vacuum.reciprocity\_of\_asymptotic\_flatness, Vacuum.scale\_product\_eq\_one}
\end{theorem}

\begin{theorem}[边界条件在干实活]\label{thm:noflat}
去掉边界条件，标度乘积只是\emph{某个}常数，而任何常数都与真空方程相容。
故渐近平直性不是约定。
\leanline{Vacuum.product\_constant\_without\_flatness}
\end{theorem}

物理读法：\textbf{真空中引力只是用时间标度换径向标度，
在 $t$--$r$ 平面里不制造任何标度}——该平面的面元恒为一。
\leanline{Vacuum.tr\_area\_element\_constant}

\begin{remark}[结清到什么程度]
\emph{已导出}：给定"真空使标度和为常数"这一后果，倒数性仅由渐近平直性即可推出，
无需任何进一步选择。\emph{未导出}：真空场方程蕴含该和为常数——
那是标准的 Ricci 组合 $R_{tt}/A+R_{rr}/B=0$，此处作为输入；
形式化它需要对角拟设的曲率计算，那仍是引力扇区唯一未完成的一块。
\end{remark}

\section{光线偏折：各向同性扇区恰好只给一半}\label{sec:deflection}

范围声明里"定量引力效应仍未导出"这句话挂了好几版。本节导出一个，
而结果是一次\emph{失败}——一次尺寸可精确测量的失败，
这比一次成功更有用，因为它说清了方向性标度必须补上多少。

光线偏折是因为它穿过标度不同的区域。在\textbf{各向同性}扇区
（旧的标量公理，现已知是 A4 的退化轨迹）里，全部效应来自单个标量 $\sigma=-\Phi$，
掠射偏折为 $\alpha_{\text{iso}}=2GM/(bc^2)$；而实测值是它的两倍
$\alpha_{\text{obs}}=4GM/(bc^2)$。

\begin{theorem}[恰好一半]\label{thm:half}
$2\,\alpha_{\text{iso}} = \alpha_{\text{obs}}$，且缺失的那一半恰等于
各向同性扇区能给出的那一半。
\leanline{Deflection.iso\_is\_exactly\_half, Deflection.missing\_half,
Deflection.iso\_ratio\_half}
\end{theorem}

\begin{mdframed}[linecolor=red!60!black,linewidth=0.9pt,backgroundcolor=red!3]
\textbf{各向同性扇区不是近似地错，而是错了恰好一个因子 2。}
缺失的一半是\emph{空间}标度各向异性的贡献——
正是第 \ref{sec:frame} 节证明标量标度\emph{根本无法表示}的那部分。
于是这个因子 2 就是同一件事以数字的形式出现：
两半分别是时间标度与空间标度，而只有一个标度的理论永远只能提供其中一半。
\end{mdframed}

\subsection{掠日光线的数值}

取 $G=6.674\times10^{-11}$、$M_\odot=1.989\times10^{30}$ kg、
$b=R_\odot=6.957\times10^{8}$ m、$c=2.998\times10^{8}$ m/s：

\begin{center}
\begin{tabular}{lll}
\toprule
 & 预言 & 实测（VLBI）\\
\midrule
观测支 $4GM/bc^2$          & $1.7515''$ & $1.7509\pm0.0002''$ \\
各向同性扇区 $2GM/bc^2$    & $0.8758''$ & —— \\
\bottomrule
\end{tabular}
\end{center}

\leanline{Deflection.solarObs\_value, Deflection.solarIso\_value,
Deflection.solar\_ratio, Deflection.iso\_fails\_by\_far}

观测支与实测符合到约 $0.04\%$；而各向同性扇区差了 $0.87''$，
是 VLBI 误差的四千倍。这不是建模精度问题，是结构性问题。

\textbf{本节的地位：} 这是本报告最尖锐的定量陈述，
而它是对标量公理的\emph{反驳}，不是对任何东西的确认。
方向性扇区要给出缺失的一半，仍需完成各向异性曲率的具体计算——那尚未完成。
本节确立的是\textbf{剩余欠债的精确大小}。

\section{把另一半付清}\label{sec:ppn}

第 \ref{sec:deflection} 节把欠债定得很精确：各向同性扇区给出 $2GM/bc^2$，
实测是 $4GM/bc^2$，缺失的一半必须来自\emph{空间}标度。本节付清它，
而付款来自一条已经证过的关系。

掠射偏折对时间与径向两个单位分别依赖，标准参数化把这依赖记进一个数：
以 $g_{tt}=-(1-2\mu)$、$g_{rr}=1+2\gamma\mu$ 计，
\[
\alpha = \frac{2GM}{bc^2}\,(1+\gamma).
\]
$\gamma=0$ 是各向同性扇区（只有时间标度），$\gamma=1$ 是实测值。
于是全部欠债化为一个问题：\textbf{什么定住 $\gamma$？}

第 \ref{sec:frame} 节已经回答过。Schwarzschild 关系在标度语言中是
$s_t\cdot s_r=1$——时间与径向单位互为倒数——而
定理 \ref{thm:frameflat} 证明标量标度不平直化就无法满足它。
互为倒数即 $g_{tt}g_{rr}=-1$，即径向因子是时间因子的逆。把那个逆展开，
线性系数就被逼定了：

\begin{theorem}[定住系数的展开]\label{thm:recexp}
\[
(1-2\mu)^{-1} = 1 + 2\mu + \frac{4\mu^2}{1-2\mu},
\]
是恒等式而非近似。$\mu$ 的系数是 $2$，即 $\gamma=1$；余项被显式写出而非丢弃。
\leanline{PPN.reciprocal\_expansion, PPN.gamma\_eq\_one\_of\_reciprocal,
PPN.gamma\_eq\_one\_first\_order}
\end{theorem}

\begin{mdframed}[linecolor=axcol,linewidth=0.9pt,backgroundcolor=axcol!4]
\textbf{倒数条件强制 $\gamma=1$，而 $\gamma=1$ 给出实测偏折。}
各向同性扇区拿不出的那一半就是径向标度，而它的大小从来就不自由。
\leanline{PPN.deflection\_reciprocal\_eq\_observed, PPN.halves\_equal}
\end{mdframed}

Cassini 把 $\gamma$ 限制在 $1$ 的约 $2\times10^{-5}$ 之内。
本框架的倒数条件预言 $\gamma=1$ 且\textbf{无自由参数}，
故这是一次真实且目前成功的检验——检验的是\emph{方向性}扇区，
各向同性扇区此前已被排除。

\begin{remark}[结清到什么程度]
缺失的一半不再是无解释的：它是径向标度，其大小由倒数关系强制而非拟合。
第 \ref{sec:deflection} 节的欠债在\textbf{系数层面}已结清。
仍未做的是逐分量曲率计算，那才能说明场方程\emph{产生}一个互为倒数的位形；
在这里，倒数性是被描述的解的性质，还不是从源导出的推论。
\end{remark}

\section{第二个数字：水星}\label{sec:precession}

第 \ref{sec:deflection}、\ref{sec:ppn} 节给出一个定量结果，
并证明各向同性扇区差了 2 倍。一个数字有巧合风险。
本节给出第二个、性质不同的数字，而各向同性扇区在其上差了 \textbf{3 倍}——
差的倍数不同，这才是有用的地方：两处偏差无法被同一个重标度同时吸收。

近日点进动（每轨道）的标准参数化为
\[
\Delta\phi = \frac{2-\beta+2\gamma}{3}\cdot\frac{6\pi GM}{a(1-e^2)c^2}.
\]
$\gamma$ 就是第 \ref{sec:ppn} 节被倒数条件 $s_ts_r=1$ 强制为 $1$ 的那个空间标度权重；
$\beta$ 度量时间标度的非线性，在 A6$'$ 的标度响应下为 $1$。

\begin{theorem}[系数]\label{thm:precoeff}
$\gamma=\beta=1$ 给出系数 $1$；各向同性扇区（$\gamma=0$）给出 $1/3$。
系数对 $\gamma$ 是仿射的，斜率 $2/3$。
\leanline{Precession.coeff\_reciprocal, coeff\_isotropic, isotropic\_is\_third,
coeff\_affine\_in\_gamma}
\end{theorem}

\begin{center}
\begin{tabular}{lll}
\toprule
 & 预言 & 实测 \\
\midrule
倒数扇区（$\gamma=1$） & $42.99''/$世纪 & $42.98\pm0.04''/$世纪 \\
各向同性扇区（$\gamma=0$） & $14.33''/$世纪 & —— \\
\bottomrule
\end{tabular}
\end{center}

\leanline{Precession.mercury\_value, Precession.mercury\_isotropic\_value}

\begin{mdframed}[linecolor=red!60!black,linewidth=0.9pt,backgroundcolor=red!3]
\textbf{真正重要的一致性检验。} 各向同性扇区在光线偏折上差 $2$ 倍，
在近日点进动上差 $3$ 倍。\emph{倍数不同}，而两者都由在\emph{同一个}参数化中
令 $\gamma=0$ 得到。\textbf{一个缺失成分——空间标度——同时解释两处偏差，
且没有任何可调余地。}
\end{mdframed}

\section{耗散之下的局域广义协变性}\label{sec:cov}

对 A6 最直接的反驳是：若基准标度永不停止漂移，局域物理岂不被摧毁？
氢原子为什么昨天和今天一样大？太阳系为什么不膨胀？

答案由 A5 强制。宇宙耗散是基准的\emph{空间常数}漂移
$\sigma_t=\sigma+c(t)$，$\dd_i c(t)=0$；而一切几何量都由 $\sigma$ 的\emph{导数}构成，
漂移逐项相消。

\begin{theorem}[局域几何与宇宙时无关]\label{thm:cov}
对任意两个宇宙时刻 $t,t'$：
\[
\sigma_i,\quad \Gamma^a{}_{bc},\quad R^a{}_{bce},\quad R_{be},\quad \Defm_{ij},
\quad e^{2\sigma}R
\]
逐一相等。
\leanline{Covariance.CosmicHistory.no\_local\_expansion}
\end{theorem}

\begin{corollary}[场方程形式不变]\label{thm:covfield}
若 $e^{2\sigma}R=\kappa\rho$ 在某一时刻成立，则它在\emph{每一}时刻以相同的 $\kappa$
与相同的局域 $\rho$ 成立。全局耗散不重整化局域引力，局域广义协变性严格保持。
\leanline{Covariance.CosmicHistory.field\_equation\_form\_invariant}
\end{corollary}

物理结论：\textbf{束缚系统不膨胀，这是定理而非假设}。膨胀不是空间对物质施加的局域力；
它只在比较\emph{两个不同基准}时才可见——而这恰恰就是一次宇宙学红移测量所做的事。

\section{奇点只是极小的尺度}\label{sec:sing}

A2 把标度写成 $s\in A^{\times}$，即\emph{单位群}的元素。这不是技术选择：
一个度量单位必须是可以拿来做除法的东西。后果是结构性的。

\begin{theorem}[无度规退化]\label{thm:sing}
$s\ne 0$ 处处成立；$\varepsilon=1/s\ne 0$ 处处成立；共形因子 $s^2$ 处处可逆。
\leanline{Singularity.ScaleField.scale\_ne\_zero, energy\_ne\_zero,
conformal\_factor\_isUnit}
\end{theorem}

\begin{theorem}[曲率只看标度的变化，不看标度的值]\label{thm:singcurv}
若两个标度场的一阶、二阶导数相同，则其 Riemann 张量相同，无论二者的标度相差多少。
\leanline{Singularity.curvature\_indep\_of\_scale\_value}
\end{theorem}

由主定理 \ref{thm:master}，Riemann 张量是 $\sigma$ 导数的\emph{多项式}，$s$ 根本不出现。
因此\textbf{标度变小本身不产生任何曲率发散}。广义相对论所报告的奇点，
在此框架下是 $\sigma\to$ 极负的区域：尺度极小、能标极高，但二者都有限且都可逆。
理论在那里并不失语。

这与"奇点是数学奇点"的图景相反，而与本纲领的物理直觉一致：
奇点是一个尺寸极小的点，不是一个洞。

\section{黑洞、引力波，以及哪些标准问题是我们的}\label{sec:horizon}

非交换曲率闭合之后，引力扇区可以往前推。
标准问题中有些在这里问得出来，有些问不出来，\textbf{把它们分开是一半的工作}。

\subsection{没有奇点——也没有视界}

第 \ref{sec:sing} 节证明标度作为单位不可消失，故无度规退化。
同一事实还说了一句此前没被引出的更强的话。

\textbf{视界是度规时间分量为零之处。}对角方向性标度下该分量是 $-s_t^2$，而 $s_t$ 是\textbf{单位}：

\begin{theorem}[无视界、无无穷红移]\label{thm:nohor}
$g_{tt}=-s_t^2$ 永不为零，故\textbf{没有视界}；
其倒数伙伴 $s_r=s_t^{-1}$ 也永不发散（第 \ref{sec:vacuum} 节把二者绑定）；
而红移是\emph{两个单位之比}，故是单位，故\textbf{既不为零也不为无穷}。
\leanline{Horizon.gtt\_ne\_zero, reciprocal\_never\_degenerate,
redshift\_never\_infinite, no\_degeneracy\_anywhere}
\end{theorem}

\begin{mdframed}[linecolor=axcol,linewidth=1.2pt,backgroundcolor=axcol!6]
\textbf{无奇点、无视界、无无穷红移——三者出自同一个事实：
标度是比值，而比值可逆。}
\end{mdframed}

\textbf{那么观测现象如何解释？}一切经典检验都由\emph{外部}算出，而外部原封未动：
从对角拟设经第 \ref{sec:riccidiag} 节到第 \ref{sec:vacuum} 节的链条不变，
故\textbf{黑洞阴影、最内稳定轨道、铃宕频率、光线偏折全部与广义相对论一致}。
黑洞在这里是\textbf{标度比值极端但\emph{有限}}的区域：\textbf{任意黑，但不全黑}。

三条与标准图景不同的推论，同源于上述定理：
\textbf{没有任何东西被因果隔断、没有信息被销毁、红移有界而非发散}。
\textbf{这三条目前都不可测，我不作相反声称。}

\subsection{引力波，以及一个尖锐的判别式}

第 \ref{sec:light} 节证明"是否类光"\textbf{不依赖于标度}。由此分出两件事。

\begin{theorem}[引力波恰好走在光锥上]\label{thm:gw}
引力波是标度的传播扰动——因为几何\emph{就是}标度花样。
而标度扰动移动不了光锥。
故\textbf{引力波与光共用光锥：同速、无色散、无质量}。
\leanline{Horizon.null\_cone\_scale\_independent, waves\_travel\_on\_the\_light\_cone}
\end{theorem}

\begin{mdframed}[linecolor=axcol,linewidth=1.2pt,backgroundcolor=axcol!6]
\begin{theorem}[禁止能量依赖的光速]\label{thm:nodisp}
由 A3，能量\emph{就是}标度之倒数，故标度无关的光锥就是\textbf{能量无关}的光锥。
\textbf{本框架禁止能量依赖的光子到达时延，任何阶。}
\leanline{Horizon.no\_energy\_dependent\_speed}
\end{theorem}

这是一个真判别式：\textbf{带最小长度的方案普遍预言一个小的能量依赖延迟，
而本框架预言恰好为零。}伽马暴计时已在约束领头效应，每一次改进都在直接检验它。
\end{mdframed}


\subsection{偏振数恰为二，而呼吸模是\emph{被禁止}的}\label{ss:pol}

波是标度花样的扰动。数它的独立模式是算术，前提是先把扇区分开。

\begin{theorem}[两个偏振，来自两个不同扇区]\label{thm:pol}
三个空间方向中，一个是传播方向，横向维数为 $2$。于是
\begin{itemize}[leftmargin=1.6em,itemsep=1pt]
\item \textbf{标度扇区}贡献 $2-1=1$ 个模（即 $+$ 模，$\delta s_x=-\delta s_y$）；
\item \textbf{转动扇区}贡献 $\dim\Lambda^2\mathbb R^2=1$ 个模（即 $\times$ 模）。
\end{itemize}
合计 $2$。\textbf{两者来源不同}——一个是标度之差，一个是转动——
这是广义相对论\emph{不}作的区分，而在此由第 \ref{sec:coupling} 节的 Cartan 拆分逼出。
\leanline{Horizon.two\_polarizations, polarization\_count}
\end{theorem}

\begin{mdframed}[linecolor=axcol,linewidth=1.2pt,backgroundcolor=axcol!6]
\textbf{那个被减掉的 $1$ 不是约定，是 A5。}
标量（呼吸）偏振就是横向的\emph{均匀}重标度——而那恰是一次基准平移，
A5 宣布它不可观测。故呼吸模\textbf{不是"在解中恰好为零"，而是不可表达}。

\textbf{这是对标量--张量理论的判别式}：那类理论一般\emph{预言}呼吸模，
而本框架预言其振幅\textbf{恒等于零}，不是小。
足够大的探测器网络的偏振检验直接测这一条。
\leanline{Horizon.no\_breathing\_mode}
\end{mdframed}

\subsection{黑洞熵是对数的，不是面积律}

熵在这里是可分辨结构的计数（第 \ref{sec:entropy} 节），
而第 \ref{sec:spectrum} 节测出其密度 $\rho$。故跨越标度比值 $R$ 的区域，熵为
\[
S = \rho \ln R .
\]

\begin{theorem}[比值加倍只增加一个常数]\label{thm:ent}
面积律会使熵\emph{翻四倍}；此处增量只依赖于比值增长的\emph{因子}，
\textbf{与已达到的比值无关}。
\leanline{Horizon.entropy\_of\_scale\_ratio, entropy\_doubling\_is\_additive,
entropy\_increment\_independent\_of\_size}
\end{theorem}

\begin{mdframed}[linecolor=red!60!black,linewidth=1pt,backgroundcolor=red!3]
\textbf{这与 Bekenstein--Hawking 面积律矛盾。}
我把它作为\emph{分歧}而非任何东西的推导来陈述：
面积律被广泛相信但\textbf{从未被测量}，
而本框架无法复现它，因为面积律的推导需要视界，而上一小节说视界不存在。
\textbf{要么框架在此处错，要么面积律错，而现有的一切都判不了。}
\end{mdframed}

\subsection{哪些问题是借来的}

\textbf{引力子。}引力子是一个具有粒子激发的场的量子。
第 \ref{sec:emergence} 节说不存在标度无关的物种清单，
而此处引力不是几何上的场，\textbf{它就是标度本身的不均匀}。
故"引力子是什么"同时引进了框架所拒绝的场本体论与物种本体论。
\textbf{取代它的是定理 \ref{thm:ncdefm}}：引力的量子内容是标度梯度的对易子，
而那不是一个粒子。

\textbf{散射振幅。}振幅需要内容固定的渐近入态与出态。
第 \ref{sec:emergence} 节说没有标度携带完整清单，故"那些"渐近态并不存在。
存活下来的是\textbf{两个被指名标度之间的映射}——
而这正是重整化振幅实际所是；\textbf{它的标度依赖不再是麻烦，而成为内容}。

\textbf{短距完备化。}完备化预设一个最后的标度。

\begin{theorem}[没有最小长度]\label{thm:nomin}
对任何被提议的最短标度都存在更细者：阈值无界，且标度是单位故分辨率无上限。
\textbf{本框架预言没有最小长度}，这与一切预设基本长度的纲领矛盾——
而这个预言正是定理 \ref{thm:nodisp} 使之可检验的那一个。
\leanline{Horizon.no\_minimum\_length, no\_final\_description}
\end{theorem}

\begin{remark}[四项的结局]
两项被回答（引力波、熵），
一项被\emph{否定地}回答且带可检验推论（短距完备化），
一项作为借来的问题被\textbf{谢绝}（引力子），并指明取代它的是什么。
\end{remark}

\clearpage


\section{引力波动力学：旋近、波形与辐射功率}\label{sec:waves}

第 \ref{sec:horizon} 节做的是运动学——波恰在光锥上、无色散、偏振恰为二、呼吸模被禁止。
本节做\textbf{动力学}，以便与脉冲双星计时和干涉仪数据对照。

\subsection{诚实的分工，先说清楚}

\begin{mdframed}[linecolor=axcol,linewidth=1.2pt,backgroundcolor=axcol!6]
\textbf{框架自己导出的：多极结构。}
定理 \ref{thm:conserv} 证明场方程\emph{只有在}其源守恒时才写得下来——
守恒在这里不是额外公设，\textbf{是方程存在的条件}。
那条定理干的正是通常靠假设 $\nabla_\mu T^{\mu\nu}=0$ 才能干的活。

\textbf{一个输入：系数 $\kappa$}（A6$'$）。但它\emph{只出现一次}——
静态 Poisson 方程与四极系数共用同一个 $\kappa$，
故由 Newton 极限定住它之后，\textbf{辐射率没有任何剩余自由度}。
这就是为什么 Hulse--Taylor 是一次\emph{检验}而不是一次拟合。

\textbf{属于标准数学、不是物理引进的：}给定 $\Box$ 与守恒源，
把推迟 Green 函数按多极展开是微积分。$32/5$ 与间距的五次方是那次展开的后果。
\end{mdframed}

\subsection{什么在辐射：守恒定理杀掉低阶}

\begin{theorem}[单极与偶极都不辐射]\label{thm:nolowmult}
守恒矩的时间导数为零，故它恒定；而恒定的矩在任何阶都不贡献辐射场。
单极是总源，偶极的一阶导是总动量——\textbf{同一条守恒律把两者都冻住}。
\leanline{Waves.monopole\_does\_not\_radiate, Waves.dipole\_does\_not\_radiate,
Waves.quadrupole\_is\_leading}
\end{theorem}

\begin{mdframed}[linecolor=axcol,linewidth=1.2pt,backgroundcolor=axcol!6]
\textbf{"无偶极辐射"在这里是预言，不是约定。}
标量--张量理论一般\emph{会}在偶极阶辐射，因为它们带一个\textbf{不守恒}的标量荷。
本框架没有这样的荷：定理 \ref{thm:pol} 去掉了标量偏振，
而定理 \ref{thm:conserv} 去掉了偶极。\textbf{脉冲双星计时测的正是这一条。}
\end{mdframed}

\subsection{圆轨道双星}

\begin{theorem}[辐射功率与旋近]\label{thm:inspiral}
$P = \tfrac{32}{5}\,G^4(m_1m_2)^2(m_1+m_2)/(c^5r^5) > 0$；
间距减半，功率乘以 \textbf{32}——这就是旋近失控的原因。
而 $E=-Gm_1m_2/2r$ 使能量损失强制 $\dot r<0$，闭形式为
\[
\dot r = -\frac{64}{5}\,\frac{G^3m_1m_2(m_1+m_2)}{c^5r^3}.
\]
\leanline{Waves.power\_pos, Waves.power\_ratio\_halving,
Waves.inspiral\_shrinks\_orbit, Waves.separation\_rate}
\end{theorem}

\subsection{啁啾：干涉仪实际测的量}

\begin{theorem}[啁啾质量]\label{thm:chirp}
频率演化只通过\textbf{一个}质量组合进入：
$\mathcal M=(m_1m_2)^{3/5}/(m_1+m_2)^{1/5}$，且 $\d f/\d t \propto \mathcal M^{5/3}f^{11/3}$。
$\mathcal M$ 对两质量对称——\textbf{领头阶波形分不出哪个是哪个}；
等质量时 $\mathcal M\cdot 2^{1/5}=m$。
\leanline{Waves.chirpMass\_symm, Waves.chirpMass\_equal,
Waves.chirp\_exponent, Waves.chirp\_mass\_exponent}
\end{theorem}

\subsection{对数据}

\begin{center}
\begin{tabular}{lll}
\toprule
观测 & 本框架 & 结果 \\
\midrule
PSR B1913$+$16 衰减／四极预言 & $1$（无偶极项） & $0.9983\pm0.0016$ \\
\textbf{PSR J0737$-$3039（双脉冲星）} & $1$（无偶极项） & \textbf{符合到 $0.013\%$} \\
偶极辐射分数 & \textbf{恒等于 $0$} & $<1.3\times10^{-4}$ \\
GW250114 啁啾质量（源系） & —— & $28.6\pm0.5\,M_\odot$ \\
\bottomrule
\end{tabular}
\end{center}

双脉冲星十六年计时把四极公式检验到 $0.013\%$，比 Hulse--Taylor 紧一个量级
（$\dot P_b$ 的相对精度 $6\times10^{-5}$）。
\leanline{Waves.doublePulsarPrecision, Waves.dipole\_bounded\_by\_double\_pulsar,
Waves.double\_pulsar\_tighter, Waves.hulse\_taylor\_agreement}

\begin{mdframed}[linecolor=red!60!black,linewidth=1.2pt,backgroundcolor=red!3]
\textbf{必须说清这些符合\emph{检验了什么}。}

\textbf{它们不区分本框架与广义相对论。}各向同性扇区里 SCD 给出 $\gamma=1$、
波在光锥上、四极领头、静态与辐射扇区共用同一个 $\kappa$——
\textbf{领头阶波动预言与广义相对论逐项相同}。
故双脉冲星同时确认两者，\textbf{对二者之间毫无分辨力}。

它\emph{确实}做到的是把 $\{$SCD, GR$\}$ 与带标量通道的理论分开。
而在这一侧分得更利落：标量--张量理论有一个可以调小的耦合，
而这里\textbf{偶极被一条定理禁止}、\textbf{呼吸模不可表达}。
\textbf{一旦测到偶极项，这里没有东西可调——框架直接失败。}

\end{mdframed}

\subsection{框架\emph{确实}分道之处：没有视界}

\begin{mdframed}[linecolor=red!60!black,linewidth=1.2pt,backgroundcolor=red!4]
定理 \ref{thm:nohor} 证明度规永不退化，故\textbf{没有视界}——
只有标度比值极端但有限的区域。
而广义相对论算铃宕时用的边界条件是\textbf{在视界上纯内向}。
没有视界就没有那个面，边界条件必须不同，
而部分反射的边界一般会在铃宕后产生\textbf{晚期回波}。

\textbf{这是波动扇区里两套理论唯一可能分开的地方，而数据目前在往反方向推。}
对 GW150914、GW231226、GW250114 的模型无关搜索\emph{未}发现显著回波信号并给出上限；
GW250114 的铃宕与 Kerr 黑洞一致。

\textbf{框架能说与不能说的：}它说反射率\emph{不}恰好为零，因为标度是单位。
它\textbf{不}预言振幅或延迟——那取决于标度逼近零的程度，
而那由具体解、最终由那一个自由单位决定。
\textbf{所以零结果尚未证伪任何东西，而框架在此也还没给出可证伪的东西。}
\textbf{登记为开放项，不是预言。}
\leanline{Waves.reflectivity\_nonzero\_but\_unbounded}
\end{mdframed}

\subsubsection*{振幅算不出来，但回波序列的\emph{结构}算得出来}

反射率随天体而异，\textbf{本框架算不了}——它由内部解决定，
而内部解需要 A6$'$ \emph{再加}一个物态方程，两者框架都不提供。
但搜索实际上模板匹配的是\emph{结构}，而结构只用度规形式就够，不需要内部模型：

\begin{theorem}[回波延迟对红移是\emph{对数}的]\label{thm:echodelay}
$\Delta t \approx \tau\ln(1/A_{\min})$，故 $A_{\min}$ 平方只把延迟\textbf{加一倍}，
十次方只乘十倍。\textbf{把 $1/A_{\min}$ 提高十个数量级，延迟只增大约十倍，不是 $10^{10}$ 倍。}
这就是为什么回波要在\textbf{毫秒}而不是 Planck 时间上找。
\leanline{Waves.echo\_delay\_logarithmic, Waves.delay\_insensitive\_to\_reflectivity}
\end{theorem}

\begin{theorem}[序列是均匀梳，振幅几何衰减]\label{thm:echocomb}
第 $n$ 次回波到达于 $n\Delta t$（每次往返代价相同），
振幅按 $R^{2n}$ 衰减。故\textbf{零结果直接给出 $R$ 的界}：
对第一次回波的振幅上限就是对 $R^2$ 的上限。
\leanline{Waves.echo\_train\_is\_a\_comb, Waves.echo\_amplitudes\_geometric,
Waves.null\_search\_bounds\_reflectivity}
\end{theorem}

\begin{mdframed}[linecolor=axcol,linewidth=1pt,backgroundcolor=axcol!5]
\textbf{框架能给的到此为止：}\emph{在哪里找}（毫秒量级，因为延迟是对数的）、
\emph{找什么形状}（均匀梳、几何衰减），
以及\emph{零结果意味着什么}（直接是 $R^2$ 的上限）。
\textbf{给不了的是 $R$ 本身。}
\end{mdframed}


\subsection{GW250114：面积定理确认了什么，又\emph{没}确认什么}\label{ss:gw250114}

迄今最响的事件给出了最锐利的铃宕数字，以及对 Hawking 面积定理的确认。
两者都直指"没有视界"这个结论，故在此正面处理。

\textbf{实测数字（取自原始论文源码，非摘要）：}网络匹配滤波信噪比 $80$；
源系分量质量 $33.6^{+1.2}_{-0.8}$ 与 $32.2^{+0.8}_{-1.3}\,M_\odot$，
遗迹 $62.7^{+1.0}_{-1.1}\,M_\odot$、自旋 $0.68\pm0.01$；
分量自旋都小（$\chi_1\le0.24$、$\chi_2\le0.26$，90\%）；
源系啁啾质量 $28.6\pm0.5\,M_\odot$。

\textbf{检验强度：}铃宕模被约束在 Kerr 谱的 $\pm30\%$ 之内；
\textbf{面积定理在所有分析窗上均达 $\ge3.4\sigma$，窗起点取 $-10M_{\rm tot}$ 时超过 $5\sigma$}，
用泛音的那一版为 $3.6\sigma$；
而扣除最佳广义相对论波形后，\textbf{残差表现得像普通探测器噪声}。
\leanline{Waves.gw250114ChirpMass, Waves.gw250114\_radiated\_fraction,
Waves.area\_law\_confirmed}

\begin{mdframed}[linecolor=red!60!black,linewidth=0.9pt,backgroundcolor=red!3]
\textbf{一处更正：}本报告先前把 GW250114 的啁啾质量引作 $30.7\,M_\odot$，
那是\emph{探测器系}的数；\textbf{源系是 $28.6\pm0.5$}。
先前的数取自摘要转述，现取自原始论文源码。
\end{mdframed}

\begin{mdframed}[linecolor=red!60!black,linewidth=1.2pt,backgroundcolor=red!4]
\textbf{面积定理的确认\emph{没有}探测到视界。}

该检验是由实测质量与自旋经 Kerr 公式\emph{推算}出各自面积，再核对
$A_f \ge A_1+A_2$。\textbf{它是一条关于\emph{被推算参数}的自洽关系。}
本框架精确复现外部（$\gamma=1$、同一个 $\kappa$、
且对称扇区在每一阶上都是 GR），故推出同样的质量与自旋，
\textbf{通过同一个核对}。

\textbf{被观测的不是视界，是一个参数化。}
这一点明说，是因为反过来读会得到一个既容易又错误的安心。
\end{mdframed}

\subsubsection*{本框架有自己的不减定理，而且是导出的}

第 \ref{sec:entropy} 节的 H 定理\textbf{只用有限分辨率}就成立——无任何概率假设；
而第 \ref{sec:horizon} 节使一个区域的熵为 $\rho\ln R$（$R$ 为标度比值）。两者合起来：

\begin{theorem}[标度比值永不减小]\label{thm:ratiomono}
熵不减且 $\rho>0$ $\Rightarrow$ $\ln R$ 不减 $\Rightarrow$ $R$ 不减。
\leanline{Waves.scale\_ratio\_never\_decreases}
\end{theorem}

\textbf{这是本框架对应面积定理的那一条，从完全不同的方向到达}——
不是从视界的因果结构，是从粗粒化。

\begin{mdframed}[linecolor=axcol,linewidth=1pt,backgroundcolor=axcol!5]
\textbf{但它\emph{不是}面积定理}：$\ln R$ 与 $A\propto M^2$ 是不同的函数
（比值加倍只增加一个常数，面积律则翻四倍）。
\textbf{目前没有任何可测量能把二者分开}，因为熵本身没有被测量——
被测的只是用来算它的那些参数。
\leanline{Waves.log\_law\_is\_not\_area\_law}
\end{mdframed}

\subsubsection*{GW250114 真正咬住我们的地方}

\textbf{残差与噪声一致，是目前对回波振幅、因而对反射率最紧的界。}
本框架预言反射率非零而量级未定，
故这个界是对它的一条\textbf{活约束，不是对任何东西的确认}。

而衰减时间（$9\%$）比频率（$2\%$）松，
\textbf{而衰减时间恰是反射边界会移动的那个量}——
如果边界条件效应要藏，它藏在那里。
\leanline{Waves.damping\_is\_the\_looser\_constraint,
Waves.null\_search\_bounds\_reflectivity}


\subsection{第二代黑洞：一条比回波更强的约束}\label{ss:secondgen}

GW241011 与 GW241110 被报告为具有\textbf{快速且被精确测定的主自旋}、
显著的自旋--轨道错位、以及不等的质量比——
这是\textbf{层级并合}的特征：较重的那个本身就是更早一次并合的遗迹。
GW241011 是迄今第三响的事件（网络信噪比 $36.0$），
主自旋 $\chi_1=0.64^{+0.06}_{-0.09}$，其为正号\emph{基本以全置信度}确立。

\begin{mdframed}[linecolor=red!60!black,linewidth=1.2pt,backgroundcolor=red!4]
\textbf{这为什么打到"没有视界"上，而且不是好消息。}

该文用快速自旋排除了质量约 $10^{-13}$ 到 $3\times10^{-12}$ eV 的超轻玻色子：
合适 Compton 波长的玻色子会通过\textbf{超辐射}把黑洞自旋降下来。

\textbf{而同一个超辐射放大，正是使带能层的\emph{无视界}致密天体不稳定的机制}——
视界本会吸收掉的模式，被反射并无界放大，
而\textbf{反射率越大，不稳定增长越快}。

所以一个被观测到的\textbf{高速自旋、且显然长寿}的致密天体，
把反射率从上方压住，\textbf{而且是随时间指数地压}，不是像回波那样线性地压。

\textbf{这是本框架此前未登记的一条真实压力。}
它来自读原始论文源码而非摘要转述，理应如此登记：
\textbf{回波搜索对 $R$ 的约束是弱的，而高速自旋遗迹的约束严厉得多。}
框架仍算不出 $R$，故\textbf{尚不能说它是否幸存}。
\leanline{Waves.gw241011\_spin\_is\_rapid, Waves.spin\_bound\_beats\_echo\_bound}
\end{mdframed}

\begin{mdframed}[linecolor=axcol,linewidth=1pt,backgroundcolor=axcol!5]
\textbf{有一件事是往对的方向走的。}玻色子排除本身与本框架相容——
它\textbf{根本不预言任何超轻标量}：标签结构是
$(\text{阈值},\mathbb Z/2,\mathbb Z)$ 且\textbf{没有第四个槽位}，
而呼吸模不可表达，正好去掉了这样一个场会携带的标量偏振。
\textbf{标量--张量理论需要解释这次未探测；这里没有东西可供探测。}
\leanline{Waves.no\_ultralight\_scalar\_available}
\end{mdframed}

\subsection{对称扇区在\emph{每一阶}上都是广义相对论}

很容易期待 SCD 与 GR 在领头阶一致、在下一阶分开。\textbf{它们不会，}
而原因是一条定理，不是展开的巧合。

\begin{mdframed}[linecolor=axcol,linewidth=1.2pt,backgroundcolor=axcol!6]
定理 \ref{thm:ncsym} 证明形变张量的\textbf{对称}部分\emph{恰是}经典那个，
\textbf{任何阶都没有对易子修正}。而每一个后 Newton 系数都由那个对称部分算出。于是
\begin{center}
\textbf{本框架的后 Newton 展开就是广义相对论的展开，到所有阶，}\\
\textbf{而不只是到已被测量的那一阶。}
\end{center}
故\textbf{计算更高阶不可能区分二者}，而理由不是差别太小——\textbf{那里没有差别}。
\leanline{NCConformal.symmetric\_part\_uncorrected}
\end{mdframed}

全部差别住在反对称部分，即 $[\sigma_i,\sigma_j]$，
而第 \ref{sec:axes} 节把它认作转动标签——\textbf{它耦合到自旋}，
且尺度是一个标度步长，故在任何引力系统中都可忽略。

\begin{mdframed}[linecolor=red!60!black,linewidth=1pt,backgroundcolor=red!3]
\textbf{值得明说而不是暗示的结论：}
\textbf{引力波扇区在任何可达精度上都不会把本框架与广义相对论分开，
花在那里的力气就是在复现广义相对论。}
唯一的例外是上一小节的铃宕边界条件，而那根本不是一个后 Newton 问题。
\end{mdframed}

\begin{remark}[仍然没有的]
数值相对论、并合、自旋进动的定量处理。
\textbf{但按上文，这些做出来只会复现广义相对论}——
它们是工程量而非物理增量。本节处理的是领头阶、圆轨道、无自旋的旋近。
\end{remark}

\clearpage
\part{第五部分\quad 各向异性扇区：标签与荷}

\begin{mdframed}[linecolor=axcol,linewidth=1.2pt,backgroundcolor=axcol!6]
\textbf{本部分的一切只存在于各向异性轨迹上。}标度花样的方向一旦处处一致，
以下全部对象——序参量、绕数、块荷、多重态、内禀对称性——\textbf{一个都不存在}，
不是取平凡值，是\emph{不存在}（第 \ref{sec:locus} 节）。

而花样是否各向异性是\textbf{标度相关}的事实。
所以本部分的标签不是世界的性质，是\emph{某个标度轨迹}的性质——
它们在各向同性破缺处一并诞生，因为它们同源。
\textbf{这就是“相互作用随能标下降而分离”，是导出的，不是假设的。}
\end{mdframed}

\section{内禀对称性无处可去}\label{sec:gauge}

规范理论挑一个李群、在每点挂一个内禀空间、让联络作用其上。
正是这份自由，使标准模型的群成为输入：形式体系本身不指定任何群。

本框架\emph{没有}这份自由，而这恰恰是有意思的地方。
\textbf{这里没有内禀空间。} 一点上唯一的结构是方向性标度 $s:\mathrm{Fin}\,n\to A^{\times}$。
因此，一个不是基底运动的对称性，只能是\emph{标度花样}的对称性——
一个使每个局域单位保持原样的方向重标号。

\begin{theorem}[内禀对称性即标度花样的稳定子]\label{thm:stab}
这些重标号构成一个群，其中没有任何选择余地。特别地：
\begin{enumerate}[label=(\roman*),leftmargin=2em]
\item 若各方向标度\emph{互不相同}，则内禀对称性\textbf{不存在}；
      规范自由不是通有的，它是简并的后果。
\item 若全部相同（各向同性扇区，即旧标量公理），则对称性是全置换群。
\end{enumerate}
\leanline{Gauge.stabilizer\_eq\_bot\_of\_injective,
Gauge.stabilizer\_eq\_top\_of\_isotropic}
\end{theorem}

\begin{theorem}[多重态就是简并块]\label{thm:multiplet}
两个方向被某个内禀对称性联系，当且仅当它们携带相同的局域单位。
故\textbf{规范结构与多重态结构同源}：都来自方向性标度之间的重合花样，
而多重态的大小就是简并块的大小。
\leanline{Gauge.sameOrbit\_iff\_eq\_scale, Gauge.internal\_symmetry\_determined}
\end{theorem}

这给出的是一条\emph{限制}，而非对 $SU(3)\times SU(2)\times U(1)$ 的推导：
框架断言内禀对称性\textbf{必须}是某个标度简并的稳定子，
故规范因子一块一个、多重态大小即块大小；
一个作用在与标度花样无关的内禀空间上的群，在这里\emph{不可用}。
这是关于"何种规范理论可以存在"的可证伪的结构性断言。

必须明说：\textbf{本节不断言观测到的花样是 $3+2+1$，也不解释代结构}。
那需要把标度花样本身导出来，而它没有被导出。

\section{点粒子回来了，而"比值"就是原因}\label{sec:axis}

第 \ref{sec:codim} 节断言三维中框架的缺陷是弦而非点。
那个结论用的是 \texttt{Defect.lean} 的\emph{标量}圆值标度——
而在公理提供方向性标度的地方使用标量标度，
正是第 \ref{sec:frame} 节早已诊断过的那个错误。\textbf{它被犯了第二次。}

A4 给的不是一个标度，是\emph{每个方向一个}标度。
所以一点上的位形不是一个相位，而是\emph{标度在各方向上的花样}，
场取值于这种花样的空间。那个空间不是圆，其缺陷也不只由回路分类。

决定性的性质是被逼出来的，而逼出它的恰恰是当初让标量图景显得不可避免的那件事：

\begin{mdframed}[linecolor=axcol,linewidth=0.9pt,backgroundcolor=axcol!4]
\textbf{标度是长度之比，所以它看不见方向的取向。}
把方向反过来，一切标度读数不变。
\end{mdframed}

\begin{theorem}[标度读数是偶的]\label{thm:even}
$Q(-v)=Q(v)$，对裸形式与被测形式都成立；
因而\textbf{没有任何标度测量能把一个方向与它的反向区分开}。
\leanline{Axis.bareForm\_neg, Axis.physForm\_neg, Axis.neg\_indistinguishable}
\end{theorem}

于是序参量是一条\textbf{无取向的轴}——一个射影对象，不是矢量。
而射影序参量众所周知同时携带两类缺陷：二阶的线缺陷，
以及\textbf{带整数荷的点缺陷}。

\begin{mdframed}[linecolor=axcol,linewidth=0.9pt,backgroundcolor=axcol!4]
\textbf{点粒子回来了，而它们回来是因为标度是一个比值。}
\end{mdframed}

\subsection{哪种花样支撑序参量}

\begin{theorem}[三种花样]\label{thm:patterns}
\begin{enumerate}[label=(\roman*),leftmargin=2em]
\item \emph{完全简并}（各向同性，即旧标量公理）：一切重标号都是对称性，
      没有任何东西逐点变化，\textbf{不支撑任何缺陷}；
\item \emph{完全非简并}：无内禀对称性，序参量是一整个标架；
\item \emph{单轴}（一个方向被挑出，其余简并）：对称性恰为块内重标号，
      序参量是\emph{一条无取向的轴}——正是携带点缺陷的那种位形。
\end{enumerate}
\leanline{Axis.isotropic\_no\_order\_parameter, Axis.nondegenerate\_full\_frame,
Axis.uniaxial\_stabilizer, Axis.uniaxial\_axis\_distinguished}
\end{theorem}

\begin{remark}[作为输入而非证明]
"三维中无取向轴场的 $\pi_2=\mathbb{Z}$、$\pi_1=\mathbb{Z}_2$，
故有整数荷点缺陷与二阶弦"是标准拓扑结论，此处作为输入。
框架提供的是\emph{射影性}，而点缺陷正是从射影性来的。
\end{remark}

\begin{mdframed}[linecolor=red!60!black,linewidth=0.9pt,backgroundcolor=red!3]
\textbf{更正记录。} 第 \ref{sec:codim} 节的定理对\emph{圆值}标度仍然成立，
被撤回的是"它们描述\emph{本}框架"这一主张：
在 A4 的方向性标度下它们不描述本框架，框架终究预言点粒子。

同一处替换（把方向性的东西当成标量的）已经犯了两次。
本节留档，以免犯第三次。
\end{mdframed}

\section{粒子：不用场的产生机制}\label{sec:defect}

量子场论用背景上的场的量子化产生粒子：粒子是激发，荷是表示标签，
质量是拉氏量的参数。这些在本框架中一样都没有。与其模仿，不如照标度图景直说。

\begin{mdframed}[linecolor=red!60!black,linewidth=0.9pt,backgroundcolor=red!3]
\textbf{本节用的是\emph{圆值}标度图景，它已被取代。}
第 \ref{sec:axis} 节表明：标度是一个\emph{比值}，故读数无向，
序参量是一根\textbf{无向轴}而非一个圆——这是方向性 A4 的后果。
本节关于绕数为整数、荷唯一、反粒子简并的定理\textbf{仍然成立}
（它们只用到"回路要么闭合要么差整数个周期"），
但由圆值图景推出的\emph{余维数}结论已被撤回（第 \ref{sec:codim} 节），
而缺陷的完整标签由第 \ref{sec:charges}、\ref{sec:particle} 节给出：
$(\text{阈值},\mathbb Z/2,\mathbb Z)$。
\end{mdframed}

若 A5 的连续标度对称性破缺到离散子群——标度只在模 $\Delta$ 的意义下有物理——
则对数标度可被当作\textbf{圆值}场处理。绕任一点的闭合回路一周，
标度必须回到\emph{模 $\Delta$ 的}自身，因而未必精确回到自身；
其偏差是 $\Delta$ 的整数倍。

\begin{mdframed}[linecolor=axcol,linewidth=0.9pt,backgroundcolor=axcol!4]
\textbf{那个整数就是粒子。} 粒子是\emph{度量单位中的缺陷}：
一处局域标度无法被梳理平整的地方。
\end{mdframed}

\begin{theorem}[荷是整数，且唯一确定]\label{thm:winding}
配置唯一决定其绕数；没有重新指派的自由。
荷之为整数，不是因为选定了某个表示，而是因为回路要么闭合、
要么差整数个周期。孤立的分数荷因此不可能出现。
\leanline{Defect.ScaleDefect.winding\_unique}
\end{theorem}

\begin{theorem}[粒子稳定]\label{thm:stable}
绕数非零的缺陷不是平凡配置：标度绕它确实不闭合。
稳定性是\emph{拓扑}的，证明中未使用任何关于力或势的假设。
\leanline{Defect.ScaleDefect.stable\_of\_winding\_ne\_zero,
winding\_eq\_zero\_of\_trivial}
\end{theorem}

\begin{theorem}[荷守恒；不能单独产生]\label{thm:pair}
绕数可加；粒子与其反粒子的总绕数为零，故可以从真空成对产生，
而单个缺陷不能。
\leanline{Defect.ScaleDefect.pair\_winding\_zero}
\end{theorem}

\begin{theorem}[反粒子质量严格相等]\label{thm:anti}
电荷共轭即取向反转，$|k|$ 不变，故
$m_{-k}=m_{k}$。这一简并不是作为对称性原理被强加的，
而是荷为绕数的推论。
\leanline{Defect.ScaleDefect.antiparticle\_same\_mass}
\end{theorem}

\begin{theorem}[几何质量塔——\emph{基于一条被假设的质量律}]\label{thm:spectrum}
\textbf{注意：} 此律是\emph{假设}而非导出，见第 \ref{sec:audit} 节的撤回。
若假定质量即逆标度且缺陷的荷每增加一就跨一个周期，则
\[
m_k=m_0e^{|k|\Delta},\qquad
\frac{m_{k+1}}{m_k}=e^{\Delta}\ \text{在每一层都相同}.
\]
\leanline{Defect.ScaleDefect.mass\_ratio\_geometric,
mass\_ratio\_const, mass\_zero}
\end{theorem}

\paragraph{与标准模型的差别，即可测试之处。}
这不重现标准模型，也不打算重现。它是一套不同的机制，具有不同的特征：
\emph{恒定的质量比}、\emph{向上无界的塔}、\emph{严格整数的荷}、
以及\emph{粒子—反粒子质量精确简并}。这四条是应当被检验的东西。

必须明说：\textbf{本节不给出规范群，不给出代结构，不产生任何具体粒子的数值}。
它给出的是一条不同于 QFT 的产生机制及其可证伪的形状。

\section{两个拓扑荷，只有一个被禁闭}\label{sec:charges}

第 \ref{sec:axis} 节让点粒子回来了：单轴标度花样的序参量是一条无取向的轴，
而这种序参量携带两类互不相同的缺陷。于是框架面临一个此前未曾面对的问题——
第 \ref{sec:color} 节的荷与点缺陷的荷是\emph{同一个}量子数，还是两个？

是两个，而它们的差别才是有意思的地方。

\begin{mdframed}[linecolor=red!60!black,linewidth=0.9pt,backgroundcolor=red!3]
\textbf{注意：}第 \ref{sec:color} 节的\emph{色}解释已被撤回（见第 \ref{sec:audit} 节）。
本节引用的只是它的\textbf{块单值性结构}，那部分成立；
而块群的取值由第 \ref{sec:particle} 节定死为 $\mathbb Z/2$。
\end{mdframed}

\begin{itemize}[leftmargin=1.6em]
\item \textbf{块荷}是绕一圈后简并方向所经历的置换。它住在一个\emph{有限}群里，
      而它正是第 \ref{sec:color} 节所约束的那个：位形可观测当且仅当标号闭合。
\item \textbf{刺猬荷}是轴在包围球面上的整数缠绕数。它住在 $\mathbb{Z}$ 里，
      而\textbf{没有任何东西约束它}。
\end{itemize}

\begin{theorem}[恰好一个被禁闭]\label{thm:onecharge}
\begin{enumerate}[label=(\roman*),leftmargin=2em]
\item 非平凡块荷不可单独观测；
\item \textbf{每一个}整数刺猬荷都被某个可观测态携带。
\end{enumerate}
\leanline{Charges.DefectCharge.block\_confined, hedgehog\_free,
exactly\_one\_confined}
\end{theorem}

\begin{theorem}[两个荷相互独立]\label{thm:indep}
固定块荷不定住刺猬荷，反之亦然。它们是两个真正独立的量子数，
不是同一个的两种读法。
\leanline{Charges.DefectCharge.charges\_independent,
observable\_pair\_any\_hedgehog}
\end{theorem}

\begin{mdframed}[linecolor=axcol,linewidth=0.9pt,backgroundcolor=axcol!4]
这份不对称不是放进去的。它来自两个荷住在\emph{同一个}序参量的\emph{不同}同伦群里：
有限的那个来自回路，整数的那个来自球面，
而可观测性条件——要求\emph{标号}闭合——只看得见前者。

由此得到的图案正是观测到的那个：\textbf{一个取有限多个值、被禁闭的荷，
外加一个不被禁闭的整数荷。}
\end{mdframed}

框架\emph{未}导出那个有限群是 $SU(3)$ 的中心，也未导出那个整数是电荷；
它导出的是：\textbf{必然恰有这两类荷，且恰有一个被禁闭}。

\section{代数允许的粒子模型}\label{sec:particle}

第 \ref{sec:algebra} 节定住了族。三项此前自由的输入随之被约束，本节把它推到头。

目标\emph{刻意}不是量子场论。第 \ref{sec:emergence} 节已论证物种清单是标度轴上某一截面的性质，
故复现它不是检验。这个层级的框架能被要求的是\textbf{标签结构}：
一个粒子是一份什么样的清单，以及两个相加会怎样。

\subsection{标度花样坍缩为一个矢量}

\begin{theorem}[缩放被方向矢量忠实标记]\label{thm:binj}
$\text{boost}\,v = \text{boost}\,w \Rightarrow v = w$。
连同定理 \ref{thm:rank1}，A4 的 $n$ 个独立单位被压成\textbf{一个矢量}，
在转动意义下压成\textbf{一个量值}。
\leanline{Particle.boost\_injective, pattern\_is\_one\_vector}
\end{theorem}

这是对一项长期输入的大幅削减——从 $n$ 个自由函数到一个数加一个方向——
但\textbf{不是消除}：那个量值仍然自由。已登记。

\subsection{恰好一个阈值速率，而这是可检验的}

第 \ref{sec:spectrum} 节由 A5 论证阈值是恒速率过程，得到 $\mathrm{CV}=1$。
秩一从结构上给出同一结论，而且尖锐得多，因为它规定了\emph{能有几个}速率。

秩一根系只有一个根，故只有\textbf{一个}指数速率。秩二以上会有多个根、多个速率，
从而是若干恒速率过程的交错——一个指数\emph{混合}。而混合必然过离散：

\begin{theorem}[恒等式]\label{thm:mix}
$p a^2 + (1-p) b^2 - (pa+(1-p)b)^2 = p(1-p)(a-b)^2$。纯代数。
\leanline{Particle.mixture\_second\_moment}
\end{theorem}

\begin{theorem}[任何两速率混合都有 $\mathrm{CV}^2 \geq 1$]\label{thm:mixge}
且\textbf{取等当且仅当两速率相同}——即恰好在秩一的情形。
\leanline{Particle.mixture\_cvSq\_ge\_one, mixture\_cvSq\_eq\_one\_iff}
\end{theorem}

\begin{mdframed}[linecolor=axcol,linewidth=1pt,backgroundcolor=axcol!5]
于是预言从"一个目标值"升级为\textbf{双侧陈述}：秩一给 $\mathrm{CV}^2=1$，
任何更高秩给 $\mathrm{CV}^2>1$。
观测值 $\mathrm{CV}^2=0.875$ \textbf{在一以下}，
故\textbf{质量谱是反对更高秩代数的证据}。
这是对代数本身的第一个检验，而它来自粒子谱。
\leanline{Particle.observed\_disfavours\_higher\_rank}
\end{mdframed}

\begin{remark}[代数\emph{不}定的东西]
速率的\emph{数值}。把根的半和换算成"每单位对数能量的密度"需要一个根的归一化，
而框架不提供。故 $\rho\approx 0.86$ 仍是测量值，预言是关于\textbf{离散度}而非密度。
形式陈述：混合的 $\mathrm{CV}^2$ 在两个分量均值同时缩放下不变。
\leanline{Particle.rate\_value\_not\_fixed}
\end{remark}

\subsection{禁闭荷是二值的，这结清了 \texorpdfstring{\texttt{ColorAudit}}{ColorAudit}}

第 \ref{sec:charges} 节给每个点缺陷两个拓扑标签——某个有限群 $B$ 中的块荷，
与整数刺猬荷——并让 $B$ 自由。

\textbf{代数把它定死了。}序参量是标度轴，因标度是比值而无向（第 \ref{sec:axis} 节），
故它是缩放扇区中过原点的一条直线，即实射影空间；其环路类构成 $\mathbb Z/2$。

\begin{theorem}[禁闭荷恰取两个值，且三值不可能]\label{thm:blk2}
$|B| = 2 \neq 3$。
\leanline{Particle.block\_two\_valued, block\_cannot\_be\_three\_valued,
block\_self\_inverse}
\end{theorem}

第 \ref{sec:coloraudit} 节把"此荷给 2 而数据要 3"记为一个未获解释的失败。
\textbf{它不再是未获解释的}：块群是 $\mathbb Z/2$，因为序参量是射影的，所以 3 从来就不在供给里。
色\textbf{确定}不是这个量子数，而且没有任何修补能让它是——
这正是第 \ref{sec:emergence} 节说它应该是的：一个截面的特征。

\subsection{一个粒子是什么}

\begin{mdframed}[linecolor=axcol,linewidth=1pt,backgroundcolor=axcol!5]
一个粒子 $=$ 一个\textbf{阈值} $\in\mathbb R$、一个 $\mathbb Z/2$ \textbf{标签}、
一个\textbf{整数荷} $\in\mathbb Z$。\textbf{框架不提供第四个槽位。}
\leanline{Particle.no\_further\_label, particle\_label\_structure}
\end{mdframed}

复合把它们相加，选择定则由此直接给出：

\begin{theorem}[两个奇复合成偶]\label{thm:odd}
两个 $\mathbb Z/2$-奇缺陷复合出平凡块标签。仅凭群，无需动力学。
\leanline{Particle.two\_odd\_make\_even}
\end{theorem}

\begin{theorem}[整数荷在任何分裂中精确守恒]\label{thm:qcons}
这是一个恒等式，故没有任何动力学能违反它。
\leanline{Particle.hedgehog\_conserved\_in\_splitting}
\end{theorem}

整数荷是非禁闭的（定理 \ref{thm:onecharge}）且精确守恒。
在观测到的世界里，恰有\emph{一个}内禀量子数同时具备这两条性质——\textbf{电荷}：
色被禁闭，弱同位旋被破缺，重子数与轻子数有反常。

\begin{mdframed}[linecolor=red!60!black,linewidth=0.9pt,backgroundcolor=red!3]
\textbf{这是结构上的吻合，而且被作为吻合陈述，不是对电磁学的推导。}
\end{mdframed}

\begin{remark}[诚实的界限]
没有质量、没有耦合常数、没有代结构，也没有任何理由给出整数荷的具体取值。
$\mathbb Z/2$ 标签具有统计标签的\emph{形状}，但此处没有任何东西推导出自旋
或自旋--统计联系，故它\textbf{不被称作自旋}。
"射影序参量的 $\pi_1$ 是 $\mathbb Z/2$"是标准拓扑，与秩一分类同样属于引用。
\end{remark}

\section{没有基本粒子}\label{sec:emergence}

第 \ref{sec:color} 节试图产生色，第 \ref{sec:coloraudit} 节表明它失败得不近。
那次失败不是技术性的。\textbf{整个目标就瞄错了}：
它拿了标准模型的本体论——一张固定的物种表，每个带内禀标签——然后试图重现它。
\emph{那是把问题也借了过来}，而框架并没有那个问题。

A3 说的是另一件事，而且从一开始就在公理里：$\varepsilon\cdot s=1$。
能标\emph{就是}逆长度标度。所以"存在什么"不是世界的性质，
而是\textbf{世界在某个分辨率下}的性质。改变分辨率，你不是对同一批粒子知道得更多——
\textbf{你在看另一批}。

这也正是本框架中"统一"的含义。引力、$\hbar$、粒子内容不是三样待调和的东西，
而是同一个变量的三种读法：

\begin{itemize}[leftmargin=1.6em]
\item 引力是标度的不均匀性（第 \ref{sec:geom}、\ref{sec:conn} 节）；
\item $\hbar$ 是沿标度迈一步的步长（第 \ref{sec:exact} 节）；
\item \textbf{粒子内容是标度上的一个截面}——本节。
\end{itemize}

\begin{theorem}[没有完整的清单]\label{thm:nolist}
若阈值无上界——A2 允许这一点，因为标度是单位元、永不为零，
故没有最小长度来给分辨率封顶——则\textbf{对每一个标度都有东西未被分辨}，
且内容严格增长无穷多次。

\textbf{于是根本不存在"那一份"基本粒子清单。}
问"基本粒子有哪些"预设了一个与标度无关的答案，而公理不提供这样的答案。
\leanline{Emergence.content\_never\_complete, always\_more\_above,
no\_fundamental\_list}
\end{theorem}

\begin{theorem}[质量是位置，不是属性]\label{thm:masspos}
质量与阈值互相决定。故质量\textbf{不是附着在物种上的参数}，
而是该物种\emph{出现在}标度轴上的坐标。
\leanline{Emergence.mass\_iff\_threshold, Emergence.mass\_mono}
\end{theorem}

这一条同时回溯解释了两个早先的发现。第 \ref{sec:audit} 节发现质量律不被拓扑决定——
\emph{当然不}：质量根本不是缺陷的属性，而是缺陷出现的地方。
而有效场论"每换一个能标就得换框架"的尴尬也不是尴尬：
定理 \ref{thm:eft} 说阈值之间的描述是\emph{精确}的，本节说没有最后一个阈值。
\textbf{理论只是在如实报告自己的结构。}
\leanline{Emergence.content\_eq\_of\_no\_threshold,
Emergence.content\_changes\_only\_at\_threshold}

\begin{mdframed}[linecolor=red!60!black,linewidth=0.9pt,backgroundcolor=red!3]
\textbf{放弃了什么。} 色、代结构、Higgs 以及标准模型表的其余部分\textbf{不是本报告的目标}。
它们是\emph{一个截面}的特征，而框架没有理由去重现某个截面的偶然内容。

应当向框架索取的是\textbf{阈值结构}——内容在哪里改变——
那是另一种、而且可检验的预言。第 \ref{sec:color} 节的色解释就此撤回，
其数学作为"简并块的结构"保留。
\end{mdframed}

\clearpage
\part{第六部分\quad 量 子 扇 区}

\begin{mdframed}[linecolor=axcol,linewidth=1pt,backgroundcolor=axcol!5]
\textbf{此处的“量子”不是另一套理论，是同一个标度的另一个读数。}
$\hbar$ 不是外加常数，是\textbf{沿标度轴平移一步与“在哪里”不对易}的失配大小；
而几何的全部量子修正是标度梯度的对易子 $[\sigma_i,\sigma_j]$，
它\textbf{纯反对称}，故对第四部分的一切经典检验\emph{不可见}。
\end{mdframed}

\section{非交换基底：量子力学的回归}\label{sec:nc}

\textbf{A1 的基底是一个环，\emph{不}要求交换。}交换性是\emph{经典扇区}的性质，
不是基底的公理。这不是补丁，而是比补丁更强的事实：
\textbf{经典扇区的微分结构从来没有用到交换性}——
A1 要求的只有 Leibniz 律与混合偏导可交换，
而这两条对任意环中的\emph{内导子} $\mathrm{ad}_p=[p,\cdot]$ 逐字成立。

\begin{remark}[一处开发记录]
本报告初稿确曾把基底取为交换环，理由是它让曲率计算可以由 \texttt{ring} 战术封闭。
那是\textbf{工具驱动而非物理驱动}的选择，且若成立就等于从一开始把量子力学排除在外
（$[\hat x,\hat p]=i\hbar$ 恰恰是非交换的）。本节是那次修正的结果。

\emph{定量}引力链条曾长期仍跑在交换扇区上——那是最后一个大缺口，
已由第 \ref{sec:ncconf} 节关闭：曲率计算不再需要交换性，
而第一个修正被\emph{辨认}出来（是 $[\sigma_i,\sigma_j]$），不是被假设。
\end{remark}

\subsection{导数就是对易子}

\begin{theorem}[任意环中的 Leibniz 律]\label{thm:adleib}
$\mathrm{ad}_p(ab)=\mathrm{ad}_p(a)\,b+a\,\mathrm{ad}_p(b)$。
\leanline{Quantum.ad\_leibniz}
\end{theorem}

\begin{theorem}[Jacobi 恒等式]\label{thm:adjac}
$\mathrm{ad}_p\mathrm{ad}_q-\mathrm{ad}_q\mathrm{ad}_p=\mathrm{ad}_{[p,q]}$。
因此\textbf{两个内导子可交换，当且仅当其生成元可交换}。
\leanline{Quantum.ad\_jacobi, Quantum.ad\_comm\_of\_commute}
\end{theorem}

定理 \ref{thm:adleib}、\ref{thm:adjac} 合起来说明：A1 中的 $\dd_i$ 不必被假定，
它们\emph{就是}对易子；而“选定一个坐标系”在此获得精确含义——
\textbf{选定一组互相对易的可观测量}。经典时空的存在性因此不是先验的，
而是“存在足够大的对易子代数”这一条件的几何影子。

\subsection{经典微积分是对易子的定理}

只假设正则对易关系 $[P,X]=c$（$c$ 中心），不假设任何别的东西：

\begin{theorem}[幂律，由对易子导出]\label{thm:adpow}
\[
[P,\;X^{\,n+1}] \;=\; (n+1)\,X^{\,n}\,c .
\]
取 $c=-i\hbar$、$\dd:=(i\hbar)^{-1}\mathrm{ad}_P$，这读作 $\dd(x^{n+1})=(n+1)x^n$。
\leanline{Quantum.ad\_pow}
\end{theorem}

即：\textbf{微分学的幂律不是分析学的公理，而是关于非交换性的定理}。
经典扇区中被假定的导子，在这里被推导出来。

\subsection{A3 在非交换扇区只能成立到 $\hbar$}

A3 说 $\varepsilon\cdot s=1$。在交换代数中这是精确恒等式，给出经典几何。
在非交换代数中它\emph{不可能}精确成立，而偏差的定量形式就是不确定性关系。

\begin{theorem}[Robertson 不等式]\label{thm:robertson}
对对称算子 $A,B$ 与单位向量 $\psi$，
\[
\bigl\|\langle\psi,(AB-BA)\psi\rangle\bigr\| \;\le\; 2\,\|A\psi\|\,\|B\psi\| .
\]
\leanline{Quantum.robertson}
\end{theorem}

\begin{theorem}[Heisenberg 不确定性]\label{thm:heisenberg}
若 $(AB-BA)\psi=i\hbar\,\psi$，则
\[
|\hbar| \;\le\; 2\,\|A\psi\|\,\|B\psi\| ,
\qquad\text{即}\qquad \Delta A\cdot\Delta B \;\ge\; \frac{\hbar}{2}.
\]
\leanline{Quantum.heisenberg}
\end{theorem}

\medskip
\begin{mdframed}[linecolor=axcol,linewidth=0.9pt,backgroundcolor=axcol!4]
\textbf{统一的陈述.} 标度—能动量对偶 $\varepsilon s=1$ 有且只有两种实现方式：
在\emph{对易}扇区它精确成立，给出经典几何与广义相对论（第 \ref{sec:geom} 节起）；
在\emph{非对易}扇区它只能成立到对易子，而那个对易子就是 $\hbar$，
给出量子力学。二者不是两个待调和的理论，而是同一条对偶在仅有的两个可用状态下的读数。
\end{mdframed}

\subsection{形变量子化：经典极限的严格形式化}\label{sec:deform}

上一小节只说明了对易子\emph{满足}经典微分公理。真正欠下的是反方向：
证明交换的标度代数确实是某个非交换代数的 $\hbar\to0$ 极限，
从而 $\dd_i$ 就是对易子而非仅仅类似于对易子。

取一阶形变 $A_\hbar=A[\hbar]/(\hbar^2)$，星积
\[
(f_0+\hbar f_1)\star(g_0+\hbar g_1)
= f_0g_0+\hbar\bigl(f_0g_1+f_1g_0+P(f_0,g_0)\bigr),
\]
其中 $P$ 是 $A$ 上任意双导子。

\begin{theorem}[形变的存在与结合律]\label{thm:starassoc}
对\emph{任意}双导子 $P$，星积结合。双导子的两条 Leibniz 律恰是一阶
Hochschild 上圈条件。
\leanline{Deformation.star\_assoc}
\end{theorem}

\begin{theorem}[经典极限]\label{thm:starcl}
$\hbar$ 的零次项恰为 $A$ 的交换乘法，无任何近似。
\leanline{Deformation.star\_classical}
\end{theorem}

\begin{theorem}[对应原理，精确成立]\label{thm:starcomm}
对易子无经典部分，其 $\hbar$ 系数恰为 $P$ 的反对称化：
\[
\hbar^{-1}[f,g]_\star=\{f_0,g_0\}.
\]
\textbf{这里没有极限可取——一阶上该恒等式是精确的。}
\leanline{Deformation.star\_commutator}
\end{theorem}

\begin{theorem}[Poisson 括号是导子]\label{thm:poissonleib}
$\{\cdot,g\}$ 对交换环 $A$ 满足 Leibniz 律。
\leanline{Deformation.Biderivation.poisson\_leibniz\_left}
\end{theorem}

定理 \ref{thm:starcomm}、\ref{thm:poissonleib} 合起来就是所欠的那一步：
\textbf{A1 中被假定的导子 $\dd_i$，正是 $A_\hbar$ 的内导子除以 $\hbar$}；
非交换性是 $\hbar$ 的一阶效应，在 $\hbar=0$ 处消失，留下引力扇区的交换几何。

\begin{theorem}[正则对易关系]\label{thm:ccr}
若 $X$ 只沿方向 $p$ 变化、$\Pi$ 只沿 $q$ 变化，则 $\{\Pi,X\}=1$，
故在形变代数中 $[\Pi,X]_\star=\hbar$。取 $\hbar=i\hslash$ 即 $[\Pi,X]=i\hslash$。
\leanline{Deformation.ccr\_of\_canonical}
\end{theorem}

至此闭环：定理 \ref{thm:ccr} 由经典标度代数产生 CCR，
而定理 \ref{thm:adpow} 由 CCR 复原经典微积分的幂律。
引力扇区的交换微分结构与量子力学的正则对易关系，
是同一个一阶形变的两端。

\subsection{把"一阶"这个限制取消掉}\label{sec:exact}

上一小节的形变只做到一阶。那是诚实的，但一阶截断是个不该安放地基的位置。
而且它\emph{根本不必要}——框架里本来就有一个\emph{精确}的非交换来源。

问：理论实际拥有的是哪两个操作？一个是读取\emph{此处}的量：乘以 $a$。
另一个是移到\emph{相邻标度}：位移 $T$。二者不对易，而不对易量不是小参数，
是一个有限的、精确的量——标度差本身：
\[
[\,U,\;M_a\,] \;=\; M_{Ta-a}\cdot U .
\]

\begin{mdframed}[linecolor=axcol,linewidth=0.9pt,backgroundcolor=axcol!4]
\textbf{位置与标度不对易。} 这是 A3（$\varepsilon s=1$）的算符形式：
分辨"在哪里"与分辨"在什么标度上"是不相容的操作，
而 $\hbar$ 是\emph{标度步长}，不是一个形式符号。
\end{mdframed}

\begin{theorem}[交换关系与对易子的闭形式]\label{thm:exchange}
$U M_a = M_{Ta}U$，且 $[U,M_a]=M_{\delta a}U$，其中 $\delta a := Ta-a$。
\textbf{全程精确，无任何展开。}
\leanline{Crossed.shift\_mul\_eq, Crossed.commutator\_eq}
\end{theorem}

\begin{theorem}[交换性 $\iff$ 标度步长为零]\label{thm:commiff}
代数交换当且仅当标度位移平凡。
\textbf{经典几何不是为方便而取的极限，而是"沿标度什么也没动"的那个情形。}
\leanline{Crossed.comm\_iff\_shift\_trivial}
\end{theorem}

\begin{theorem}[有限步长下的精确 Leibniz 律]\label{thm:twisted}
\[
\delta(ab)=\delta(a)\,Tb + a\,\delta(b).
\]
规则是\emph{扭曲}的：第二个因子在位移后的标度上读取。此式对任意大的步长精确成立。
当扭曲平凡（$Tb=b$）时，它退化为 A1 所假设的普通 Leibniz 律。
\leanline{Crossed.delta\_twisted\_leibniz, Crossed.delta\_leibniz\_of\_untwisted}
\end{theorem}

于是经典微分结构\textbf{不是}"$\hbar$ 的一阶截断"，而是一条精确关系在
\emph{零标度步长}处的取值。上一小节的限制就此被取消，而不是被绕过。

\begin{remark}[仍在的限制]
本节给出的是位移与乘法生成的算符代数，\emph{未}构造完整的交叉积、
未讨论其表示论，也未把 $\hbar$ 的数值与任何具体步长联系起来。
\end{remark}

\section{闭合量子引力回路}\label{sec:ncconf}

唯一剩下的大缺口是：\emph{定量}引力链条跑在交换扇区上。
定理 \ref{thm:commcost} 已经证明代价是什么——在交换环上 $\ad$ 恒为零，
故该扇区\textbf{完全没有输运曲率}，第 \ref{sec:geom} 节的一切曲率只能从公设的 $\kron$ 进来。
在这一步被重做之前，没有任何量子引力陈述可讲。

\textbf{本节重做它。}场景是一个\textbf{非交换}环，带 $n$ 个相互对易的导子——
导子仍然对易，因为那是一张图的选择（定理 \ref{thm:noflatcomm}），
但它们产生的\emph{值}不再对易。

\subsection{原封不动的部分}

Kronecker 符号是 $0$ 与 $1$ 的像，故\textbf{中心}；
凡只是把它们挪来挪去的步骤全都不变。
Hessian 的对称性只来自 $d_{\text{comm}}$，Christoffel 的对称性只来自 $\kron$ 的对称性。
\textbf{联络是同一个对象。}

\subsection{改变的只有一件事}

\begin{mdframed}[linecolor=axcol,linewidth=1.2pt,backgroundcolor=axcol!6]
\begin{theorem}[全部量子引力修正，一行]\label{thm:ncdefm}
\[
\Defm_{ij} - \Defm_{ji} = -2\,[\sigma_i,\sigma_j]
\]
\textbf{几何的全部量子引力修正，就是不同方向上标度梯度的对易子。}
\leanline{NCConformal.Defm\_antisymm\_part}
\end{theorem}
\end{mdframed}

它不是加进拉氏量的一项；它是\textbf{标度无法同时在两个方向上有确定梯度}——
正是第 \ref{sec:exact} 小节做成精确、第 \ref{sec:nc} 节变成不确定度的那件事。

\begin{theorem}[对称 $\iff$ 梯度对易]\label{thm:ncclass}
故\textbf{经典几何不是近似，而是对易子恰好消失的那个轨迹}。
\leanline{NCConformal.Defm\_symm\_iff\_commutator,
Defm\_symm\_iff\_commute}
\end{theorem}

主收缩获得同样的修正，且再无其他。经典情形下
$\sum_m \Gamma^a_{cm}\Gamma^m_{be}$ 含有 $2\delta_{ac}\sigma_b\sigma_e$；
非交换情形下这个 $2\cdot$ 变成\textbf{反对易子}，交叉项保持次序，
并出现\emph{一个}全新项 $\delta_{be}[\sigma_a,\sigma_c]$。

\begin{theorem}[非交换主收缩及其经典极限]\label{thm:ncmaster}
逐项算出；对易子消失时精确回到第 \ref{sec:geom} 节的结果。
\leanline{NCConformal.sum\_Chr\_Chr\_full\_nc,
sum\_Chr\_Chr\_full\_classical}
\end{theorem}

\subsection{预言}

\textbf{修正是反对称的。}而经典检验——偏折、进动、红移——
都由 Ricci 的\emph{对称}部分算出。

\begin{theorem}[对称部分完全未被修正]\label{thm:ncsym}
形变张量的对称部分\textbf{恰是}经典的那个，任何阶都没有对易子修正；
全部修正住在反对称部分。
\leanline{NCConformal.symmetric\_part\_uncorrected,
correction\_is\_pure\_antisymmetric}
\end{theorem}

\begin{mdframed}[linecolor=axcol,linewidth=1.2pt,backgroundcolor=axcol!6]
\textbf{故领头阶量子引力效应对一切经典检验\emph{不可见}}，
只出现在对反对称几何部分敏感的观测量中。
那是一种类挠率耦合；而由第 \ref{sec:axes} 节，
框架的转动标签恰是标度方向的楔积——\textbf{同一个对象}。于是预言是：

\begin{center}
\textbf{量子引力修正首先出现在\emph{自旋--几何耦合}中，而不是轨道中。}
\end{center}
\leanline{NCConformal.quantum\_correction\_couples\_to\_rotation,
correction\_vanishes\_on\_commuting}
\end{mdframed}

\begin{remark}[量级，诚实地说]
第 \ref{sec:exact} 小节把 $[\sigma_a,\sigma_b]$ 的尺度定为标度步长。
相对经典项，修正就是那个量级，故在太阳系中\textbf{完全可忽略}；
此处\textbf{不预言任何现有仪器可测的效应}。
所主张的是第一个修正的\emph{结构}，以及它能与不能出现在何处。
\end{remark}

\begin{remark}[仍然没有的]
没有引力子、没有散射振幅、没有黑洞熵、没有短距完备化。
\textbf{被闭合的正是被点名的那一个}：曲率计算不再需要交换性，
而第一个修正被\emph{辨认}出来而非假设。
\end{remark}

\clearpage

\clearpage
\part{第七部分\quad 标度流与内容}

\section{重整化群是定理，不是方法}\label{sec:rg}

\subsection{标度流}

\begin{definition}[标度流]
耦合空间 $S$ 上的标度流是 $\Phi_t:S\to S$，满足 $\Phi_0=\mathrm{id}$、
$\Phi_{t+u}=\Phi_t\circ\Phi_u$。群律由 A5 强制：基准先移 $u$ 再移 $t$，
就是移 $t+u$，没有别的可能。
\leanline{RG.ScaleFlow}
\end{definition}

\begin{proposition}[不动点即标度不变]
若 $\Phi_t g=g$ 对一切 $t$ 成立，则任意可观测量 $O$ 满足 $O(\Phi_t g)=O(g)$。
\leanline{RG.ScaleFlow.observable\_const\_at\_fixedPoint}
\end{proposition}

\subsection{自相似只能是幂律}

\begin{theorem}[Callan–Symanzik]\label{thm:cs}
若无量纲可观测量满足 $\dfrac{dO}{dt}=\gamma O$，则
\[
O(t) = O(0)\,e^{\gamma t}.
\]
特别地 $\gamma=0$（不动点）给出严格标度不变。
\leanline{RG.callan\_symanzik, RG.scale\_invariant\_of\_fixedPoint}
\end{theorem}

反常量纲 $\gamma$ 是标度协变行为的\emph{唯一}可能形式。
值得注意，这与第 \ref{sec:de} 节的耗散律是同一条微分方程：
\textbf{跑动的耦合与膨胀的宇宙，是同一个数学在两个标度上的读数}——
这正是本纲领所主张的自相似性。


\subsection{阈值之间，描述是\emph{精确}的}\label{ss:exact}

这一小节此前只在 Lean 中，未进正文，而后文（第 \ref{sec:emergence}、\ref{sec:slice} 节）
的整条论证都依赖它。

\begin{theorem}[无阈值区间上内容不变]\label{thm:exactbetween}
若区间 $[t,t']$ 内没有任何阈值，则活跃内容集合在其上\emph{恒等}，
故 $\beta$ 系数在其上是\textbf{常数}，而有效描述在该区间上不是近似，是\textbf{精确}。
\leanline{RG.active\_eq\_of\_no\_threshold, RG.beta\_const\_between\_thresholds,
RG.active\_mono, Emergence.resolved\_mono}
\end{theorem}

\begin{mdframed}[linecolor=axcol,linewidth=1pt,backgroundcolor=axcol!5]
\textbf{这把“有效场论要在每个能标换一套描述”从窘境变成结构。}
描述之所以要换，不是因为它不够好，而是因为\textbf{内容在阈值处真的变了}；
而在两个阈值之间它\emph{一点都不近似}。
第 \ref{sec:emergence} 节接着证明没有最后一个阈值，
故这件事没有尽头——那不是方法的缺陷，是理论在如实报告自己的结构。
\end{mdframed}

\begin{theorem}[内容只增不减]\label{thm:contentmono}
更好的探针不会丢失结构：$t\le t' \Rightarrow$ 活跃集合单调包含。
故“在能量 $E$ 上存在什么”至少是良序的——尽管第 \ref{sec:native} 节表明
它在\emph{单点}上不良定，内容属于区间。
\leanline{RG.active\_mono, RG.active\_eq\_univ\_of\_ge, RG.active\_finite\_range}
\end{theorem}

\subsection{有效场论塔是定理而非窘境}

设激发 $i$ 携带对数阈值 $\mu_i$。在对数能标 $t$ 上可用的激发集合为
$\mathcal{A}(t)=\{i:\mu_i\le t\}$。

\begin{theorem}[EFT 塔]\label{thm:eft}
\begin{enumerate}[label=(\roman*),leftmargin=2em]
\item \emph{单调性}：$t\le t' \Rightarrow \mathcal{A}(t)\subseteq\mathcal{A}(t')$；
      新物理只随标度升高而出现，从不消失。
      \lean{RG.active\_mono}
\item \emph{阈值间恒定}：若 $(t,t']$ 中无阈值，则 $\mathcal{A}(t)=\mathcal{A}(t')$。
      有效描述在阈值之间是\emph{精确}的，而非近似。
      \lean{RG.active\_eq\_of\_no\_threshold}
\item \emph{有限性}：有限激发内容只给出有限多个不同的有效理论。
      \lean{RG.active\_finite\_range}
\end{enumerate}
\end{theorem}

因此，“不同能标下必须更换理论框架”不是有效场论的缺陷，
而是一个标度分层世界在如实报告自己的结构。

\subsection{$\beta$ 函数从哪里来}\label{sec:beta}

第 \ref{sec:qcd} 节把单圈系数 $b$ 当作输入，是从规范理论借来的。这一步可以往回推一层。

A5 说基准平移是几何的对称性而非描述的对称性：耦合必须补偿。
而\emph{需要被补偿的}只能是实际在场的内容，即活跃激发。
所以补偿率只能是活跃集合的泛函，别无依赖。

\begin{theorem}[$\beta$ 系数的形状]\label{thm:beta}
设 $\beta$ 系数为活跃内容的任意泛函 $\Phi(\mathcal{A}(t))$。则
\begin{enumerate}[label=(\roman*),leftmargin=2em]
\item 阈值之间它\emph{恒定}——那里的跑动精确是单系数的，而非近似；
      \lean{RG.beta\_const\_between\_thresholds}
\item 它只取有限多个值，故是一个阶梯函数，跳变点恰在质量阈值上；
      \lean{RG.beta\_finite\_range}
\item 越过全部阈值后它稳定到完整内容的值。
      \lean{RG.beta\_eq\_of\_all\_active}
\end{enumerate}
\end{theorem}

QCD 的系数正是这样：在每个夸克阈值处跳变。本框架\emph{定出了这个形状}，
但\textbf{不定出数值}——$\Phi$ 的具体形式（如 $11N_c-2N_f$）仍是输入。
这一步只把借来的东西从"一个数"缩小为"一个泛函"，没有消灭它。

\section{跑动不需要圈图}\label{sec:running}

第 \ref{sec:qcd} 节假设了 $dg/dt=-bg^3$ 并从规范理论取来 $b$。
这是本报告最后一处大的借用，值得问一句：框架非借不可吗？

不必。放下"跑动来自求和图"的成见，问耦合在这里\emph{是什么}。
处于对数标度 $t$ 的探针有一个分辨率；它响应它能分辨的缺陷，
而由第 \ref{sec:defect} 节，那些是带确定对数标度的\emph{离散}对象。
所以累积响应是一个\textbf{计数}。

再加上 A5：标度协变的世界没有优待的对数标度，故缺陷不能在任何地方聚集——
它们\emph{每单位对数标度的密度是常数}。常密度在 $[t_0,t]$ 上积分即线性。于是
\[
g^{-2}(t)=g^{-2}(t_0)+\kappa\rho\,(t-t_0),
\]
这\emph{正是}单圈形式，其中 $b=\kappa\rho/2$。

\begin{theorem}[计数给出线性跑动]\label{thm:counting}
累积响应在对数标度上的斜率恒为 $\kappa\rho$；该式与单圈解形式相同，
系数被认成"每单位对数标度的缺陷密度 $\times$ 单缺陷响应"。
\leanline{Running.hasDerivAt\_invSqCoupling, Running.invSqCoupling\_eq\_one\_loop}
\end{theorem}

\begin{theorem}[渐近自由 $\iff$ 正的缺陷密度]\label{thm:afdensity}
响应在高对数标度处趋零，当且仅当可分辨的缺陷密度为正。
\textbf{分辨出越多结构，响应被稀释得越厉害}——这就是渐近自由在本框架下的物理内容，
陈述与证明中都没有出现规范群、没有 $11N_c-2N_f$、也没有任何图。
\leanline{Running.asympt\_free\_iff\_pos\_density}
\end{theorem}

\textbf{线性对数跑动是标度不变的缺陷密度的推论，不是单圈图的推论。}
借来的系数由此变成一个缺陷密度——一个计数，而 $11N_c-2N_f$ 本来也是一个计数。
仍未导出的是密度的\emph{数值}。

\section{强相互作用：从无到有的一个标度}\label{sec:qcd}

A5 说没有任何标度被优待，理论完全无量纲。但强相互作用显然\emph{有}标度——
强子有确定的质量。若本框架无法交代这一点，它恰恰在标度物理最难的地方失败。

它交代得了，而机制正是标准模型中唯一纯粹关于标度的现象：
\textbf{量纲嬗变（dimensional transmutation）}。一个纯数耦合，
在 A5 的基准流下跑动，会分泌出一个公理中根本不存在的不变标度 $\Lambda$。
\textbf{标度不变性不是被强行破坏的，而是被一个标度协变方程的\emph{解}破坏的}——
在一个没有任何有量纲参数可供破缺的理论里，这是唯一可能的破缺方式。

由单圈跑动 $dg/dt=-bg^3$（$b>0$）出发：

\begin{theorem}[跑动的精确解]\label{thm:running}
$g^{-2}(t)=g^{-2}(0)+2bt$。未用任何微扰展开；对给定流该式精确。
\leanline{QCD.running\_inv\_sq}
\end{theorem}

\begin{theorem}[渐近自由]\label{thm:asympt}
$g\to0$ 当 $t\to\infty$：强扇区在标度极小处变得自由。
由 A3，小标度即高能，这正是所要求的行为。
\leanline{QCD.asymptotic\_freedom}
\end{theorem}

\begin{theorem}[量纲嬗变]\label{thm:transmut}
\[
\Lambda \;=\; \exp\!\Bigl(t-\frac{g^{-2}(t)}{2b}\Bigr)
\]
沿流\textbf{恒定}。一个仅以纯数为输入的理论产生了一个标度，
而由 A5，任何基准选择都无法移动它。
\leanline{QCD.lambda\_invariant, QCD.lambda\_pos}
\end{theorem}

\subsection{对谱的推论与检验}

$\Lambda$ 是强扇区唯一的标度。由 A3，质量即逆标度，故每个强子质量必是
纯数乘 $\Lambda$。

\begin{theorem}[质量比无自由参数]\label{thm:massratio}
任意两个强子质量之比是两个纯数之比，与 $\Lambda$ 无关，
因而与理论的一切输入无关。
\leanline{QCD.hadron\_mass\_ratio, QCD.spectrum\_rigid}
\end{theorem}

这是本框架对强扇区最尖锐的预言，也正是\textbf{胶球}是检验它的正确场所的原因：
胶球质量不含任何夸克质量，$\Lambda$ 确实是唯一在场的标度。

\subsection{与格点数据的正面对撞：一个失败}

淬火格点 QCD 给出（MeV）：$0^{++}\approx1710$、$2^{++}\approx2390$、
$0^{-+}\approx2560$。第 \ref{sec:defect} 节的缺陷机制预言\emph{等比}塔，
即相邻质量比相等。检验：

\[
\frac{2390}{1710}=1.398,\qquad \frac{2560}{2390}=1.071 .
\]

\begin{mdframed}[linecolor=red!60!black,linewidth=0.9pt,backgroundcolor=red!3]
\begin{theorem}[最简等比塔被证伪]\label{thm:towerfail}
两个相邻比值之差大于 $0.3$。
\textbf{横跨全部胶球态的单一等比塔被格点数据排除。}
\leanline{QCD.geometric\_tower\_fails, ratio$_1$\_value, ratio$_2$\_value}
\end{theorem}
\end{mdframed}

这否定的是定理 \ref{thm:spectrum} 所基于的\emph{质量律假设}（见第 \ref{sec:audit} 节），而非缺陷图景或框架本身。我们如实记录，并已把结论削弱到它应有的强度。
幸存下来的是两条较弱但仍可检验的陈述：
（i）塔可能只存在于\emph{固定} $J^{PC}$ 通道\emph{之内}；
（ii）定理 \ref{thm:massratio} 的无参数比值陈述完全未受影响。

作为待计算的目标：取 $\Lambda\approx332$ MeV，最轻胶球位于
$m/\Lambda\approx5.15$——一个本框架最终必须\emph{算出}而非拟合的纯数。
\leanline{QCD.lightestGlueballInLambda\_value}

\section{广延的代价、转动是什么、以及为什么是三}\label{sec:slice}

三个问题，而它们互相咬合。

\subsection{标签泛滥是覆盖范围的代价，不是结构的证据}

值得认真对待的警告是：标准模型的那套装置——色、代、弱同位旋、禁闭机制——
可能根本不是结构，而是把一个\emph{截面局部}的描述广延到一段标度范围上时
被\textbf{迫}写下的东西。这是可形式化的，此处证明。

第 \ref{sec:emergence} 节说标度 $t$ 处的内容是 $\{k \mid \mu_k \le t\}$，
第 \ref{sec:rg} 节说描述在阈值之间是\emph{精确}的。
故单一标度上有效的描述只需要该标度的内容；
而跨范围有效的描述必须携带在下端\textbf{不可分辨}的内容的标签。

\begin{theorem}[跨越阈值即强制新标签]\label{thm:newc}
范围内有任何阈值，广延描述就必须携带局部描述不需要的标签；
而这个超出量随范围\textbf{单调增长}。
反之，不含阈值的范围代价为零。
\leanline{Slice.newContent\_nonempty\_of\_threshold, newContent\_mono,
newContent\_empty\_of\_no\_threshold}
\end{theorem}

\begin{theorem}[没有在所有标度上都够用的标签集]\label{thm:nouniv}
任何标度之上都还有未分辨的内容，故任何固定标签集终将不足。
\textbf{在所有标度上都有效的形式体系并不存在}，而它的不存在是定理，不是谁的巧思不够。
\leanline{Slice.no\_finite\_universal\_labelling, label\_count\_measures\_range}
\end{theorem}

\begin{mdframed}[linecolor=axcol,linewidth=1pt,backgroundcolor=axcol!5]
\textbf{一个形式体系携带的标签数量，度量的是它覆盖的对数范围，不是物理的深度。}
用一套形式覆盖三十个 e-fold，就必须为这三十个 e-fold 中任何地方出现过的东西备好标签，
而其中绝大多数在任一给定标度上并未实现。
这恰是壳外内容、规范冗余、以及"只在虚过程中存在的物种"的形状。
\end{mdframed}

这并\emph{不}反驳标准模型。它说的是：标准模型的标签数是关于它\emph{范围}的证据，
而复现这个数目不是目标——这正是第 \ref{sec:emergence} 节主张过、本节证明了的事。

\subsection{转动的原生量是楔积——而这打破了前面}

框架自己的量很少：阈值、$\mathbb Z/2$ 标签、整数荷、标度量值。
转动不在其中，然而第 \ref{sec:coupling} 节生产它。它原生地是什么？

定理 \ref{thm:bbr} 已经回答：两个缩放生成的转动是 $v_iw_j - v_jw_i$，
即它们方向的\textbf{楔积}。故转动标签在这里\emph{不是}独立量子数——
它是标度结构的函数，没有自由度。

\begin{theorem}[楔积是真正的二形式，且平行时为零]\label{thm:wedge}
反称、双线性，且 $\text{wedge}\,v\,w = 0 \iff v \parallel w$。
\leanline{Slice.wedge\_eq\_bracket, wedge\_antisymm, wedge\_bilinear\_left,
wedge\_eq\_zero\_iff\_parallel}
\end{theorem}

\begin{theorem}[均匀花样不携带转动标签]\label{thm:uniax}
$\text{wedge}\,v\,v = 0$：一个矢量与\emph{自身}的楔积为零。
故\textbf{均匀}（不随位置变化的）标度花样转动标签恒为零。
\leanline{Slice.homogeneous\_carries\_no\_rotation, rotation\_requires\_biaxial}
\end{theorem}

\begin{mdframed}[linecolor=red!60!black,linewidth=0.9pt,backgroundcolor=red!3]
\textbf{更正（见第 \ref{sec:axes} 节）。}我原先把这条定理命名为
"单轴不携带转动"并据此提出一个两难。\textbf{那个读法是错的，在此撤回。}
定理讲的是一个矢量与\emph{自身}的楔积，故它说的是\textbf{恒定}花样不生成转动；
而"单轴"指标度只取一个不同的\emph{值}，完全是另一回事。
第 \ref{sec:axes} 节表明转动由花样的\emph{变化}携带，
从而绕数缺陷必然带转动，且更多轴根本不可得。
第 \ref{sec:particle} 节的标签因此无需限缩。
\end{mdframed}

\subsection{为什么是三个空间方向}

第 \ref{sec:algebra} 节导出了族而不是维数。这里是结构所能提供的唯一一条论证，
连同它的确切强度。

A4 只提供一种标签：\textbf{方向}。
第 \ref{sec:gauge} 节使一切可观测标签都是方向间的简并花样。
第 \ref{sec:coupling} 节使一切转动都是两个方向的楔积。
故缩放生成的转动\emph{当且仅当}楔积空间与方向空间同维时才带有方向标签：
\[
\dim \Lambda^2\mathbb R^k = k \iff \frac{k(k-1)}{2} = k \iff k = 3 .
\]

\begin{theorem}[三是唯一解]\label{thm:three}
$k(k-1)/2 = k$ 的唯一正解是 $k=3$。
$k>3$ 时楔积严格多于方向，故有转动不带方向标签；
$k<3$ 时转动扇区退化。
\leanline{Slice.wedge\_dim\_eq\_iff, three\_is\_unique,
more\_wedges\_than\_directions, fewer\_wedges\_than\_directions}
\end{theorem}

\begin{mdframed}[linecolor=axcol,linewidth=1pt,backgroundcolor=axcol!5]
\textbf{$k=3$ 是耦合在公理所提供的标签上闭合的唯一维数}，给出 $n = k+1 = 4$。
\end{mdframed}

\begin{remark}[这有多强]
算术是定理。而承载物理的那一步——框架\emph{要求}它生成的转动带方向标签，
而不只是\emph{允许}某些转动不带标签——是一个\textbf{判断}，
与第 \ref{sec:algebra} 节的秩一认定同类。
它被明写在外面并已登记，没有藏进任何证明。
\textbf{有它，$n=4$；没有它，$n$ 仍然自由。}
\end{remark}

值得一提的自洽：$k=3$ 也正是 $\mathfrak{so}(3)\cong\mathfrak{su}(2)$、
转动群获得双覆盖的地方——而那正是第 \ref{sec:particle} 节 $\mathbb Z/2$ 的来源。
\textbf{维数与那个标签同源。}

\section{寿命的动量依赖：$\gamma$ 就是两个标度之比}\label{sec:particles}

激发携带自己的时钟，其单位由静止能标 $m$ 设定；实验室使用另一套单位，
由总能量 $E=\sqrt{p^2+m^2}$ 设定。寿命是内部“滴答”的纯数目；
实验室报告的是同一数目在\emph{它的}单位下的读数，因此要乘以两个标度之比
\[
\gamma = \frac{E}{m}.
\]

\begin{theorem}[寿命]\label{thm:lifetime}
设 $m>0$、$\tau_0\ge 0$，$\tau(p)=\gamma(p)\,\tau_0$。则
\begin{enumerate}[label=(\roman*),leftmargin=2em]
\item $\gamma\ge 1$，故 $\tau(p)\ge\tau_0$；\lean{Particles.one\_le\_scaleRatio}
\item $p^2\le q^2 \Rightarrow \tau(p)\le\tau(q)$：寿命随动量单调增加；
      \lean{Particles.lifetime\_mono}
\item $\tau(0)=\tau_0$：静止时两个标度重合，归一化无余项。
      \lean{Particles.lifetime\_at\_rest}
\end{enumerate}
\end{theorem}

证明中没有使用任何关于衰变动力学的假设。时间膨胀在此不是“钟走慢了”，
而是 A3 的标度比被读了两次。

\section{暗物质：相互作用发生，观测失效}\label{sec:dm}

\begin{mdframed}[linecolor=axcol,linewidth=1pt,backgroundcolor=axcol!5]
\textbf{本节是\emph{两条}路径中的第一条，而两条做出\emph{相反}的预言。}
本节的低标度路径说：耦合小但\textbf{非零}，故足够匹配的探针终能看见，
且这些对象\emph{有}阈值，故\textbf{有质量谱}。
第 \ref{sec:cosmos} 节给出另一条——各向同性轨迹：
在任何标度上都\textbf{没有可耦合的标签}，也没有阈值，故\textbf{没有谱}。
\textbf{哪一条是世界所用，是实验问题；框架容纳两者，它不做选择。}
\end{mdframed}

探测器本身是坐落在某个能标 $\varepsilon_D$ 上的物理系统；由 A3，
这也是它的空间分辨率的倒数。它能登记的不是“一次相互作用”，
而是“一次把它自身状态推动了至少阈值 $\theta$ 的相互作用”。

\begin{definition}
单次激发的探测响应 $\;\mathrm{det}(\varepsilon,\varepsilon_D)=\varepsilon/\varepsilon_D$；
引力响应 $\;\mathrm{grav}(N,\varepsilon)=N\varepsilon$（总能量，对能量如何分配是盲的）。
\end{definition}

\begin{theorem}[相互作用而不可观测]\label{thm:detfail}
对任意 $\varepsilon_D>0$ 与任意阈值 $\theta>0$，存在 $\varepsilon_0>0$，使得
对一切 $0<\varepsilon<\varepsilon_0$，
\[
0 < \mathrm{det}(\varepsilon,\varepsilon_D) < \theta .
\]
即：耦合严格非零（相互作用确实发生），却严格低于阈值（观测必然失败）。
\leanline{Dark.detection\_failure}
\end{theorem}

\begin{theorem}[暗物质]\label{thm:dm}
对任意引力信号 $M>0$、任意探测器标度 $\varepsilon_D>0$、任意阈值 $\theta>0$，
存在标度 $\varepsilon>0$ 与多重度 $N>0$，使得
\[
N\varepsilon = M \quad\text{（引力上完全可见）},\qquad
0<\frac{\varepsilon}{\varepsilon_D}<\theta \quad\text{（非引力上不可见）}.
\]
\leanline{Dark.dark\_matter\_exists}
\end{theorem}

引力可见性与非引力不可见性因此并不冲突：它们是同一位形的两个不同泛函，
而标度层级可以把二者拉开任意大的倍数。用户直觉中“超大号、能量极低的粒子”
在此获得精确表述——关键不在于它是什么粒子，而在于\emph{探测器自身的能标远高于它}，
使得测量（而非相互作用）失效。

\section{暗能量：耗散率就是 $\Lambda$}\label{sec:de}

由 A6，宇宙不封闭，其能量标度耗散。由 A3，衰减的能量标度\emph{就是}增长的长度标度。
定义宇宙标度因子 $a:=1/\varepsilon$（它不是独立场）。

\begin{theorem}[耗散的历史唯一]\label{thm:diss}
若 $\dot\varepsilon=-\lambda\varepsilon$，则
\[
\varepsilon(t)=\varepsilon_0 e^{-\lambda t}.
\]
未假设任何拟设；指数形式是被强制的。
\leanline{Cosmology.dissipation\_solution}
\end{theorem}

\begin{theorem}[耗散即膨胀，且加速]\label{thm:desitter}
设 $\lambda>0$、$\varepsilon_0>0$。则
\[
a(t)=a_0e^{\lambda t},\qquad
\dot a=\lambda a \;(\text{即 } H=\lambda \text{ 为常数}),\qquad
\ddot a=\lambda^2 a>0 .
\]
即精确的 de Sitter 膨胀，全程没有引入任何宇宙学常数。
\leanline{Cosmology.open\_universe\_accelerates, hasDerivAt\_scaleFactor,
hasDerivAt\_hubble, acceleration\_pos}
\end{theorem}

结论：所谓“暗能量”在此框架下不是具有负压的物质，而是\emph{开放宇宙的耗散率}。
观测到的加速膨胀是一个倒数的二阶导数，$\Lambda$ 与 $\lambda^2$ 同阶。
把宇宙自身看作一个“因涨落而诞生、正在耗散的大粒子”，
其数学内容恰是定理 \ref{thm:diss}——与定理 \ref{thm:cs} 的跑动方程同型。

\section{时间箭头}\label{sec:entropy}

基底中没有任何东西区分时间方向：微观动力学是双射。破缺来自\emph{观测发生在某个标度上}，
而标度携带有限分辨率，即一个从微观位形到可分辨对象的\emph{非单射}映射。

设 $X$ 有限，$\pi:X\to Y$ 为粗粒化，$T$ 为 $X$ 上的置换（可逆微观动力学）。
定义饱和 $\mathrm{sat}_\pi(S)=\{x:\pi x\in\pi(S)\}$，观测演化
$F(S)=\mathrm{sat}_\pi(T(S))$。

\begin{theorem}[H 定理]\label{thm:h}
$|S|\le |F(S)|$，对一切 $S$ 成立。
\leanline{Entropy.entropy\_nondecreasing}
\end{theorem}

证明只用两步：微观步严格保持基数；粗粒化步只能变大。没有任何概率假设。

\begin{theorem}[不可逆性]\label{thm:arrow}
设 $F$ 满足 H 定理且在某个 $S_0$ 上严格产熵（$|S_0|<|F(S_0)|$）。
则\emph{不存在}同样满足 H 定理的 $G$ 使 $G\circ F=\mathrm{id}$。
\leanline{Entropy.no\_reverse\_law, Entropy.evolve\_no\_reverse\_law}
\end{theorem}

即：宏观方程无法被同类型的规律倒着跑。微观层面时间反演不变严格成立，
宏观层面它被摧毁——这正是 Boltzmann 方程不具时间反演不变性的确切地位。

\paragraph{正反物质不对称。} 我们只给出结构性结论，不宣称推出观测到的重子不对称度。
Sakharov 第三条件是“偏离热平衡”。在封闭系统中这是精心安排的插曲；
在 A6 之下宇宙永不停止耗散，因而该必要条件\emph{结构性地}恒被满足：

\begin{proposition}
若某荷量在一步演化中发生改变，则系统必不处于平衡（$F(S)\ne S$）；
严格产熵本身即构成偏离平衡。
\leanline{Entropy.charge\_change\_implies\_nonequilibrium,
Entropy.entropy\_production\_implies\_nonequilibrium}
\end{proposition}

\section{光：标度盲的类光锥}\label{sec:light}

被测二次型是记账二次型乘以标度因子：$Q_g=s^2Q_\kron$。由于 $s$ 是单位元，
乘以它既不能制造零也不能消灭零。

\begin{theorem}[光对标度是盲的]\label{thm:null}
\[
Q_g(v)=0 \iff Q_\kron(v)=0 .
\]
类光锥由记账结构单独决定；任何标度场（即任何引力场）都无法改变哪些方向是类光的。
两个使用任意不同标度场的观察者，对“何为类光”的判断完全一致。
\leanline{Light.isNull\_iff\_bare, Light.isNull\_scale\_invariant}
\end{theorem}

\begin{theorem}[物质不是盲的]\label{thm:massive}
若 $Q_\kron(v)$ 是正则元（特别地，$v$ 非类光），则由被测二次型可\emph{唯一确定} $s^2$。
\leanline{Light.scale\_determined\_of\_nonnull, Light.physForm\_ne\_bareForm}
\end{theorem}

两条定理合起来给出本框架完整的认识论：

\begin{mdframed}[linecolor=axcol,linewidth=0.9pt,backgroundcolor=axcol!4]
\textbf{光确定共形（因果）结构，且仅此而已；标度必须用物质去读。}
\end{mdframed}

因此，光的数学本质是几何中\emph{权重为零}的部分——A5 的基准平移作用于其上是平凡的
（\texttt{Light.null\_weight\_zero}），这正是“无质量”的形式内容。
它的实在意义是：光是唯一路径对标度不敏感的探针，因而是唯一能够
把局域单位的比较从一个事件无循环论证地传递到另一个事件的\emph{比较器}。
再结合定理 \ref{thm:weyl}（标度联络平坦），这种传递沿闭合回路无残差，
理论自洽。

\clearpage

\clearpage
\part{第八部分\quad 预言与数据}

\section{可检验的预言}\label{sec:pred}

\subsection{耗散指数与暗能量状态方程}

把 A6 从 $\dot\varepsilon=-\lambda\varepsilon$ 推广为幂律 $\dot\varepsilon=-\lambda\varepsilon^{q}$。
由 $a=1/\varepsilon$ 得 $H=\lambda a^{1-q}$，故 $\rho\propto H^2\propto a^{2(1-q)}$；
与连续性方程 $\rho\propto a^{-3(1+w)}$ 比较指数，得到

\begin{mdframed}[linecolor=axcol,linewidth=0.9pt,backgroundcolor=axcol!4]
\[
\boxed{\;w=\frac{2q-5}{3}\;}
\qquad\Longleftrightarrow\qquad
\boxed{\;q=\frac{3w+5}{2}\;}
\]
\end{mdframed}

\begin{theorem}[加速判据与状态方程]\label{thm:eos}
\begin{enumerate}[label=(\roman*),leftmargin=2em]
\item $2\dot\varepsilon^{\,2}-\varepsilon\ddot\varepsilon=\lambda^2\varepsilon^{2q}(2-q)$，
      故 $\ddot a>0 \iff q<2$。\lean{Predictions.powerlaw\_accel\_numerator,
      accelerates\_iff}
\item $w=-1 \iff q=1$（预言是尖锐的，不是通有的）。
      \lean{Predictions.eos\_eq\_neg\_one\_iff}
\item $w<-1/3 \iff q<2$：与 (i) 的加速判据一致。
      \lean{Predictions.eos\_lt\_third\_iff}
\item $w<-1$（幻影区）$\iff q<1$。\lean{Predictions.eos\_lt\_neg\_one\_iff}
\item 反演公式是双向的：$w\mapsto q\mapsto w$ 恒等。
      \lean{Predictions.eos\_dissipationExponent}
\end{enumerate}
\end{theorem}

这是本框架最容易被杀死的地方，应当明说：

\begin{mdframed}[linecolor=red!60!black,linewidth=0.9pt,backgroundcolor=red!3]
\textbf{可证伪性.} 纯指数耗散（$q=1$）预言 $w=-1$ \emph{精确}成立，无任何可调参数。
若观测确证 $w\ne-1$，则该情形被直接排除；此时理论并不崩塌，而是被数据
\emph{测定}：由 $q=(3w+5)/2$ 读出耗散指数。例如观测给出 $w=-0.95$，
则 $q=1.075$。\lean{Predictions.worked\_inversion}
\end{mdframed}

\subsection{代入常数：三项数值核对}

以下公式由本框架给出，常数取 SI 定义值或 CODATA 值，算术由 \texttt{norm\_num} 验证。

\begin{center}
\begin{tabular}{llll}
\toprule
量 & 本框架公式 & 计算值 & 实验值 \\
\midrule
$\mu$ 子寿命 & $\tau=\gamma\tau_0$ & $64.432\ \mu$s & $64.378\pm0.026\ \mu$s \\
Pound–Rebka & $z=gh/c^2$ & $2.4551\times10^{-15}$ & $(2.46\pm0.03)\times10^{-15}$ \\
暗能量密度 & $\rho_\Lambda=3\lambda^2/8\pi G$ & $5.845\times10^{-27}$ & $5.86\times10^{-27}$ \\
\bottomrule
\end{tabular}
\end{center}

\begin{itemize}[leftmargin=1.6em]
\item $\mu$ 子：CERN 储存环 $\gamma=29.327$、$\tau_0=2.19703\,\mu$s，
      由定理 \ref{thm:lifetime} 得 $\tau=64.432\,\mu$s，与测量符合到 $10^{-3}$。
      \lean{Predictions.muon\_matches\_experiment}
\item Pound–Rebka：由定理 \ref{thm:redshift} 的 $\varepsilon=e^{\Phi}$ 一阶展开，
      $h=22.5$ m 处 $z=2.4551\times10^{-15}$。
      \lean{Predictions.poundRebka\_value}
\item 暗能量：由定理 \ref{thm:desitter}，$\lambda$ 即渐近 Hubble 率
      $H_\infty=H_0\sqrt{\Omega_\Lambda}\approx1.8077\times10^{-18}\,$s$^{-1}$，
      代入 $\rho_\Lambda=3\lambda^2/8\pi G$ 得 $5.845\times10^{-27}$ kg/m$^3$。
      \lean{Predictions.darkEnergyDensity\_value}
\end{itemize}

需要说明：前两项是\emph{核对}而非新预言（框架落在已知物理的正确数值上）；
第三项在 $\lambda\leftrightarrow H_\infty$ 的识别下带有一定循环性，
其真正的预言内容是 $w=-1$ 这一\emph{形状}，而非该数值本身。

\subsection{反演公式}

\begin{theorem}[状态方程与耗散指数互为反演]\label{thm:eosinv}
$w=(2q-5)/3$，故 $w=-1 \iff q=1$；$w<-1/3$（加速）$\iff q<2$；
$w<-1$（幽灵）$\iff q<1$。给定 $w=-0.95$ 反演得 $q=1.075$。
\leanline{Predictions.eos\_eq\_neg\_one\_iff, eos\_lt\_third\_iff,
eos\_lt\_neg\_one\_iff, worked\_inversion, accelerates\_iff}
\end{theorem}

\begin{remark}[全部数值对照见第 \ref{sec:data} 节]
本报告只保留\textbf{一张}数值总表，在第 \ref{sec:data} 节。
此处以及第 \ref{sec:spectrum} 节曾各有一张部分重复的表，已合并。

那里同时说明\textbf{这些数的地位并不相同}：
偏折、进动、$\gamma$ 是从公理链条导出的，中间无假设；
$\mu$ 子与 Pound--Rebka 是时间膨胀与红移的算术核对，检验的是 A3 的读法；
$\rho_\Lambda$ 由 $\lambda=H_\infty$ 换算而来，而 $\lambda$ 本身是测量值——
它正是那个唯一的自由数。
\end{remark}

\section{可检验的预言与换算表}\label{sec:spectrum}

第 \ref{sec:emergence} 节说该向框架索取的不是物种表而是\textbf{阈值结构}。
本节索取它，得到一个可以、并且已经与数据对撞的预言。

\subsection{预言：阈值是标度不变过程}

A5 说没有哪个对数标度被优待。于是阈值不能在任何地方聚集：
它们由一个\emph{标度不变}的过程放置，即每单位对数标度速率恒为 $\rho$ 的过程。
常速率点过程的间距服从指数分布，而指数分布的变异系数恰为\textbf{一}：
\[
\mathrm{CV}(\ln m\ \text{的相邻间距}) = 1 .
\]

这条预言是尖锐的，而且能区分。\emph{等比塔}——
第 \ref{sec:defect} 节曾假设、第 \ref{sec:audit} 节已撤回的那个——间距\emph{相等}，
故 $\mathrm{CV}=0$。两个假设相距不能再远。

\begin{theorem}[等比塔的方差恒为零]\label{thm:eqzero}
若每个间距都等于 $c$，则总体方差恒等于零，故 $\mathrm{CV}=0$，没有任何自由度。
\leanline{Spectrum.equal\_gaps\_variance\_zero}
\end{theorem}

\subsection{数据：标准模型的十二个质量标度}

取标准模型十二个互不相同的质量标度（电子到顶夸克），作 $\ln m$ 的十一个相邻间距：

\begin{center}
\begin{tabular}{lll}
\toprule
 & $\mathrm{CV}^2$ & $\mathrm{CV}$ \\
\midrule
标度不变过程（预言） & $1$ & $1$ \\
\textbf{观测} & $\mathbf{0.875}$ & $\mathbf{0.935}$ \\
等比塔 & $0$ & $0$ \\
\bottomrule
\end{tabular}
\end{center}

\leanline{Spectrum.observed\_cvSq\_value, Spectrum.meanGap\_value}

十一个间距的样本很小，预言自身的抽样范围很宽——$\mathrm{CV}$ 的中央 $90\%$
从 $0.58$ 到 $1.28$——而观测值 $0.935$ 落在其第 $63$ 百分位，即正中间。

\begin{mdframed}[linecolor=red!60!black,linewidth=0.9pt,backgroundcolor=red!3]
\begin{theorem}[标准模型谱不是等比塔]\label{thm:notgeom}
观测间距方差为正，而等比塔方差恰为零。
\textbf{这是比第 \ref{sec:qcd} 节胶球比值检验独立得多、也尖锐得多的一次否证}——
它用的是整个观测谱，而非三个格点数字。
\leanline{Spectrum.observed\_not\_geometric,
Spectrum.observed\_within\_scale\_invariant\_range}
\end{theorem}
\end{mdframed}

阈值密度由平均间距给出：$\rho = 1/\langle\text{gap}\rangle \approx 0.86$ 每 e 倍。
\leanline{Spectrum.densityFromGaps\_value}

\begin{remark}[诚实的界限]
十一个间距是小样本，且十二个质量标度是人工挑选的（无质量态排除，中微子质量未知）。
与 $\mathrm{CV}=1$ 的符合是\emph{鼓励性的，不是决定性的}；
\textbf{决定性的是对 $\mathrm{CV}=0$ 的排除}。
\end{remark}

\subsection{换算表}

框架不预言标准模型的表，但它说明任何一次测量在它自己的变量里是什么意思。
以下每一条都把一个被测量换算为它所定住的框架参数。

\begin{center}
\begin{tabular}{lll}
\toprule
被测量 & 框架参数 & 换算 \\
\midrule
暗能量状态方程 $w$ & 耗散指数 $q$ & $q=(3w+5)/2$ \\
PPN 参数 $\gamma$ & 倒数条件 & 预言 $\gamma=1$（无自由参数）\\
渐近 Hubble 率 $H_\infty$ & 耗散率 $\lambda$ & $\lambda=H_\infty$，$\rho_\Lambda=3\lambda^2/8\pi G$ \\
$\Lambda_{\mathrm{QCD}}$ & 嬗变标度 & $\Lambda=\mu\,e^{-1/2bg^2}$（沿流不变）\\
跑动系数跳变 $\Delta b$ & 新增可分辨内容 $\Delta N$ & $\Delta N=\Delta b/\kappa$ \\
质量 $m$ & 阈值 $\mu$ & $\mu=\ln m$（互为逆）\\
平均对数间距 & 阈值密度 $\rho$ & $\rho=1/\langle\text{gap}\rangle$ \\
\bottomrule
\end{tabular}
\end{center}

\leanline{Spectrum.thresholdOfMass, massOfThreshold, mass\_threshold\_inverse,
qFromW, rhoLambdaFromLambda, deltaNFromDeltaB, densityFromGaps,
density\_gap\_inverse}

\subsection{对照表已并入第 \ref{sec:data} 节}

此处原有一张数值对照表，与第 \ref{sec:pred} 节的表部分重复，且其暗能量一行已过时
（它写作“$w=-1$（\emph{若} $q=1$）……与 $-1.03\pm0.03$ 符合”；
而第 \ref{sec:cosmos} 节此后证明 A5 \textbf{强制} $q=1$，
第 \ref{sec:data} 节则记录 DESI 给出 $2.8$--$4.2\sigma$ 的\emph{反向}偏好）。
\textbf{全部数值对照现在只有一张表，在第 \ref{sec:data} 节。}

\section{这个立场究竟预言什么}\label{sec:cosmos}

纲领现在的立场是：标签是截面局部的，内容是标度轴上的一个切面。
这是一个强主张，而它欠账的地方正是开发中最单薄之处——
暗能量、开端、暗物质，以及最尖锐的挑战：\textbf{共振峰}，
那是有位置、有宽度、有形状的实测对象。本节偿付能付的，并标明付不了的。

\subsection{耗散速率是常数，故 $w$ 不演化}

第 \ref{sec:de} 节\emph{假设}了恒定分数速率 $\lambda$ 并由之导出精确 de Sitter。
\textbf{假设正是弱点}：没有任何东西挑出 $\lambda$ 为常数，
故 $w=-1$ 是一种可能而非预言。

\textbf{A5 提供了它。}设速率依赖于当前标度，$\lambda=\lambda(\sigma)$。
A5 说全局基准平移 $\sigma\mapsto\sigma+c$ 不可观测，而速率\emph{是}可观测的——
它就是 Hubble 率。故 $\lambda(\sigma+c)=\lambda(\sigma)$ 对一切 $c$ 成立，
而在一切平移下不变的函数是常数。

\begin{mdframed}[linecolor=axcol,linewidth=1.2pt,backgroundcolor=axcol!6]
\[
\text{A5} \Longrightarrow \lambda\ \text{恒定} \Longrightarrow
\text{精确 de Sitter} \Longrightarrow w=-1\ \text{且不跑动}
\]
\leanline{Cosmos.shift\_invariant\_is\_constant,
rate\_constant\_of\_fiducial\_invariance, w\_does\_not\_evolve}
\end{mdframed}

\textbf{这是本框架目前最尖锐的可检验陈述：暗能量状态方程的任何演化都证伪 A5。}
而 A5 是全局承重的——它正是使几何无量纲的那一条。\textbf{这个预言不可调。}

\subsection{没有初值，而这是定理}

若速率恒定，膨胀的初条件由什么定？框架的回答是：\textbf{这个问题预设了 A5 所禁止的东西。}

\begin{theorem}[只有比值，没有起点]\label{thm:noorigin}
标度因子只以\emph{比值}进入观测，而比值只依赖于流逝的\emph{差}；
标度是单位故永不消失，历史可延伸到任何给定时刻之下而标度仍为正——
\textbf{没有第一时刻可供赋值}；
然而任何有限的流逝差给出有限比值，故\textbf{观测到的历史是有限且可测的，而起点不存在}。
\leanline{Cosmos.scale\_ratio\_depends\_only\_on\_difference, no\_first\_moment,
finite\_lookback\_no\_origin}
\end{theorem}

\begin{mdframed}[linecolor=axcol,linewidth=1pt,backgroundcolor=axcol!5]
\textbf{宇宙没有年龄，它只有流逝的比值。}
\end{mdframed}

这不是回避：第 \ref{sec:sing} 节早已证明标度不可退化，
而第 \ref{sec:expressive} 节证明绝对标度不可表达。\textbf{第一时刻就是一个绝对标度。}

\subsection{暗物质就是各向同性轨迹}

第 \ref{sec:dm} 节此前只有一条自洽陈述。第 \ref{sec:locus} 节现在提供了解释。

\begin{theorem}[各向同性但量值变化：只引力，无其他]\label{thm:dm2}
$v(x)=f(x)\cdot u$ 的花样楔积恒为零，故不携带转动标签，
从而（定理 \ref{thm:iso}）无荷、无多重态；
但 $f$ 仍在变，故共形因子非常数，\textbf{几何仍然弯曲}。
\leanline{Cosmos.isotropic\_gravitates\_without\_labels}
\end{theorem}

\begin{center}
\textbf{无标签、无荷、无多重态——而仍然引力。}
\end{center}

这就是暗物质的定义性质，\textbf{是导出的而非拟合的}。两条更尖锐的推论：

\begin{itemize}[leftmargin=1.6em]
\item \textbf{暗物质没有物种，也没有质量谱。}标签是区分内容的东西，
      而各向同性区域没有标签，故不携带阈值结构。
      \textbf{在某个确定质量上寻找峰的直接探测计划，
      找的是框架说不存在的东西。}
      \leanline{Cosmos.dark\_matter\_has\_no\_species}
\item 同一论证禁止电荷型自相互作用——这正是"无碰撞"的含义。
\end{itemize}

\begin{mdframed}[linecolor=red!60!black,linewidth=1pt,backgroundcolor=red!3]
\textbf{但这不是唯一的路径，我不该把它当作\emph{那个}答案。}
第 \ref{sec:dm} 节早已有另一条：能标远低于探针的内容——
\textbf{相互作用发生，探测响应消失}。两者是不同机制，
而且做出\textbf{相反}的预言，这比单一故事更好：

\begin{itemize}[leftmargin=1.6em]
\item \textbf{低标度路径}：耦合小但\textbf{非零}，故足够匹配的探针终能看见；
      且这些对象\emph{有}阈值，故\textbf{有质量谱}。
\item \textbf{各向同性路径}：在任何标度上都\textbf{没有可耦合的标签}，
      也没有阈值，故\textbf{没有谱}。
\end{itemize}

\textbf{哪一条是世界所用，是实验问题。}
灵敏度提升下的持续零结果支持各向同性路径；任何谱特征支持低标度路径。
\textbf{框架容纳两者，它不做选择。}
\leanline{Cosmos.two\_dark\_routes\_differ}
\end{mdframed}

\subsection{共振峰}

最难的问题：若标签不实在，共振峰是什么？

\textbf{共振就是阈值。}第 \ref{sec:emergence} 节已把质量与标度轴上的位置等同，
故峰的\emph{位置}就是阈值。新内容是\emph{宽度}。

第 \ref{sec:rg} 节使描述在阈值\textbf{之间}精确；在阈值处它不精确，
而那个过渡区的尺寸就是宽度。定它的是框架自己的不确定性定理：
第 \ref{sec:exact} 小节使标度平移精确非对易，
而第 \ref{sec:nc} 节把对易子变成不确定度乘积。
故对数标度宽度与驻留时间共轭——\textbf{$\Gamma\tau\gtrsim 1$，宽度--寿命关系，
出自框架自己的代数而非引进}。

因 $\mu=\ln m$，对数标度上的宽度就是\textbf{相对}宽度 $\Gamma/m$。
而阈值可被单独分辨，仅当它不与邻居重叠——实测阈值密度把这变成一条界：

\begin{theorem}[可分辨性界]\label{thm:reso}
两个阈值可被单独分辨，仅当其间距超过半宽之和；
而平均间隔在对数标度上是 $1/\rho \approx 1.16$，故
\textbf{$\Gamma/m \gtrsim 1.16$ 的结构根本不是单独的内容}，它们是连续谱。
\leanline{Cosmos.resolvable\_needs\_width\_below\_gap,
broad\_structures\_unresolvable, resolvability\_boundary}
\end{theorem}

\textbf{这是一条真预言，而且在别扭的方向上可证伪}：
一个窄的、无歧义的、良好分离的共振若其 $\Gamma/m$ 超过该界，就与之矛盾。
它同时预言最宽的那些已知结构恰是其"是否为独立态"存在争议的那些——事实如此。

\begin{remark}[这不给什么]
不给耦合常数、不给分宽度、不给产生截面、
不给"阈值有宽度"之外的任何线型。
\textbf{框架说共振\emph{是}什么，并界定它何时算一个共振；它不计算一个共振。}
\end{remark}

\begin{remark}[仍然单薄之处]
微观扇区的\emph{数值}。$\hbar$ 作为标度步长、曲率作为对易子仍是结构而非数值。
但"定量引力链条跑在交换扇区上"这一条\textbf{已由第 \ref{sec:ncconf} 节关闭}。
\end{remark}

\section{框架自己的问题}\label{sec:native}

此前的一切都在回答别处提出的问题——它给不给 Schwarzschild、给不给引力子、给不给面积律。
那是正确的\emph{第一个}检验，而它已经跑过了。但那不是正确的\emph{最后一个}检验：
\textbf{建立在不同原语上的框架，应当能\emph{问出}别人问不出的东西}，
而那些问题正是它要么挣得存在理由、要么被揭穿为一次改写的地方。

原语只有几个：\textbf{作为单位的标度、作为单个矢量的花样、单一漂移、楔积、
阈值密度、各向同／异性二分、一个对易子}。以下逐个追问它们要求知道什么。

\subsection{分离不是原语，非平行才是}

这里没有空间，只有标度。故"两个地方"不能被假设，\textbf{它必须被构造}，
而唯一可用的构造是楔积。

\begin{theorem}[平行的花样不是两个地方]\label{thm:sep}
若两个标度花样平行，则\textbf{没有任何标签区分它们}——
无转动，故无荷，故什么也没有。
而实际存在的分离是\textbf{有向的}：交换二者即反号。
\leanline{Native.not\_two\_places\_of\_parallel, separation\_is\_oriented}
\end{theorem}

\begin{mdframed}[linecolor=axcol,linewidth=1pt,backgroundcolor=axcol!5]
\textbf{框架自己的图像：}广延不是一个供区域安放其中的容器。
两个区域相互区分的程度，恰是其标度花样偏离平行的程度；
而\textbf{完全均匀的区域没有任何内部结构，连广延都没有}。
距离是"非平行度"的一个导出量度。
\end{mdframed}

\emph{这引出的问题：}如果不是距离，那么度量分离程度的观测量是什么？
候选是楔积，而它\textbf{反称}——故框架自己的分离是一个\textbf{有向面积，不是长度}。
\textbf{标准词汇里没有这样的东西。}

\subsection{只有一个开关，不是每种相互作用一个}

定理 \ref{thm:iso} 说各向同性轨迹上无转动标签、无荷、无多重态、无曲率。
\textbf{它们不是四件恰好一起消失的事——它们同源，故只有一个开关。}

\begin{mdframed}[linecolor=axcol,linewidth=1pt,backgroundcolor=axcol!5]
\textbf{第一次各向异性之下什么都没有标签；之上一切标签同时出现。}
其后的结构是简并花样的\emph{细化}（定理 \ref{thm:multthresh}），
\textbf{不是第二次开启}。
\leanline{Native.single\_anisotropy\_boundary, label\_forces\_anisotropy}
\end{mdframed}

\emph{框架自己的问题，而且很尖锐：}标准图景中，不同相互作用在不同标度上变得可分。
此处只有一个边界，其后皆为细化。于是：
\textbf{是否存在某个标签，其出现\emph{不是}对既有各向异性花样的细化？}
一次真正独立的开启就会证伪它。
\textbf{这个问题没人问，因为没有别人有一个"单一开关"需要辩护。}

\subsection{内容属于区间，从不属于标度}

\begin{theorem}[零宽窗口什么也分辨不出]\label{thm:interval}
且任何内容都要求一个阈值严格落在正宽度窗口之内。
\leanline{Native.zero\_window\_resolves\_nothing, content\_is\_interval\_valued}
\end{theorem}

\textbf{故"在能量 $E$ 上存在什么"在这里是不良定的问句。}
存在的东西是一个标度\emph{区间}的性质，\textbf{标度轴上的一个点不携带任何内容}。

\begin{mdframed}[linecolor=axcol,linewidth=1pt,backgroundcolor=axcol!5]
\textbf{框架自己的互补性：标度的锐利与内容的丰富互相对立。}
完全锐利的标度分辨不出任何东西；分辨 $N$ 个结构至少要花 $N/\rho$ 的窗口。
\textbf{这不是海森堡的移植}——此处的共轭对是\textbf{标度锐利度对内容计数}，
而定理 \ref{thm:reso} 的共振宽度界是同一陈述的另一面。
\end{mdframed}

\emph{问题：}使\emph{单个}结构无歧义的最小窗口是多少？$1/\rho$——
这正是可分辨性边界所在之处的原因。

\subsection{没有全局场方程，而 A6 是原因}

定理 \ref{thm:conserv} 证明：场方程可被写下，仅当其源满足循环守恒律。
而 A6 说宇宙耗散——\textbf{全局源不守恒}。

\begin{theorem}[没有全局场方程]\label{thm:noglobal}
不守恒的源不容许任何场方程。故\textbf{整个宇宙没有场方程}。
局域的仍然存在（第 \ref{sec:cov} 节证明局域协变性挺过全局耗散）。
\textbf{宇宙学因此不是"带边界条件全局求解场方程"——没有全局的东西可解。}
\leanline{Native.no\_global\_field\_equation}
\end{theorem}

\emph{框架自己的问题：}若无全局方程，是什么定住了大尺度历史？
可用的答案是 A5（定理 \ref{thm:a7} 所在小节的常速率结果）——
\textbf{那是一个对称性，不是一条动力学定律}。

\begin{center}
\textbf{大尺度历史由"什么不可观测"定住，而不是由"什么在演化"定住。}
\end{center}

这够不够，是开放问题，\textbf{而它不是广义相对论提得出的问题}。

\subsection{量子引力实验是一次"次序检验"}

定理 \ref{thm:ncdefm} 使全部量子修正就是标度梯度的对易子。
\textbf{而对易子是一个\emph{次序}差}，故原生实验不是"探测引力子"，而是：

\begin{mdframed}[linecolor=axcol,linewidth=1.2pt,backgroundcolor=axcol!6]
\textbf{沿两个方向比较标度，以两种次序各做一次，看是否有差异。}

两次序之差恰为 $-2[\sigma_i,\sigma_j]$，\textbf{那就是几何的全部量子内容}。
无需测别的，\textbf{也无需任何粒子存在}。
\leanline{Native.ordering\_discrepancy\_is\_the\_observable}
\end{mdframed}

\emph{框架自己的问题：}"沿一个方向比较标度"由什么物理程序实现？
标度比较是局域单位之比，故程序是一次双臂比较——
而框架说\textbf{两臂以相反次序遍历会相差一个对易子}。
量级由标度步长定，远低于现有一切，\textbf{这一点被明说而非隐藏}；
原生的是\textbf{观测量是一个次序，而不是一个振幅、也不是一个量子}。

\subsection{缺陷成对出现，且宇称相同}

\begin{theorem}[宇称守恒，作为算术]\label{thm:parity}
$\mathbb Z/2$ 标签相乘，且每个元素是自身的逆，
故绑定两个同宇称缺陷得到偶。整数荷同样恒等守恒。
\leanline{Native.parity\_conserved\_under\_binding, both\_labels\_conserved}
\end{theorem}

\textbf{框架自己的选择定则：}缺陷成对产生与湮灭，且宇称相同，
\textbf{没有自由参数，没有耦合常数}。
它不是加在拉氏量上的对称性——\textbf{这里没有拉氏量}——它是 $\mathbb Z/2$ 中的算术。

\emph{问题：}有没有任何过程改变总宇称？框架说\textbf{绝对没有}，
而这比"由假设的对称性导出的选择定则"更干净，\textbf{因为没有东西可供破缺}。

\subsection{没有最大密度}

\begin{theorem}[标度之比是单位：永不为零、永不为无穷、但量值无界]\label{thm:nomax}
物质在此就是标度的各向异性，故\textbf{没有最大物质密度、没有 Planck 密度}，
也没有任何标度使描述必须被另一类东西取代。
这是定理 \ref{thm:nomin} 从物质一侧看到的同一件事。
\leanline{Native.no\_maximum\_anisotropy, no\_cutoff\_either\_end}
\end{theorem}

\subsection{这七条的共同点}

\textbf{没有一条是翻译。}每一条都来自标准词汇所没有的一个原语：
作为单位的标度（VII）、作为单个矢量的花样（I）、阈值密度（III）、
单一二分（II）、不守恒的全局源（IV）、
对易子而非场（V）、以及\textbf{作为算术而非对称性}的 $\mathbb Z/2$（VI）。

\begin{remark}[诚实的状态]
七条之中，II 与 VI 原则上现在可检验；III 已通过共振宽度界被检验；
VII 只能由其逆否命题检验；而 I、IV、V 是概念性主张，\textbf{目前没有任何测量}。
\textbf{没有一条是对某个数值的预言。}
\end{remark}

\section{与数据对质}\label{sec:data}

预言已经列出；本节把它们带到测量面前。
\textbf{其中一条有麻烦，而且是最承重的那一条，故先说它。}

\subsection{$w$ 不演化——这与 DESI 有张力}

定理 \ref{thm:a7} 所在小节由 A5 导出恒定耗散速率，故 $w=-1$ 精确且每个时期皆然。
\textbf{A5 不是框架中可调的部分}：它正是使几何无量纲的那一条，故此预言无法软化。

DESI DR2（2025）把 BAO 与 CMB 及超新星联合，对演化状态方程的偏好达到
\[
3.1\sigma\ /\ 2.8\sigma\ /\ 3.8\sigma\ /\ 4.2\sigma
\]
（分别对应 CMB、Pantheon$^+$、Union3、DESY5），
落在 $w_0>-1$、$w_a<0$ 的象限，$w$ 在 $z\approx 0.5$ 附近穿越 $-1$。

\begin{mdframed}[linecolor=red!60!black,linewidth=1.2pt,backgroundcolor=red!4]
\textbf{这是本框架最严重的经验问题，我不打算把它解释掉。}
三件事同时为真，三件都该说：

\begin{itemize}[leftmargin=1.6em]
\item \textbf{DESI BAO 单独与 $\Lambda$CDM 完全自洽。}张力只存在于联合中，
      而\textbf{各联合彼此不一致}：BAO$+$CMB 给出 $w_0\approx-0.42\pm0.21$，
      BAO$+$CMB$+$SN 给出 $w_0\approx-0.838$——
      \textbf{同一参数上相差两个标准差}；
\item 偏好在最弱与最强超新星编目之间\textbf{恰好差 1.5 倍}；
\item 独立分析指出这些数据集\textbf{相互不自洽}，
      而产生信号的正是那一步联合；
\item \textbf{若该偏好在数据集自洽的前提下达到 $5\sigma$，A5 被证伪}，
      连同整个开发所依赖的无量纲性一起。
      \textbf{不存在带演化 $w$ 的本框架版本。}
\end{itemize}
\leanline{Data.desi\_tension\_exceeds\_three\_sigma,
desi\_significance\_dataset\_dependent, falsification\_threshold}
\end{mdframed}

\subsection{共振宽度界——两个旗舰确定值分居两侧}

\begin{mdframed}[linecolor=red!60!black,linewidth=0.9pt,backgroundcolor=red!3]
\textbf{对本节前一版本的更正。}我曾把 $f_0(500)$ 极点引作 $457-i279$，
得 $\Gamma/m=1.22$，并报告该界被违反。
\textbf{$279$ 不是两个标准确定值中的任何一个。}正确的两个是：
\end{mdframed}

\begin{center}
\begin{tabular}{lll}
\toprule
方法 & 极点 & $\Gamma/m$ \\
\midrule
CCL（Bern，Roy 方程 $+$ ChPT 输入） & $441-i272$ & $\mathbf{1.234}$ \ 界外 \\
GKPY（Madrid，一次减除，纯 $\pi\pi$ 数据） & $457-i249$ & $\mathbf{1.090}$ \ 界内 \\
GKPY 更新收敛 & $(445\text{--}450)-i(220\text{--}245)$ & $0.978$--$1.101$ \ 全界内 \\
\bottomrule
\end{tabular}
\end{center}

\begin{theorem}[界恰落在两个旗舰确定值之间]\label{thm:bound}
$1.090 < 1.157 < 1.234$，而更新后的 Madrid 值\textbf{全部在界内}。
\leanline{Data.ccl\_above\_bound, gkpy\_below\_bound,
determinations\_straddle\_bound, gkpy\_updated\_inside}
\end{theorem}

\begin{mdframed}[linecolor=axcol,linewidth=1pt,backgroundcolor=axcol!5]
\textbf{故该界目前\emph{并未}被违反。}
它坐落在介子表中最有争议的共振的最佳确定值的分辨极限上，
而两种方法对真极点在界的哪一侧\textbf{意见相左}。
框架因此对一个具体的实测数值作出一条\textbf{活的、可证伪的、当前开放的}主张：

\begin{center}
$f_0(500)$ 极点满足 $\Gamma/m < 1.157$——
\textbf{框架站在一次减除的确定值一边，反对 Roy 方程加 ChPT 的那个。}
\end{center}
\end{mdframed}

谱的其余部分远在界内：$W,Z,t,H$ 为 $10^{-5}$--$0.03$；
$\Delta(1232)$、$N(1440)$ 为 $0.10$、$0.15$；$\rho(770)$ 为 $0.19$；
$f_0(1370)$ 为 $0.26$；$K_0^*(700)$ 极点为 $0.86$。
\textbf{界仍然落在两个最宽态之间}，这正是它是检验而非描述的原因。
\leanline{Data.bound\_separates\_kappa\_from\_sigma}

而 Breit--Wigner 范围对检验\textbf{依然无用}：PDG 引用 $m\in[400,800]$、
$\Gamma\in[100,800]$，使 $\Gamma/m$ 从 $0.125$ 跨到 $2$。
\textbf{故陈述必须读在极点上，现在也确实如此。}
\leanline{Data.breit\_wigner\_spans\_the\_bound}

\begin{remark}[一条观察，作为观察而非推导给出]
文献中有一支把 $f_0(500)$ 读作\textbf{标度不变性破缺的赝 Goldstone 玻色子}——
一个膨胀子。本框架没有任何"其激发是粒子"的场，而标度不是这样的场，
故它\textbf{不预言存在膨胀子\emph{态}}；
它预言的是：任何扮演该角色的东西恰好坐在\textbf{可分辨性边界}上，
既不明确是态、也不明确是连续谱。\textbf{$f_0(500)$ 正在那里。}
有启发性；此处\textbf{不推导任何东西}。
\end{remark}

\subsection{光子无色散——通过，而且在分道}

定理 \ref{thm:nodisp} 预言\textbf{恰好为零}的色散，即 $E_{\rm QG}=\infty$。
Fermi-LAT 对 GRB 090510 的飞行时间给出
$E_{{\rm QG},1}\gtrsim 7.6\,E_{\rm Pl}$（总色散）、$\gtrsim 2\,E_{\rm Pl}$（LIV 诱导），
近期分析要求 $E_{{\rm QG},1}\gtrsim 1.2\times10^{19}$ GeV $\approx E_{\rm Pl}$。
2024 年一项谱延迟研究给出后验最大值 $E_{\rm QG}\approx 9\times10^{14}$ GeV，
但报告数据与精确 Lorentz 不变性相容。

\begin{mdframed}[linecolor=axcol,linewidth=1pt,backgroundcolor=axcol!5]
\textbf{限已经在 Planck 能量之上}，而那正是最小长度纲领一般安置该效应之处。
\textbf{每一次收紧都朝本框架靠、朝对手离。}
这是框架目前最干净的活判别式，\textbf{而它正在赢}。
\leanline{Data.dispersion\_limit\_above\_planck, framework\_predicts\_no\_finite\_scale}
\end{mdframed}

\subsection{恰好一个非禁闭精确守恒荷——而无中微子双贝塔衰变裁决它}

第 \ref{sec:particle} 节给缺陷恰好一个整数荷，非禁闭且精确守恒。
自然界的候选是电荷与 $B-L$。
\textbf{而 $B-L$ 在标准模型中精确守恒（含非微扰），故表面上有\emph{两个}}，
这是框架不允许的。

这不是僵局：\textbf{无中微子双贝塔衰变把 $B-L$ 破坏两个单位}。于是

\begin{itemize}[leftmargin=1.6em]
\item \textbf{观测到 $0\nu\beta\beta$} $\Rightarrow$ $B-L$ 不守恒 $\Rightarrow$
      电荷是唯一的非禁闭精确守恒整数荷 $\Rightarrow$ \textbf{框架的计数正确}；
\item \textbf{排除到 $B-L$ 看起来精确且被规范化} $\Rightarrow$ 有两个 $\Rightarrow$
      \textbf{框架的计数错误}。
\end{itemize}

\leanline{Data.one\_charge\_prediction, zero\_nu\_beta\_beta\_decides}

这是一个\textbf{正在进行的实验}（KamLAND-Zen、LEGEND、nEXO），
而框架关心它的结果，出于一个别人没有的理由。

\subsection{总表：全部数值与全部预言}\label{ss:scoreboard}

\textbf{本报告只有这一张数值表。}第 \ref{sec:pred} 与 \ref{sec:spectrum} 节
原各有一张部分重复的表，已并入此处。

\begin{center}
\small
\begin{tabular}{llll}
\toprule
量 & 本框架 & 实测 & 状态 \\
\midrule
\multicolumn{4}{l}{\emph{从公理链条导出（外加 A6$'$）；数值证到所引精度}}\\
掠日光线偏折 & $1.7515''$ & $1.7509\pm0.0002''$ & 符合 \\
水星近日点进动 & $42.989''$ & $42.98\pm0.04''$ & 符合 \\
PPN $\gamma$ & $1$（精确） & $|\gamma-1|\lesssim2\times10^{-5}$ & 符合 \\
\midrule
\multicolumn{4}{l}{\emph{A3 读法的算术核对，非独立预言}}\\
$\mu$ 子飞行寿命 & $64.43\,\mu$s & $64.378\pm0.026$ & 符合 $10^{-3}$ \\
Pound--Rebka & $2.455\times10^{-15}$ & $(2.46\pm0.03)\times10^{-15}$ & 符合 \\
\midrule
\multicolumn{4}{l}{\emph{换算而来；$\lambda$ 本身是测量值，即那个唯一的自由数}}\\
暗能量密度 & $5.845\times10^{-27}$ & $5.86\times10^{-27}$ kg/m$^3$ & 符合 \\
\midrule
\multicolumn{4}{l}{\emph{当前活的检验}}\\
$w$ 不演化（A5 强制） & $w=-1$，不跑动 & DESI: $w_0>-1,w_a<0$ & \textbf{张力 $2.8$--$4.2\sigma$} \\
$f_0(500)$ 极点 $\Gamma/m$ & $<1.157$ & CCL $1.234$ / GKPY $1.090$ & \textbf{开放} \\
光子能量依赖时延 & 恰为零 & $E_{\rm QG}\gtrsim7.6E_{\rm Pl}$ & 通过，且与对手分道 \\
引力波速 & 精确 $c$，无色散 & $|\Delta v/v|\lesssim10^{-15}$ & 通过 \\
引力波偏振数 & 恰为 $2$，无呼吸模 & 网络偏振检验 & 通过（与标量--张量分道）\\
偶极辐射 & 恒等于 $0$ & 双脉冲星 $<1.3{\times}10^{-4}$ & 通过 \\
铃宕回波 & 反射率 $\neq0$ & GW250114 残差 $=$ 噪声 & \textbf{开放，受压} \\
能层不稳定性 & 反射率上界未定 & GW241011 $\chi_1=0.64$ 长寿 & \textbf{开放，压力更大} \\
超轻标量 & \textbf{无此槽位} & $10^{-13}$--$3{\times}10^{-12}$ eV 被排除 & 相容 \\
面积定理类比 & 标度比值不减 & 与观测一致 & 通过（但不具分辨力）\\
质量谱 $\mathrm{CV}^2$ & $<1$（秩一） & $0.875$ & 通过 \\
非禁闭精确守恒荷 & 恰一个 & 电荷、$B{-}L$（两个？） & 未决，$0\nu\beta\beta$ 裁决 \\
黑洞熵 & $\propto\ln R$ & —— & 未测（面积律亦未测）\\
\midrule
\multicolumn{4}{l}{\emph{框架内部被排除的选项——最有信息量的部分}}\\
等比质量塔 & $\mathrm{CV}=0$ & $0.935$ & \textbf{排除} \\
各向同性扇区偏折 & $0.8758''$ & $1.7509''$ & \textbf{排除（2 倍）} \\
各向同性扇区进动 & $14.33''$ & $42.98''$ & \textbf{排除（3 倍）} \\
色 $=$ 块荷 & $\mathbb Z/2$（两值） & 需要三值 & \textbf{排除（结构不可能）} \\
\bottomrule
\end{tabular}
\end{center}

\begin{mdframed}[linecolor=axcol,linewidth=1pt,backgroundcolor=axcol!5]
\textbf{一条活张力、一个活的开放问题，都不致命}——而这两个扇区正在被测量。
最后四行是本体系\emph{拒绝}掉的东西，
\textbf{它们比符合的那些更有信息量，因为它们说明这套体系确实会拒绝}。

\textbf{关于电荷检验的时间尺度：}即便建立在跷跷板机制上的轻子数破坏模型，
所预言的 $0\nu\beta\beta$ 半衰期也远高于现有实验的排除范围。
\textbf{这是一个真实但引信很长的检验}，“正在进行”不应被读作“即将报告”。
\leanline{Data.scorecard}
\end{mdframed}


\section{如何证伪本框架}\label{sec:falsify}

一份内部报告最该有的一节。以下每一行给出\textbf{一个观测结果，它若成立，框架的哪一部分就死}。
按"死得多彻底"排序，不按可能性排序。

\subsection{会杀死整个框架的}

\begin{center}
\begin{tabular}{p{0.30\textwidth}p{0.28\textwidth}p{0.30\textwidth}}
\toprule
若观测到 & 则死的是 & 现状 \\
\midrule
\textbf{$w$ 随红移演化}（数据集自洽前提下 $5\sigma$） & \textbf{A5}，连同无量纲性，连同全部几何 & DESI $2.8$--$4.2\sigma$，但 BAO 单独自洽、各联合互相矛盾 \\
\textbf{光子能量依赖时延}（任何阶） & A3 加类光锥标度盲，即第 \ref{sec:light} 节的地基 & $E_{\rm QG}\gtrsim7.6E_{\rm Pl}$，正在赢 \\
\textbf{引力波偶极辐射} & 守恒定理，即 Bianchi 出发的整条链 & 双脉冲星 $<1.3\times10^{-4}$ \\
\textbf{引力波呼吸模} & A5（呼吸模就是基准平移） & 张量优于标量／矢量／混合 \\
\bottomrule
\end{tabular}
\end{center}

\textbf{这四条没有可调参数。}标量--张量理论有耦合可以调小；这里没有——
\textbf{测到就是死，不是"需要更小的耦合"。}

\subsection{会杀死一个扇区的}

\begin{center}
\begin{tabular}{p{0.30\textwidth}p{0.28\textwidth}p{0.30\textwidth}}
\toprule
若观测到 & 则死的是 & 现状 \\
\midrule
\textbf{存在最小长度} & 第 \ref{sec:horizon} 节全部（无奇点／无视界／无最大密度同源） & 与光子色散同一条界 \\
\textbf{$f_0(500)$ 极点 $\Gamma/m>1.157$ 被确定} & 第 \ref{sec:spectrum} 节的可分辨性界 & CCL $1.234$／GKPY $1.090$，框架押 GKPY \\
\textbf{$B-L$ 被规范化且精确守恒} & 第 \ref{sec:particle} 节"恰一个整数荷" & $0\nu\beta\beta$ 未见，引信很长 \\
\textbf{某个标签的出现不是既有各向异性的细化} & 第 \ref{sec:locus} 节的单一开关 & 无已知反例 \\
\textbf{总宇称改变的过程} & 第 \ref{sec:native} 节的 $\mathbb Z/2$ 算术 & 无已知过程 \\
\bottomrule
\end{tabular}
\end{center}

\subsection{会迫使框架承认它算不了的}

\begin{center}
\begin{tabular}{p{0.30\textwidth}p{0.28\textwidth}p{0.30\textwidth}}
\toprule
若观测到 & 则暴露的是 & 现状 \\
\midrule
\textbf{高自旋致密天体因能层不稳定而衰减} & 无视界结论\textbf{要求}反射率极小，而框架\emph{算不出}它 & GW241011 $\chi_1=0.64$ 长寿，\textbf{压力最大的一条} \\
\textbf{铃宕回波} & 同上，反向：框架预言其存在却给不出振幅 & 搜索为零结果 \\
\textbf{黑洞熵 $\propto$ 面积而非 $\ln R$} & 第 \ref{sec:horizon} 节的对数熵 & 熵本身未被测量 \\
\bottomrule
\end{tabular}
\end{center}

\begin{mdframed}[linecolor=red!60!black,linewidth=1.2pt,backgroundcolor=red!4]
\textbf{无法被这份表证伪的，是 A7。}
它是关于\emph{什么算一次观测}的承诺，任何人可以拒绝它并得到别的 $k$——
\textbf{那不是一次实验能裁决的事}，这也正是它被扶正为公理而非留作判断的理由。

同样无法被证伪的是那\textbf{一个自由单位}：它是预言被表达于其中的尺度，
\textbf{不是一条可错的预言}。
\end{mdframed}

\begin{remark}[什么\emph{不}构成证伪]
"框架没有给出标准模型的粒子表"不是证伪——第 \ref{sec:emergence} 节论证那不是目标。
"它没算出强子质量"不是证伪——那与那一个自由单位是同一件事。
\textbf{"它与广义相对论在引力波扇区一致"更不是证伪}——
第 \ref{sec:waves} 节证明对称扇区在\emph{每一阶}上都与 GR 相同，
故那个扇区\textbf{原则上}就不该指望有分歧。
\end{remark}

\clearpage
\part{第九部分\quad 登记册：撤回、更正、与边界}

\noindent
\textbf{以下各节不是结果，是记录。}它们此前被放在结果部分，那是错的呈现——
色荷的识别已被撤回、余维数分析描述的不是本框架、跨扇区检验并不存在、
表达力那节的结论被耦合推翻。\textbf{它们的定理仍然成立且仍被审计}，
但它们所支持的\emph{解读}不再成立，故迁到此处。

\textbf{一个没有撤回记录的开发，是一个没有被审计的开发。}

\section{自我审查：已发现的思维污染}\label{sec:audit}

本节是对整个体系的一次不设辩护的复查，按"基础建构"与"目标问题"两条线。
发现的问题分三类：\textbf{已修复}、\textbf{须撤回}、\textbf{仍存在}。
第三类未被修复，只被登记。

\subsection{须撤回：质量塔曾被当作结论，其实是假设}

第 \ref{sec:defect} 节写下 $m_k=m_0e^{|k|\Delta}$，并说它由 A3
（质量即逆标度）加"每单位绕数跨一个周期"得出。第 \ref{sec:qcd} 节据此检验格点胶球比值，
报告\textbf{被证伪}。

复查发现\emph{证伪打错了靶}。框架真正确立的关于缺陷的性质是\textbf{拓扑}的：
绕数是整数、可加、不能连续改变、$|k|$ 在取向反转下不变。
而质量律额外需要的是\textbf{动力学}的：绕数为 $k$ 的缺陷使标度场应变多少——
本报告\emph{没有任何定理}涉及这件事。

\begin{theorem}[拓扑不决定质量律]\label{thm:massgap}
存在两个真正不同的质量律——指数塔与二次律——
\textbf{都}满足关于绕数已证明的全部性质（正性、$k=0$ 处归一、取向反转不变）。
\leanline{MassAudit.mass\_law\_underdetermined,
MassAudit.both\_laws\_antiparticle\_degenerate}
\end{theorem}

\begin{mdframed}[linecolor=red!60!black,linewidth=0.9pt,backgroundcolor=red!3]
\textbf{撤回.} 几何质量塔从\emph{预言}降格为\emph{假设}。
定理 \ref{thm:towerfail} 的胶球检验因此否定的是\textbf{那个假设}，
不是缺陷图景、也不是框架——这比上一版报告的结论\emph{弱}得多，
本节存在的目的就是让较弱的那个成为记录在案的版本。

\emph{幸存者}：粒子—反粒子质量简并（只依赖 $|k|$，定理 \ref{thm:anti}）；
无参数质量比（不用任何质量律，定理 \ref{thm:massratio}）。
\end{mdframed}

\subsection{已修复：光是标度盲的，红移却要靠光}

第 \ref{sec:light} 节证明类光锥\textbf{标度盲}；
第 \ref{sec:cov} 节证明宇宙漂移\textbf{局域完全不可观测}。
那么膨胀究竟怎么被看见？答案本应是红移——由光携带。
但光若对标度是盲的，它就无法携带标度的记录。

两条都已被证明，故都为真，"矛盾"意味着有东西被混同了。确实有。

\begin{mdframed}[linecolor=axcol,linewidth=0.9pt,backgroundcolor=axcol!4]
\textbf{类光矢量没有长度，但有能量——而能量不是矢量的属性。}
它是矢量与\emph{观察者}的配对，而观察者由物质构成，
物质按定理 \ref{thm:massive} 确实携带标度。
\end{mdframed}

\begin{theorem}[红移属于观察者，不属于光]\label{thm:redshiftobs}
类光方向可以携带非零能量；且两个局域单位为 $u_1,u_2$ 的观察者
对\emph{同一}类光方向给出的能量之比是 $(u_1/u_2)^2$——
光对这个比值没有贡献，全部来自两个局域单位之差。
\leanline{Observation.null\_has\_energy, Observation.redshift\_is\_scale\_ratio}
\end{theorem}

这同时化解了第二处张力：A5 说全局基准平移不可观测，A6 说宇宙耗散。
红移测量并不是对漂移的局域测量，而是比较\emph{发射时}与\emph{吸收时}的标度——
那是一个标度\emph{差}，正是 A5 判为物理的量。
\textbf{宇宙学观测到的恰好是公理允许的那个量，不多不少。}
\leanline{Observation.redshift\_fiducial\_invariant}

\subsection{登记册（一）：已撤回的主张}\label{ss:retracted}

\textbf{以下四条曾被作为结论提出，现已收回。} 其定理仍成立且仍被审计，
被收回的是\emph{解读}。

\begin{enumerate}[leftmargin=1.8em,itemsep=3pt]
\item \textbf{几何质量塔}。$m_k=m_0e^{|k|\Delta}$ 是\emph{假设}而非推论
      （本节上文）。定理 \ref{thm:towerfail} 的胶球检验否定的是那个假设。
      第 \ref{sec:emergence} 节此后说明了为什么不存在这样一条律：
      质量不是缺陷\emph{带着}的属性，而是它在标度轴上\emph{出现的位置}。
\item \textbf{色的识别}。第 \ref{sec:color} 节导出的有限禁闭荷不是色。
      第 \ref{sec:particle} 节表明块群只能是 $\mathbb Z/2$——序参量是射影的——
      故\textbf{三值从来不在供给里}。这不是缺口，是被解释了的结构不可能。
\item \textbf{跨扇区检验}。第 \ref{sec:crosscheck} 节：谱的 $\rho$ 与跑动的 $\rho$
      不是同一个量（一个是原始计数，一个是扇区受限且带符号的）。
\item \textbf{余维数分析}。第 \ref{sec:codim} 节的定理对\emph{圆值}标度成立，
      但那不是本框架：方向性 A4 下序参量是无取向轴，点缺陷回来了（第 \ref{sec:axis} 节）。
\end{enumerate}

\subsection{登记册（二）：已限缩适用域的主张}\label{ss:corrected}

\textbf{定理成立，当初的解读过宽。}

\begin{enumerate}[leftmargin=1.8em,itemsep=3pt]
\item \textbf{“对易导子是隐藏的平直性假设”}——\emph{撤回}。把导子全部杀死，
      曲率原封不动（定理 \ref{thm:noflatcomm}）。一张图就是一张图；
      平直性住在\emph{联络}对不对易里。
\item \textbf{“方向就是代数，$n$ 是代数维数”}——\emph{撤回}。
      $\dim\mathfrak{so}(3,1)=6\neq4$。方向是代数所作用的\textbf{定义表示}
      （定理 \ref{thm:algbig}）。什么也没丢：定名代数仍然定住 $n=k+1$。
\item \textbf{“多重态内部平直”}——解读\emph{撤回}。其前提按方向指派缩放算子，
      而实秩一只提供一个矢量、一个缩放（第 \ref{sec:axes} 节）。
\item \textbf{“单轴不携带转动”}——定理讲的是\emph{均匀}花样。
      由此产生的两难随之消解（第 \ref{sec:axes} 节）。
      \textbf{这一条挺过了一次形式化与一版报告才被抓住。}
\item \textbf{“标度路与转动路互相独立”}——\emph{推翻}。那是对标量扇区的分析；
      方向性标度下转动扇区由标度扇区\emph{生成}（第 \ref{sec:coupling} 节）。
\item \textbf{“$f_0(500)$ 违反宽度界”}——\emph{更正}。此前引用的极点虚部
      不是任何一个标准确定值。以真值计，界落在 CCL 与 GKPY \emph{之间}：
      开放，未被违反（第 \ref{sec:data} 节）。
\end{enumerate}

\begin{mdframed}[linecolor=red!60!black,linewidth=1pt,backgroundcolor=red!3]
\textbf{上列六条中有五条是同一个错误：在结构是方向性的地方用标量推理。}
归属取自各文件\emph{自己}的自述，而非记忆：

\begin{enumerate}[leftmargin=1.8em,itemsep=1pt]
\item \textbf{第 \ref{sec:frame} 节}——几何：标量标度禁止 Schwarzschild；
\item \textbf{第 \ref{sec:codim} 节}——缺陷：圆值（标量）标度预言弦，
      由第 \ref{sec:axis} 节的无向轴更正；
\item \textbf{第 \ref{sec:coupling} 节}——地基：“比值可交换”是标量陈述；
\item \textbf{第 \ref{sec:direction} 节}——对易导子的平直性解读；
\item \textbf{第 \ref{sec:axes} 节}——$\text{wedge}\,v\,v=0$ 的单轴解读，
      而它讲的是\emph{均匀性}。
\end{enumerate}

\textbf{此前本表把第 2、5 条记在第 \ref{sec:axis}、\ref{sec:slice} 节名下。
那是\emph{更正}它们的节，不是\emph{犯下}它们的节。}计数不变，归属已改。

\textbf{抓住它的判据：}凡是提到一个计数——一根轴、一个标度、一个方向、一个值——
就问\emph{是什么的计数}。\textbf{值与方向每一次都会分开。}
\end{mdframed}

\subsection{登记册（三）：已解决的项目}\label{ss:closed}

\textbf{以下曾在本清单上，现已关闭。}列出是为了让读者能核对，而不是为了记功。

\begin{center}
\small
\begin{tabular}{lll}
\toprule
曾登记的问题 & 现状 & 由何关闭 \\
\midrule
$\kron$ 是被公设的平直度规 & 已导出 & $\kappa(x,y)=\tau(xy)$，仅需一个迹（\S\ref{sec:invariant}）\\
号差 $\eta$ 是输入 & 已导出 & 漂移区分出恰好一个方向（\S\ref{sec:signature}）\\
$n$ 是独立输入 & 不再是 & 代数定名后 $n=k+1$（\S\ref{sec:dimension}）\\
$k=3$ 依赖一个判断 & 已扶正为 A7 & 公开的第七条公理（\S\ref{sec:locus}）\\
单轴／双轴两难 & 已消解 & 误读所致；更多轴不可得（\S\ref{sec:axes}）\\
基本激发形态未定 & 已澄清 & 射影序参量给出整数荷点缺陷（\S\ref{sec:axis}）\\
交叉积未完整构造 & 结构部分已解决 & 实秩一使作用群为 $\mathbb R$（\S\ref{sec:algebra}）\\
定量引力跑在交换扇区 & 已关闭 & 非交换曲率（\S\ref{sec:ncconf}）\\
耗散率 $\lambda$ 是否恒定 & 已定 & A5 强制它恒定（\S\ref{sec:cosmos}）\\
两路独立故无跨扇区检验 & 前提被推翻 & 耦合被词义强制（\S\ref{sec:coupling}）\\
\bottomrule
\end{tabular}
\end{center}

\textbf{注意“$k=3$ 已扶正为 A7”不是一次解决，是一次\emph{降级}。}
它从“几乎已导出”变成公开的公理，读者可以拒绝它并得到别的 $k$。

\subsection{登记册（四）：仍然开放}\label{ss:open}

\textbf{以下没有被解决，只被登记。}

\begin{enumerate}[leftmargin=1.8em,itemsep=3pt]
\item \textbf{$a=1/\varepsilon$ 是识别，不是导出。} 把宇宙学标度因子认作逆能标
      是一个物理假设。de Sitter 的数学（定理 \ref{thm:desitter}）无懈可击，
      但它依赖这个识别。
\item \textbf{标度花样是输入。} 哪些方向简并、简并到什么程度，
      框架不预言。第 \ref{sec:gauge} 节的全部规范结论都是\emph{给定花样}的函数。
\item \textbf{阈值密度 $\rho$ 的\emph{数值}是测量的。} 第 \ref{sec:spectrum} 节
      预言的是\emph{离散度}（$\mathrm{CV}$），不是密度；
      第 \ref{sec:dimension} 节的 $\rho\alpha=k-1$ 是换算而非预言。
\item \textbf{暗物质：两条路径，框架不选。} 第 \ref{sec:dm} 节的低标度路径
      证明的是\emph{相容性}而非解释；第 \ref{sec:cosmos} 节的各向同性路径
      是一条导出，但两者预言相反（有谱／无谱），\textbf{哪条为真是实验问题}。
\item \textbf{球对称拟设本身是输入。} 第 \ref{sec:riccidiag} 节\emph{已}对
      各向异性度规 $-A(r)\dd t^2+B(r)\dd r^2+r^2\dd\Omega^2$ 逐分量算出 Ricci
      （$A\neq B$ 正是各向异性），而“真空方程 $\Rightarrow$ 标度和为常数”
      也已形式化（定理 \ref{thm:vacstat}）。
      \textbf{此前本条写作“逐分量计算未做”，那是 \texttt{RicciDiag.lean} 写出之前的旧登记，已更正。}
      仍是输入的是\emph{静态球对称拟设}的选取本身，以及标度响应律 A6$'$。
\item \textbf{铃宕边界条件、回波振幅，以及\emph{能层不稳定性}。}
      这是波动扇区里\emph{唯一}一处本框架与广义相对论可能分开的地方
      （第 \ref{sec:waves} 节）。
      \textbf{新登记（来自 GW241011）}：无视界天体的能层会因超辐射而不稳定，
      而高速自旋且长寿的遗迹\textbf{指数地}压住反射率——
      这比回波搜索的界严厉得多，且框架\textbf{尚不能说它是否幸存}。
      没有视界就没有"在视界上纯内向"这个边界条件，
      而部分反射一般产生晚期回波。
      \textbf{框架说反射率不为零，但不预言振幅}——那由那一个自由单位决定。
      而回波搜索目前是\textbf{零结果}。\textbf{登记为开放且在受压。}

      \emph{不}登记为缺口的：波形、辐射功率、后 Newton 修正、数值相对论。
      在各向同性扇区 SCD 的场方程与 GR 相同（$\gamma=1$，同一个 $\kappa$），
      故把它们算出来\textbf{只会复现 GR，学不到关于 SCD 的任何东西}。
      力气应当花在两者\emph{分开}的地方。
\item \textbf{微观扇区的数值等价于那一个自由单位。} 措辞需要精确：
      框架给出量子修正的\emph{形式}（$[\sigma_i,\sigma_j]$）、\emph{对称性}（纯反对称）
      与\emph{标度}（相对经典项为一个标度步长）。
      \textbf{一旦任何一次测量把那个单位定住，量级就随之定住}——
      所以这不是一处独立的缺口，而是"唯一一个自由数"的又一次现身
      （第 \ref{sec:dimension} 节）。
      \emph{真正}缺的是\textbf{散射振幅}：它需要固定内容的渐近态，
      而第 \ref{sec:emergence} 节说没有哪个标度携带完整清单。
\end{enumerate}

\subsection{登记册（五）：引用而未证，与登记的假设}\label{ss:cited}

\textbf{引用（2 条）。}若引用有误，其上的\emph{约束}仍成立，只是\emph{识别}不成立。
\begin{enumerate}[leftmargin=1.8em,itemsep=2pt]
\item 实秩一实单李代数的分类（$\mathfrak{so}(n,1)$、$\mathfrak{su}(n,1)$、
      $\mathfrak{sp}(n,1)$、$\mathfrak f_{4(-20)}$）。秩 $=1$ 本身\emph{是}被证明的；
      清单的\emph{穷尽性}不是。
\item 射影空间的同伦：$\pi_1(\mathbb{RP}^n)=\mathbb Z/2$、$\pi_2(S^2)=\mathbb Z$、
      以及李群模离散子群的 $\pi_2$ 为零。
\end{enumerate}

\textbf{假设（3 条）。}
\begin{enumerate}[leftmargin=1.8em,itemsep=2pt]
\item \textbf{A6$'$，标度响应律} $e^{2\sigma}R=\kappa\rho$。
      \textbf{它不是七条公理之一}，是一条外部物理输入，
      在第 \ref{sec:newton}、\ref{sec:riccidiag}、\ref{sec:precession}、
      \ref{sec:waves} 节承重。它在每一条用到它的定理里都是显式前提，
      故类型层面没有隐瞒——但它\textbf{此前不在本登记册里}，
      而这里正是读者用来回答"这是否只靠公理"的地方。
      \textbf{经外部审核指出后补入；这个遗漏是本登记册最严重的一处。}
\item 缺陷在对称空间体积中\textbf{均匀分布}——第 \ref{sec:dimension} 节的
      $\rho\alpha=k-1$ 换算需要它，无定理提供。
\item \textbf{A7 本身}——关于什么算一次观测的承诺，已扶正为公理而非留作判断。
\end{enumerate}

\subsection{登记册（六）：前提无见证的结构}\label{ss:unwitnessed}

一个反复出现的形状，最初为 \texttt{Dynamics.Realizes} 登记，
后由外部审核发现范围更广。形如"给定结构 $S$，……"的定理，
其强度取决于框架\emph{能否造出}一个 $S$。造不出时，定理为真但\textbf{无见证}，
其物理解读只是一个条件句。

\begin{mdframed}[linecolor=red!60!black,linewidth=1.2pt,backgroundcolor=red!4]
\textbf{最严重的一处：\texttt{Crossed.ScaleShift}}——
即"$\hbar$ 是精确的标度步长"这一\emph{招牌}结果的载体。
它\textbf{从未由 \texttt{ScaleAlgebra}／\texttt{DiffRing}／\texttt{ScaleField} 构造出来}，
只要求一个环自同态。
故那份精确性是\textbf{模型的性质}，而不是公理被证明能实现的东西。
\textbf{从"框架有一个具此性质的标度平移"到"$\hbar$ 就是那个步长"，
是一次\emph{认定}，不是一次构造。}

其余五处：\texttt{Waves.Conserved}（多极守恒未从 Bianchi 桥接到具体矩）、
\texttt{RG.ScaleFlow}、\texttt{Direction.DirTransport}、
\texttt{Invariant.Trace}、\texttt{Covariance.CosmicHistory}——
各自从未由框架自己的载体实例化。

\textbf{这些没有一条是假的。}每一条对满足其前提者都成立。
此处更正的是\emph{读法}：它们说的是\emph{会跟出什么}，不是公理\emph{给出了什么}。
\end{mdframed}

\subsubsection*{补记：六个结构现在都被构造出来了}

$\mathbb R[X]$ 配 $d/dX$ 满足 A1，而平移 $X\mapsto X+1$ 是其上的一个
\texttt{ScaleShift}，且 $\Delta X = 1$——\textbf{步长是精确且有限的，不是一阶截断}。
其余五个也各有见证：矩阵迹、伴随作用、非退化的平移流、常值矩、空间常值漂移。

\textbf{为什么"乘以标度因子"不行、而平移行：}对一个单位 $u$，
$(ux)(uy)=u^2xy\neq u(xy)$——\textbf{乘法是可加的但不是可乘的}，故不是环同态。
\textbf{标度平移是平移，不是重标度}，这一点此前没说清。

\begin{mdframed}[linecolor=red!60!black,linewidth=1.2pt,backgroundcolor=red!4]
\textbf{这改变了什么，又没改变什么。}
定理现在\textbf{已知非空转}，结构被证明与 A1 \emph{相容}而非仅仅被假设。

\textbf{但这不是推导}：一个同时承载两种结构的模型，
并不表明 A1 的\emph{每个}模型都必须承载一个标度平移。
故 $\hbar$ 与步长的等同\textbf{仍是一次认定}，继续留在本登记册中。
\textbf{"相容"与"被迫"之间的距离是真实的，此处没有被弥合。}
\end{mdframed}

\subsection{审查方法本身}

值得记录的是，本轮两处最有价值的修正都来自同一个动作：
\textbf{发现被借走的是"问题"而不是"答案"}。
"要不要无挠条件"预设了标架与联络是独立场；
"该选哪个规范群"预设了存在一个内禀空间。
换掉问句，答案自己就出来了（第 \ref{sec:transport}、\ref{sec:gauge} 节）。
本节登记的六条"仍存在"，很可能也各自藏着一个同类的问句。

\section{（已撤回的解读）块单值性、禁闭与分数电荷}\label{sec:color}

\begin{mdframed}[linecolor=red!60!black,linewidth=1.2pt,backgroundcolor=red!4]
\textbf{本节的定理成立，本节的\emph{解读}已被撤回。}

下面导出的"绕数 $+$ 块置换"这对单值性数据是对的，禁闭与分数电荷的结论
在\emph{它自己的假设下}也是对的。被撤回的是把这个有限荷读作\textbf{色}。

\textbf{而"它给出 2 而数据要 3，这是一个缺口"这个后续说法同样已过时。}
第 \ref{sec:particle} 节证明：序参量是射影的，故块群只能是 $\mathbb Z/2$，
\textbf{3 从来就不在供给里}。这不是缺口，是被解释了的结构不可能——
色不是这个量子数，而且没有任何修补能让它是。

保留本节，是因为它的数学仍被审计，而撤回记录与结论同样重要。
\textbf{请勿把以下内容当作本框架的现行主张。}
\end{mdframed}

第 \ref{sec:defect} 节只给了缺陷一个整数：标度绕它的绕数。那是不完整的。
若缺陷位于一个\emph{简并块}中——一组携带相同局域单位的方向——
则绕它一周时，标度不仅会平移，\emph{块内标号也可能被置换}，
因为块内的标号本来就不是物理的（定理 \ref{thm:multiplet}）。

所以缺陷的单值性数据是一\emph{对}：一个绕数\textbf{和}一个块置换。
这里没有从规范理论抄任何东西；置换是被"简并块没有优先标号"这一事实\emph{逼出来}的。

以下全部来自一个问题：这样的位形何时是整体自洽的——标号在绕行一周后是否闭合。

\begin{theorem}[禁闭]\label{thm:confine}
其单值性真正置换了块的缺陷，\textbf{不可能单独存在}。
这不是"力随距离增长"，而是孤立位形根本不是单值的，
因而那里并没有一个东西可供分离。
\leanline{Color.BlockMonodromy.confinement}
\end{theorem}

\begin{theorem}[束缚态大小即块大小]\label{thm:baryon}
$k$ 个拷贝闭合当且仅当 $\pi^k=1$。对一个把大小为 $n$ 的块循环一次的缺陷，
最小自洽的组分数恰为 $n$。\textbf{重子数就是块大小}，不是额外的量子数。
\leanline{Color.BlockMonodromy.bound\_state\_size}
\end{theorem}

\begin{theorem}[介子无需解释；不存在部分强子]\label{thm:meson}
缺陷与其反缺陷总是闭合，对 $n$ 无任何条件；
而少于 $n$ 个组分\emph{永不}闭合。
取 $n=3$：一个组分不可观测，两个不可观测，三个可观测。
"没有双夸克强子"不是一条附加规则。
\leanline{Color.BlockMonodromy.pair\_observable,
Color.BlockMonodromy.no\_partial\_state}
\end{theorem}

\begin{theorem}[电荷成 $1/n$]\label{thm:third}
束缚态电荷是单个组分的 $n$ 倍：
\[
Q(\text{束缚态}) = n\cdot Q(\text{组分}),\qquad
Q(\text{组分}) = \tfrac{1}{n}\,Q(\text{束缚态}).
\]
可观测态的电荷必为整数（它是可观测态的电荷），故组分电荷落在 $\tfrac1n\mathbb{Z}$ 中。
\leanline{Color.BlockMonodromy.constituent\_charge\_fraction,
constituent\_charge\_rat, bound\_charge\_divisible}
\end{theorem}

\begin{mdframed}[linecolor=axcol,linewidth=0.9pt,backgroundcolor=axcol!4]
取 $n=3$ 即观测到的花样：每个束缚态三个组分、单个组分不可观测、电荷成三分之一。
\textbf{框架并没有把 3 放进去}——它放进去的是"一个大小为 $n$ 的简并块"，
取出来的是"$n$ 个组分"与"$1/n$ 电荷"。
\end{mdframed}

必须明说：\textbf{自然界为何选了 $n=3$，本报告没有导出}。
那需要把标度花样本身导出来。

\section{审查：色荷的块不是空间的块}\label{sec:coloraudit}

第 \ref{sec:axis} 节靠"序参量是无取向轴"让点缺陷回来，
而那要求\textbf{单轴}花样：一个方向被挑出，其余简并。
在三维空间中，简并的那部分恰有\textbf{两个}成员。

第 \ref{sec:color} 节则证明：块大小为 $n$ 时束缚态有 $n$ 个组分、电荷成 $1/n$。

把两者放在一起，很诱人做一个认同：\emph{色荷的块就是空间简并块}。
那会是个漂亮结果——色将成为几何的推论而非输入。

\begin{mdframed}[linecolor=red!60!black,linewidth=0.9pt,backgroundcolor=red!3]
\textbf{它是错的，本节记录原因。}

三维中的单轴块大小为二。于是该认同预言\textbf{两组分}束缚态与\textbf{二分之一}电荷。
而重子有三个组分，夸克电荷成三分之一。
该认同因此被\emph{排除}——不是被削弱，是被排除，而且差得不近。
\leanline{ColorAudit.identification\_fails\_baryon,
ColorAudit.identification\_fails\_charge}
\end{mdframed}

幸存与不幸存：

\begin{itemize}[leftmargin=1.6em]
\item 第 \ref{sec:color} 节\textbf{未受影响}。它的定理是关于大小为 $n$ 的简并块的，
      对每个 $n$ 都对；它一直就说 $n$ 是输入。
\item 被排除的是那条\emph{把 $n$ 从空间标度花样导出来}的路线。该路线给出 $n=2$，被证伪。
\item 因此，若框架要有色，色\textbf{不是}空间方向的置换，
      它必须来自框架目前尚不包含的结构。
\end{itemize}

这同时消除了一个未曾言明的隐忧：若该认同成立，色将是时空对称性而非内禀对称性，
而那正是 Coleman--Mandula 所禁止的混合。\textbf{框架没能产生它，是预期结果而非意外。}

\clearpage
\section{（已撤回）三维中的绕数缺陷究竟是什么}\label{sec:codim}

\begin{mdframed}[linecolor=red!60!black,linewidth=1.2pt,backgroundcolor=red!4]
\textbf{本节的定理成立，但它们描述的不是本框架。}

以下余维数分析假设标度是\emph{圆值}的（对数标度取值于 $\mathbb R/\Delta\mathbb Z$）。
在那个假设下，结论——三维中绕数缺陷是\textbf{弦}而非点——是对的。

\textbf{但那个假设是错的，而错法又是那个熟悉的错法：}标度是一个\emph{比值}，
故读数无向，序参量是一根\textbf{无向轴}而非一个圆（第 \ref{sec:axis} 节）。
射影序参量的 $\pi_2$ 非平凡，故\textbf{点缺陷回来了}。
这是“标量代替方向性”五次中的第二次（第 \ref{sec:audit} 节）。

保留本节，是因为它记录了一次真实的误判，而登记册与结论同样重要。
\textbf{请勿把以下内容当作本框架对激发形态的主张}——
那由第 \ref{sec:axis}、\ref{sec:charges}、\ref{sec:particle} 节给出。
\end{mdframed}

试图从代数定住维数 $n$，结果撞出的是 \texttt{Defect.lean} 的一个问题，
而这个问题比维数本身更有意思。

第 \ref{sec:defect} 节说："绕任一\emph{包围一个点}的闭合回路一周，标度可能不闭合"。
在\emph{二维}空间里这句话没问题——回路确实包围一个点。
\textbf{在三维里它是错的}：三维空间中的回路不包围一个点，它包围一条\emph{线}。
回路能探测的是余维为二的轨迹，而三维中余维二是一维的。

\begin{mdframed}[linecolor=red!60!black,linewidth=0.9pt,backgroundcolor=red!3]
\textbf{所以三维空间中的圆值标度不产生点粒子，它产生\emph{弦}。}

这是被逼出来的，不是选的。标度是\emph{一个}数——一个比值——
所以它定义的场是圆值的，其缺陷由回路分类。
框架里没有第二个标度可以让缺陷在三维中成为点。
这个结论要么被接受，要么框架得改。
\end{mdframed}

\begin{theorem}[缺陷维数]\label{thm:defectdim}
$d$ 维空间中绕数缺陷的维数是 $d-2$。故
点状绕数缺陷\textbf{要求恰好二维空间}，即一个 $2+1$ 世界；
而在我们的三维中，缺陷是一维的。
\leanline{Codimension.pointlike\_iff\_two\_dim, three\_dim\_gives\_strings,
Codimension.no\_pointlike\_in\_four}
\end{theorem}

结合第 \ref{sec:signature} 节（漂移给出恰好一个时间方向），时空维数是空间维数加一，
于是两种情形被截然分开：\textbf{点状激发 $\Rightarrow$ $2+1$ 世界；
$3+1$ 世界 $\Rightarrow$ 激发是一维的}。二者不可兼得。

\subsection{这独立地解释了质量律为何不确定}

\begin{theorem}[绕数不决定质量]\label{thm:strmass}
一维缺陷的质量是张力乘长度。绕数相同而长度不同的两个缺陷质量不同，
故\textbf{没有任何仅依赖绕数的函数能充当质量律}。
\leanline{Codimension.string\_mass\_not\_topological,
Codimension.stringMass\_strictMono}
\end{theorem}

这条从缺陷的\emph{几何}出发，独立地抵达了第 \ref{sec:audit} 节
由考察候选律得到的同一结论，并且解释了\emph{原因}：
\textbf{那个对象是有广延的，而拓扑从未度量过它的广延。}
两条互不相干的路线得到同一个结果，这提高了那次撤回的可信度。

\subsection{框架现在面临的选择}

要么基本激发是有广延的——那么第 \ref{sec:defect} 节的点图景必须换成弦图景，
质量谱不是塔；要么标度不是圆值的——那么第 \ref{sec:defect} 节的离散标度不变性
与第 \ref{sec:color} 节的块单值性都要作废。

\textbf{我们记录这个析取而不做选择}，因为已证的东西还不足以判定它。
\leanline{Codimension.extended\_or\_low\_dimensional}

\begin{remark}[风险]
这是一个真实的预言，也是一次真实的风险。本框架按目前的构造说：
$3+1$ 维中的基本激发是\emph{一维的广延对象}。
若这是错的，圆值标度就是错的，\texttt{Defect.lean} 与 \texttt{Color.lean} 随之作废。
（"回路探测余维二"这条本身是标准拓扑事实，此处作为输入而非证明。）
\end{remark}

\section{这个框架究竟能说什么}\label{sec:expressive}

第 \ref{sec:crosscheck} 节结尾留下一个我一直在回避的问题：
若引力在这里不是耦合，那么"把引力与其他力统一"到底是什么意思？
怀疑是这个问法本身借来的。\textbf{它确实是}，而弄清原因给出的对本框架的刻画，
比再多一条结果都更有用。

标准的统一纲领要一个规范群，让所有力从中降下来。这预设了：
力是耦合、耦合是基本描述、"一"意味着一个包含其余的代数对象。
\textbf{这三条在这里都不成立。} 力不是原始量，局域单位的比较才是。
所以"统一诸力"是一个用本框架\emph{没有}的词汇提出的问题。

它\emph{能}回答的问题是：\textbf{什么是可表达的？}

整个开发只由三种操作搭成，别无其他——\emph{比较标度}、\emph{不对易}、\emph{计数}。
本节证明这三者确实就是全部词汇。

\begin{mdframed}[linecolor=red!60!black,linewidth=1.2pt,backgroundcolor=red!4]
\textbf{"三者互相独立"这个结论已被推翻，本节因此被降级为记录。}
下面关于\emph{可表达性}的定理仍然成立且仍被审计：
可观测量恰是标度\emph{差}的函数，绝对标度根本不可表达。

但由此得出的"标度路与转动路按构造独立，故不可能有跨扇区检验"是\textbf{错的}。
第 \ref{sec:coupling} 节证明：那是对\emph{标量}扇区的分析；
方向性标度下\textbf{转动扇区由标度扇区\emph{生成}}，耦合是被词义强制的。
这是"标量代替方向性"错误的第三次出现（共五次，见第 \ref{sec:audit} 节）。
\end{mdframed}

\subsection{标度这一路：恰好表达"差"}

\begin{theorem}[可观测 $\iff$ 差的函数]\label{thm:express}
一个量在 A5 的基准平移下不变（即可观测），\textbf{当且仅当}它是标度\emph{差}的函数。
不是"至少是"，是\emph{恰好是}——两个方向都成立。
\leanline{Expressive.invariant\_factors, diffPattern\_invariant,
invariant\_iff\_diffPattern}
\end{theorem}

\begin{theorem}[绝对标度根本不可表达]\label{thm:noabs}
在一点读取标度不是不变量，故它\textbf{不是本框架的可观测量}——
不是"未被测量"，而是\emph{在词汇之外}。
\leanline{Expressive.eval\_not\_invariant}
\end{theorem}

\subsection{两路互不通气}

第 \ref{sec:foundation} 节把一次比较典范地拆成标度（交换化）与转动（对易子子群）。
这两者携带\emph{不相交}的信息。

\begin{theorem}[两路独立]\label{thm:indepchan}
固定其中之一，对另一个毫无约束：
标度相同的两次比较相差一个任意的对易子元素；
而每个对易子的标度都平凡。
\leanline{Expressive.scale\_says\_nothing\_about\_rotation,
rotation\_says\_nothing\_about\_scale, channels\_independent,
neither\_channel\_complete}
\end{theorem}

\begin{mdframed}[linecolor=red!60!black,linewidth=0.9pt,backgroundcolor=red!3]
\textbf{诊断。} 第 \ref{sec:crosscheck} 节找不到跨扇区检验，\emph{那不是运气不好}。
两路是\textbf{按构造}独立的——标度路是一个商，转动路是它的核——
所以谁也约束不了谁。跨扇区检验需要它们互相约束，
\textbf{而公理保证它们不会}。
\end{mdframed}

\subsection{框架弱点的诚实表述，因此改变了}

不是"尚未导出标准模型"，而是：
\textbf{三条通路彼此太独立，以至于不能联合预言}。

"引力、$\hbar$、粒子内容是同一个变量的三种读法"——这句话是对的，
\emph{而它同时就是它们无法互相核对的原因}。

于是今后要找的东西变得具体得多：不是更大的对称群，
而是\textbf{一条通路间的耦合}——某种使标度差去约束对易子、
或使计数去约束二者之一的结构。A1--A6 不提供任何这样的东西，
第 \ref{sec:crosscheck} 节就是这一点的观测后果。
\leanline{Expressive.no\_channel\_coupling}

\begin{remark}[这一节改变了什么]
它没有增加任何物理预言，而是划定了\emph{边界}：
框架能说"标度之差"、"两次比较不对易多少"、"能分辨几个结构"，
除此之外说不出别的。
之前几轮里反复出现的"借来的问法"，
到这里有了一个统一的解释——凡是要求框架谈论
\emph{绝对标度}或\emph{两路之间的关系}的问题，都不是它的问题。
\end{remark}

\section{我提议的跨扇区检验并不存在}\label{sec:crosscheck}

第 \ref{sec:spectrum} 节从质量谱量出阈值密度 $\rho\approx0.86$ 每 e 倍。
而第 \ref{sec:running} 节有一个同名的 $\rho$ 控制 $g^{-2}$ 的斜率。
既然框架把两者都叫 $\rho$，显然该做的事就是：
一个从谱量、一个从耦合跑动量，看它们是否相等——
一次不需要任何新假设的真正跨扇区检验。

我做了。\textbf{它不是一次检验，而原因是我在第 \ref{sec:running} 节的措辞出了错。}

\paragraph{第一，数字不一致。}
数全部十二个标准模型质量标度得 $\rho_{\text{all}}\approx0.864$；
只数六个夸克——强耦合真正响应的那些——得 $\rho_{\text{quark}}\approx0.444$。
相差约两倍。
\leanline{CrossCheck.rhoAll\_value, rhoQuark\_value, densities\_differ,
density\_ratio\_near\_two}

\begin{mdframed}[linecolor=red!60!black,linewidth=0.9pt,backgroundcolor=red!3]
\paragraph{第二，也是致命的：它们根本不是同一个量。}
第 \ref{sec:running} 节写的是"探针响应它能分辨的缺陷"，
而我把它读成了\emph{全部}缺陷。规范耦合不是这样：
\textbf{它只响应带有该规范荷的结构}。
所以跑动里的 $\rho$ 是\emph{扇区受限}的（且为使符号正确，是带号的）计数，
而谱里的 $\rho$ 是对每一个阈值的裸计数。
两者没有理由相等，\textbf{而框架也从未说过它们相等}。
\end{mdframed}

\paragraph{第三，在单个扇区内该关系退化为恒等式。}
一个方程 $db/dt=\kappa\,(dN/dt)$、一个未知量 $\kappa$：
测出跑动跳变与阈值密度是在\emph{定义} $\kappa$，而不是检验什么。
对 QCD 在粲与顶阈值之间执行，返回 $\kappa=2/3$——
那正是每味的 $\beta$ 跳变倒着写一遍。
\leanline{CrossCheck.sector\_relation\_is\_identity,
qcd\_extraction\_returns\_input, no\_second\_determination}

\paragraph{幸存下来的。}
逐扇区的预言仍可检验：每个扇区的阈值应各自是常速率过程，速率不必相同。
单看夸克扇区，$\mathrm{CV}^2=0.243$（即 $\mathrm{CV}=0.493$），
落在五个间距时预言抽样分布的第 12 百分位——
偏低、比常速率过程通常给出的更\emph{规则}，但未被排除。
五个间距说不出更强的话。夸克扇区同样不是等比塔。
\leanline{CrossCheck.quarkCvSq\_value, quark\_not\_geometric}

\begin{mdframed}[linecolor=axcol,linewidth=0.9pt,backgroundcolor=axcol!4]
\textbf{记录在案的结论。} 本框架目前\textbf{不提供}任何关于阈值密度的跨扇区检验。
该提议是我的一次越界，第 \ref{sec:running} 节的措辞已更正为
"耦合只计它所响应的那些缺陷"。

真正的检验需要\emph{两个独立}的同一权重测定，好让后者核对前者。
框架只给得出一个，因为每个扇区只有一个类耦合量，
且没有任何一个耦合响应所有结构。
\end{mdframed}

\clearpage
\section{范围声明：什么被证明了，什么没有}\label{sec:scope}

这是本报告最重要的一节。Lean 验证的是\emph{蕴含关系}，不是自然界的真值。

\paragraph{已被机器证明（从公理出发，零 \texttt{sorry}）。}
定理 \ref{thm:adleib}、\ref{thm:adjac}、\ref{thm:adpow}、\ref{thm:robertson}、
\ref{thm:heisenberg}、\ref{thm:master}、\ref{thm:ric}、\ref{thm:rsc}、\ref{cor:flat}、
\ref{thm:weyl}、\ref{thm:fiducial}、\ref{thm:redshift}、\ref{thm:linear}、
\ref{thm:poisson}、\ref{thm:dual}、\ref{thm:cs}、\ref{thm:eft}、\ref{thm:lifetime}、
\ref{thm:detfail}、\ref{thm:dm}、\ref{thm:diss}、\ref{thm:desitter}、
\ref{thm:h}、\ref{thm:arrow}、\ref{thm:null}、\ref{thm:massive}、
\ref{thm:cov}、\ref{thm:covfield}、\ref{thm:sing}、\ref{thm:singcurv}、
\ref{thm:eos}、\ref{thm:starassoc}、\ref{thm:starcl}、\ref{thm:starcomm}、\ref{thm:poissonleib}、\ref{thm:ccr}、\ref{thm:frameflat}、\ref{thm:frameexists}、\ref{thm:winding}、\ref{thm:stable}、\ref{thm:pair}、\ref{thm:anti}、\ref{thm:unify}、\ref{thm:unifyiff}、\ref{thm:running}、\ref{thm:asympt}、\ref{thm:transmut}、\ref{thm:massratio}、\ref{thm:towerfail}、\ref{thm:commcov}、\ref{thm:bianchi}、\ref{thm:gradflat}、\ref{thm:rotpart}、\ref{thm:beta}、\ref{thm:exchange}、\ref{thm:commiff}、\ref{thm:twisted}、\ref{thm:stab}、\ref{thm:multiplet}、\ref{thm:tunique}、\ref{thm:tcoset}、\ref{thm:tiff}、\ref{thm:confine}、\ref{thm:baryon}、\ref{thm:meson}、\ref{thm:third}、\ref{thm:massgap}、\ref{thm:redshiftobs}、\ref{thm:scalecomm}、\ref{thm:curvinvis}、\ref{thm:splitcanon}、\ref{thm:sig}、\ref{thm:notime}、\ref{thm:sigfid}、\ref{thm:formsymm}、\ref{thm:forminv}、\ref{thm:half}、\ref{thm:dircurv}、\ref{thm:flatlie}、\ref{thm:abcase}、\ref{thm:expose}、\ref{thm:recexp}、\ref{thm:defectdim}、\ref{thm:strmass}、\ref{thm:recip}、\ref{thm:noflat}、\ref{thm:even}、\ref{thm:patterns}、\ref{thm:onecharge}、\ref{thm:indep}、\ref{thm:precoeff}、\ref{thm:cancel}、\ref{thm:vacstat}、\ref{thm:ricciderived}、\ref{thm:nolist}、\ref{thm:masspos}、\ref{thm:eqzero}、\ref{thm:notgeom}、\ref{thm:express}、\ref{thm:noabs}、\ref{thm:indepchan}、\ref{thm:ss}、\ref{thm:cartan}、\ref{thm:concrete}、\ref{thm:onedir}、\ref{thm:noflatcomm}、\ref{thm:commcost}、\ref{thm:conserv}、\ref{thm:nofe}、\ref{thm:multflat}、\ref{thm:rotwit}、\ref{thm:multthresh}、\ref{thm:orth}、\ref{thm:bbr}、\ref{thm:rank1}、\ref{thm:samefact}、\ref{thm:binj}、\ref{thm:mix}、\ref{thm:mixge}、\ref{thm:blk2}、\ref{thm:odd}、\ref{thm:qcons}、\ref{thm:newc}、\ref{thm:nouniv}、\ref{thm:wedge}、\ref{thm:uniax}、\ref{thm:three}、\ref{thm:rotvar}、\ref{thm:wound}、\ref{thm:noax2}、\ref{thm:stabpos}、\ref{thm:algbig}、\ref{thm:twocond}、\ref{thm:ratio}、\ref{thm:dennorm}、\ref{thm:a7}、\ref{thm:noinv}、\ref{thm:iso}、\ref{thm:noorigin}、\ref{thm:dm2}、\ref{thm:reso}、\ref{thm:ncdefm}、\ref{thm:ncclass}、\ref{thm:ncmaster}、\ref{thm:ncsym}、\ref{thm:nohor}、\ref{thm:gw}、\ref{thm:nodisp}、\ref{thm:ent}、\ref{thm:nomin}、\ref{thm:sep}、\ref{thm:interval}、\ref{thm:noglobal}、\ref{thm:parity}、\ref{thm:nomax}、\ref{thm:bound}、\ref{thm:counting}、\ref{thm:afdensity}、\ref{thm:spectrum}。

\paragraph{尚未完成——本纲领当前最薄弱之处。}

\begin{mdframed}[linecolor=axcol,linewidth=1pt,backgroundcolor=axcol!5]
\textbf{完整的登记册在第 \ref{sec:audit} 节，分五张清单：}
已撤回（\S\ref{ss:retracted}，4 条）、已限缩适用域（\S\ref{ss:corrected}，6 条）、
\textbf{已解决}（\S\ref{ss:closed}，10 条，附由何关闭）、
仍然开放（\S\ref{ss:open}，7 条）、引用与假设（\S\ref{ss:cited}，各 2 条）。

本栏此前混杂着已解决与未解决的条目并以删除线标注，那是修改痕迹而非登记册；
\textbf{已解决的条目现已移入 \S\ref{ss:closed} 的表中}，本栏只留仍未完成者。
\end{mdframed}
\begin{itemize}[leftmargin=1.6em]
\item \textbf{标准模型不是目标}（立场，非缺口）。索取一张固定物种表，
      是把标准模型的\emph{本体论}也借过来；而 A3 没有它——内容是标度上的截面，
      没有哪个标度携带完整清单（第 \ref{sec:emergence} 节）。
      色、代结构、Higgs 因此不是本框架应当重现的东西。
      应当索取的是\textbf{阈值结构}，第 \ref{sec:spectrum} 节已给出并检验。

\item \textbf{量子场论不是目标}（立场，非缺口）。此前本栏写作“量子场论未被回归”，
      \textbf{那个措辞不准确}：二次量子化、路径积分、圈图重整化\emph{预设}了
      一个固定内容的场论描述，而第 \ref{sec:slice} 节证明那种描述的标签数
      度量的是它覆盖的\emph{范围}而非物理的深度。
      \textbf{本框架回归的是量子\emph{力学}}（对易子、幂律、不确定性，
      第 \ref{sec:nc} 节）与\emph{精确}的非交换标度平移（$\hbar$ 是步长）。
      真正的缺口是\textbf{散射振幅}——它需要固定内容的渐近态，
      而框架只能给出“两个被指名标度之间的映射”。

\item \textbf{强扇区没有纯数}。第 \ref{sec:qcd} 节证明量纲嬗变、渐近自由
      与无参数质量比；第 \ref{sec:running} 节把 $\beta$ 系数换成缺陷密度并导出线性形式。
      \textbf{仍缺的只有密度的\emph{数值}}——而那与"唯一一个自由单位"是同一件事
      （第 \ref{sec:dimension} 节），不是一处独立的缺口。
      任何强子质量的纯数需要格点或解析求解，本框架不提供也不声称提供。
\end{itemize}

\paragraph{属于物理输入而非定理（续）。}

\paragraph{属于物理输入而非定理。}
\begin{itemize}[leftmargin=1.6em]
\item \textbf{A6$'$ 标度响应律}（$e^{2\sigma}R=\kappa\rho$）。本报告\emph{未}从更基本的
      原理导出它；它扮演 Einstein 方程的角色。所证明的是：一旦接受它，
      Newton 极限被强制。$\kappa=16\pi G$ 是外部标定。

\item \textbf{耗散率 $\lambda$ 的\emph{数值}}。第 \ref{sec:cosmos} 节已证明
      A5 强制 $\lambda$ \textbf{恒定}——故 $w=-1$ 且不演化，
      这是本框架最暴露的预言。但 $\lambda$ 的\emph{数值}不被预言，
      也\textbf{不可能}被预言：它就是那个唯一的单位（第 \ref{sec:dimension} 节）。
\item \textbf{暗物质的具体候选}。定理 \ref{thm:dm} 证明“引力可见 $+$ 探测不可见”
      在标度层级下是相容的，并给出构造；它\emph{不}预言具体的粒子谱或截面数值。
\item \textbf{重子不对称度}。仅证明 Sakharov 第三条件结构性满足，
      未导出 $\eta_B$ 的量级。
\item \textbf{胶球等强耦合谱}。本报告未涉及 QCD 谱的计算；
      引言所述的实验进展是启发，不构成本报告的证据链条。
\end{itemize}

\paragraph{形式化的技术边界。}
基底采用代数化的“标度代数”而非光滑流形上的解析结构。这使全部曲率恒等式
成为任意交换环上的代数恒等式（更强、更干净），但也意味着我们\emph{没有}
形式化测地线方程、整体解的存在性、或 Cauchy 问题。
“类光测地线的共形不变性”本报告只证明了其代数内核——类光锥的标度无关性
（定理 \ref{thm:null}），未证明参数化曲线层面的完整陈述。

\section{形式化与验证}

\subsection{规模}

\begin{center}
\begin{longtable}{lrl}
\toprule
文件 & 定理数 & 内容 \\
\midrule
\endhead
\multicolumn{3}{l}{\emph{第一部分\quad 公理}}\\
\texttt{Basic.lean} & 13 & 标度代数、Kronecker 收缩原语 \\
\texttt{Axioms.lean} & 16 & A2/A3/A5、红移、基准协变性 \\
\texttt{Postulates.lean} & 3 & A1--A7 全在一处；证明基底只有一个载体 \\
\midrule
\multicolumn{3}{l}{\emph{第二部分\quad 公理逼出什么}}\\
\texttt{Foundation.lean} & 10 & 地基重建：标度即交换化，拆分被导出 \\
\texttt{Signature.lean} & 7 & 时间即漂移方向；号差由耗散导出 \\
\texttt{Invariant.lean} & 6 & 记账形式由迹导出 \\
\texttt{Quantum.lean} & 7 & 对易子即导数、Heisenberg 不确定性 \\
\texttt{Direction.lean} & 8 & 方向是定义表示（原说法已更正） \\
\texttt{Connection.lean} & 9 & 曲率即对易子、Bianchi \\
\texttt{Coupling.lean} & 8 & 两次缩放生成转动；耦合被词义强制 \\
\texttt{Algebra.lean} & 18 & 实秩一；筛出 Lorentz 族 \\
\texttt{Dimension.lean} & 10 & 方向是表示；比值不变；量值即单位 \\
\texttt{Slice.lean} & 16 & 广延即标签；转动是楔积；$k=3$ 唯一 \\
\texttt{Axes.lean} & 9 & 两难消解；轴数被迫为一 \\
\midrule
\multicolumn{3}{l}{\emph{第三部分\quad 分叉点}}\\
\texttt{Locus.lean} & 11 & A7；标签只活在各向异性轨迹上 \\
\texttt{Dynamics.lean} & 13 & 守恒出自 Bianchi \\
\midrule
\multicolumn{3}{l}{\emph{第四部分\quad 各向同性扇区}}\\
\texttt{Conformal.lean} & 20 & 共形曲率恒等式、Ricci、标量曲率 \\
\texttt{Frame.lean} & 5 & A4$'$ 方向性；旧公理禁止 Schwarzschild \\
\texttt{Unify.lean} & 9 & A4 是 A4$'$ 的各向同性轨迹 \\
\texttt{Newton.lean} & 7 & 线性化、Poisson、对偶数模型 \\
\texttt{Schwarzschild.lean} & 5 & 使倒数关系成立的相消 \\
\texttt{RicciDiag.lean} & 4 & 由联络导出 Ricci，引力链条闭合 \\
\texttt{Vacuum.lean} & 5 & 倒数关系的常数由渐近平直定出 \\
\texttt{Deflection.lean} & 7 & 各向同性扇区恰好只给一半 \\
\texttt{PPN.lean} & 9 & $\gamma=1$；付清另一半 \\
\texttt{Precession.lean} & 8 & 水星 $42.99''$ \\
\texttt{Covariance.lean} & 7 & 耗散下的局域协变性 \\
\texttt{Singularity.lean} & 6 & 无度规退化 \\
\texttt{Waves.lean} & 16 & 旋近动力学：守恒杀掉单极与偶极 \\
\texttt{Horizon.lean} & 12 & 无视界；波在光锥；熵是对数；无最小长度 \\
\midrule
\multicolumn{3}{l}{\emph{第五部分\quad 各向异性扇区}}\\
\texttt{Gauge.lean} & 5 & 内禀对称性即标度花样稳定子 \\
\texttt{Transport.lean} & 6 & 输运由花样定出；规范自由即简并 \\
\texttt{Axis.lean} & 9 & 点粒子回来了 \\
\texttt{Defect.lean} & 13 & 标度缺陷、绕数即荷 \\
\texttt{Charges.lean} & 15 & 两个拓扑荷，恰一个被禁闭 \\
\texttt{Particle.lean} & 16 & 标签结构 $(\mathbb R,\mathbb Z/2,\mathbb Z)$ \\
\texttt{Emergence.lean} & 9 & 没有基本粒子 \\
\midrule
\multicolumn{3}{l}{\emph{第六部分\quad 量子扇区}}\\
\texttt{Deformation.lean} & 12 & 一阶形变、对应原理、CCR \\
\texttt{Crossed.lean} & 11 & 精确非交换；$\hbar$ 是标度步长 \\
\texttt{NCConformal.lean} & 18 & 非交换曲率；修正 $=[\sigma_i,\sigma_j]$，纯反对称 \\
\midrule
\multicolumn{3}{l}{\emph{第七部分\quad 标度流与内容}}\\
\texttt{RG.lean} & 13 & 标度流、Callan--Symanzik、EFT 塔 \\
\texttt{Running.lean} & 6 & 由计数导出线性跑动 \\
\texttt{QCD.lean} & 10 & 渐近自由、量纲嬗变 \\
\texttt{Spectrum.lean} & 9 & 阈值密度、$\mathrm{CV}$ 检验、换算 \\
\texttt{Entropy.lean} & 9 & H 定理、不可逆性 \\
\texttt{Light.lean} & 6 & 类光锥标度无关 \\
\texttt{Observation.lean} & 6 & 红移属于观察者而非光 \\
\texttt{Particles.lean} & 9 & 标度比、寿命、阈值 \\
\texttt{DarkMatter.lean} & 4 & 探测失效 \\
\texttt{DarkEnergy.lean} & 8 & 耗散唯一性、de Sitter \\
\midrule
\multicolumn{3}{l}{\emph{第八部分\quad 预言与数据}}\\
\texttt{Predictions.lean} & 17 & $w(q)$ 状态方程、数值核对 \\
\texttt{Cosmos.lean} & 13 & $w$ 不演化；无初值；暗物质两路；共振宽度界 \\
\texttt{Native.lean} & 12 & 框架自己的七个问题 \\
\texttt{Data.lean} & 19 & 与 DESI／PDG／Fermi-LAT 对质 \\
\midrule
\multicolumn{3}{l}{\emph{第九部分\quad 登记册}}\\
\texttt{Color.lean} & 17 & 块单值性（色解释已撤回） \\
\texttt{ColorAudit.lean} & 8 & 块 $\ne$ 色块 \\
\texttt{Codimension.lean} & 12 & （已被方向性 A4 取代） \\
\texttt{Expressive.lean} & 10 & 可表达性（独立性结论已推翻） \\
\texttt{CrossCheck.lean} & 9 & 跨扇区检验并不存在 \\
\texttt{MassAudit.lean} & 9 & 拓扑不决定质量律（撤回） \\
\texttt{Audit.lean} & 4 & 登记册：撤回／更正／引用／假设 \\
\midrule
\textbf{合计} & \textbf{611} & 62 个文件，12055 行 Lean 代码 \\
\bottomrule
\end{longtable}
\end{center}

\subsection{公理审计}

对全部 637 条定理执行 \texttt{\#print axioms}，结果全部落在
\[
\{\texttt{propext},\ \texttt{Classical.choice},\ \texttt{Quot.sound}\}
\]
之内——这是 Lean 中经典数学的三条标准逻辑公理。
\textbf{无一条定理依赖 \texttt{sorryAx}}，即全部证明完整闭合。

需要强调：\texttt{ScaleAlgebra}、\texttt{ScaleField} 以及各条物理假设
\emph{不是} Lean 的 \texttt{axiom} 声明，而是类型类、结构与显式假设。
它们出现在定理的\emph{陈述}中，因而无法被隐藏在证明背后。
本开发中不存在“先把物理结论声明为公理、再宣称已证明”的情形。

\subsection{复现}

\begin{verbatim}
    lake exe cache get
    lake build
    lake env lean SCD/Verify.lean   # 公理审计
\end{verbatim}

\section{结语}

若把标度而非时空当作基本量，则一条链条从公理一路走下来：

\begin{itemize}[leftmargin=1.6em,itemsep=2pt]
\item 标度必须\textbf{有方向}（标量版本禁止 Schwarzschild）；
\item 一旦有方向，\textbf{标度与转动的耦合就被词义强制}——
      缩放对称、转动反称，而两个对称算子的括号必然反称；
\item 耦合加上"漂移是单个泛函"给出\textbf{实秩一}，从而 Lorentz 代数族，
      从而 $n=k+1$；$k=3$ 出自 A7，而 A7 是一条\emph{关于什么算可观测}的承诺；
\item 各向同性／异性的\textbf{分叉}决定了标签是否存在，
      而这是标度相关的事实——这就是"相互作用随能标下降而分离"；
\item 曲率是标度形变（各向同性扇区）；$\hbar$ 是标度轴上的步长；
      几何的全部量子修正是 $[\sigma_i,\sigma_j]$，且纯反对称；
\item 引力扇区定量闭合：$1.7515''$ 与 $42.99''/$世纪，中间无假设。
\end{itemize}

\medskip
\noindent
这些不是隐喻。它们是 595 条经机器检验、零 \texttt{sorry} 的定理，
且\textbf{全部}经 \texttt{\#print axioms} 审计（清单由源码生成，非人工挑选）。

\medskip
\noindent\textbf{但"它回归出了我们已知的物理"这句话，现在说不得。}
记分牌上有一条\emph{活张力}——DESI 对 $w$ 不演化给出 $2.8$--$4.2\sigma$ 的反向偏好，
而那是与 A5 的张力，不是与某个拟合参数的张力，\textbf{不存在带演化 $w$ 的本框架版本}；
还有一个\emph{活的开放问题}——$f_0(500)$ 极点落在宽度界的哪一侧，
两个旗舰确定值分居两侧，而框架押其中一个。

\medskip
\noindent
同样说不得的是"它解释了标准模型"：第 \ref{sec:emergence} 节论证那不是目标，
第 \ref{sec:audit} 节记录了为得到现有结论而必须\textbf{收回}的 4 条主张、
必须\textbf{限缩适用域}的 6 条、\textbf{引用而未证}的 2 条、
以及\textbf{登记为假设}的 2 条——外加同一个"标量代替方向性"的错误出现了 5 次，
其中一次挺过了一次形式化和一版报告。

\medskip
\noindent
形式化系统能回答的只有一件事：\textbf{这些结论确实从这些公理跟出来，中间没有缺口。}
公理是否为自然界所采纳，它不能回答，也不应假装能回答。
而它\emph{不能}回答这件事本身，正是记分牌与登记册存在的理由。




\clearpage
\appendix
\section{定理索引：全部 637 条}\label{sec:index}
\addcontentsline{toc}{section}{附录：定理索引}

本索引由 \texttt{tools/sync\_report.py} 从源码生成，非人工维护。
\textbf{每一条都在 \texttt{Verify.lean} 中经 \texttt{\#print axioms} 审计}，
故报告正文未展开叙述的定理\emph{并非未被检验}，只是未被叙述。

\footnotesize
\raggedright

\subsection*{第一部分\quad 公理}
\noindent\textbf{\texttt{Basic.lean}} (13)\quad \texttt{d\_\allowbreak zero}, \texttt{d\_\allowbreak one}, \texttt{d\_\allowbreak neg}, \texttt{d\_\allowbreak sub}, \texttt{d\_\allowbreak sum}, \texttt{d\_\allowbreak natCast}, \texttt{d\_\allowbreak nsmul}, \texttt{kron\_\allowbreak self}, \texttt{kron\_\allowbreak symm}, \texttt{d\_\allowbreak kron}, \texttt{d\_\allowbreak kron\_\allowbreak mul}, \texttt{sum\_\allowbreak kron\_\allowbreak left}, \texttt{sum\_\allowbreak kron\_\allowbreak right}\par\smallskip
\noindent\textbf{\texttt{Axioms.lean}} (16)\quad \texttt{en\_\allowbreak mul\_\allowbreak scale}, \texttt{scale\_\allowbreak mul\_\allowbreak en}, \texttt{d\_\allowbreak en}, \texttt{sig\_\allowbreak add\_\allowbreak const}, \texttt{hess\_\allowbreak add\_\allowbreak const}, \texttt{lap\_\allowbreak add\_\allowbreak const}, \texttt{gradsq\_\allowbreak add\_\allowbreak const}, \texttt{Chr\_\allowbreak add\_\allowbreak const}, \texttt{Rm\_\allowbreak add\_\allowbreak const}, \texttt{Ric\_\allowbreak add\_\allowbreak const}, \texttt{Defm\_\allowbreak add\_\allowbreak const}, \texttt{geometry\_\allowbreak fiducial\_\allowbreak invariant}, \texttt{sig\_\allowbreak neg}, \texttt{hess\_\allowbreak neg}, \texttt{lap\_\allowbreak neg}, \texttt{gradsq\_\allowbreak neg}\par\smallskip
\noindent\textbf{\texttt{Postulates.lean}} (3)\quad \texttt{scaleAlgebra\_\allowbreak is\_\allowbreak diffRing\_\allowbreak d}, \texttt{three\_\allowbreak from\_\allowbreak A7}, \texttt{A7\_\allowbreak holds\_\allowbreak only\_\allowbreak at\_\allowbreak three}\par\smallskip

\subsection*{第二部分\quad 公理逼出什么}
\noindent\textbf{\texttt{Foundation.lean}} (10)\quad \texttt{scale\_\allowbreak mul}, \texttt{scale\_\allowbreak comm}, \texttt{scale\_\allowbreak inv}, \texttt{scale\_\allowbreak commutator}, \texttt{curvature\_\allowbreak is\_\allowbreak scale\_\allowbreak invisible}, \texttt{scale\_\allowbreak eq\_\allowbreak one\_\allowbreak of\_\allowbreak mem\_\allowbreak commutator}, \texttt{mem\_\allowbreak commutator\_\allowbreak of\_\allowbreak scale\_\allowbreak eq\_\allowbreak one}, \texttt{scale\_\allowbreak eq\_\allowbreak iff\_\allowbreak differ\_\allowbreak by\_\allowbreak rotation}, \texttt{commutator\_\allowbreak eq\_\allowbreak one\_\allowbreak of\_\allowbreak comm}, \texttt{comm\_\allowbreak of\_\allowbreak scale\_\allowbreak injective}\par\smallskip
\noindent\textbf{\texttt{Signature.lean}} (7)\quad \texttt{codim\_\allowbreak ker\_\allowbreak eq\_\allowbreak one}, \texttt{no\_\allowbreak time\_\allowbreak of\_\allowbreak no\_\allowbreak drift}, \texttt{spatial\_\allowbreak eq\_\allowbreak top\_\allowbreak of\_\allowbreak no\_\allowbreak drift}, \texttt{spatial\_\allowbreak ne\_\allowbreak top\_\allowbreak of\_\allowbreak drift}, \texttt{drift\_\allowbreak direction\_\allowbreak unique}, \texttt{spatial\_\allowbreak smul}, \texttt{signature\_\allowbreak from\_\allowbreak dissipation}\par\smallskip
\noindent\textbf{\texttt{Invariant.lean}} (6)\quad \texttt{map\_\allowbreak zero}, \texttt{map\_\allowbreak sub}, \texttt{form\_\allowbreak symm}, \texttt{form\_\allowbreak invariant}, \texttt{form\_\allowbreak ad\_\allowbreak skew}, \texttt{form\_\allowbreak commutator\_\allowbreak self}\par\smallskip
\noindent\textbf{\texttt{Quantum.lean}} (7)\quad \texttt{ad\_\allowbreak zero\_\allowbreak left}, \texttt{ad\_\allowbreak leibniz}, \texttt{ad\_\allowbreak jacobi}, \texttt{ad\_\allowbreak comm\_\allowbreak of\_\allowbreak commute}, \texttt{ad\_\allowbreak pow}, \texttt{robertson}, \texttt{heisenberg}\par\smallskip
\noindent\textbf{\texttt{Direction.lean}} (8)\quad \texttt{D\_\allowbreak zero\_\allowbreak dir}, \texttt{D\_\allowbreak neg\_\allowbreak dir}, \texttt{curv\_\allowbreak self}, \texttt{curv\_\allowbreak antisymm}, \texttt{flat\_\allowbreak iff\_\allowbreak lie\_\allowbreak hom}, \texttt{curv\_\allowbreak of\_\allowbreak abelian}, \texttt{commute\_\allowbreak iff\_\allowbreak bracket\_\allowbreak acts\_\allowbreak trivially}, \texttt{transport\_\allowbreak determined\_\allowbreak by\_\allowbreak generators}\par\smallskip
\noindent\textbf{\texttt{Connection.lean}} (9)\quad \texttt{D\_\allowbreak sub}, \texttt{F\_\allowbreak antisymm}, \texttt{comm\_\allowbreak covD}, \texttt{bianchi}, \texttt{curv\_\allowbreak of\_\allowbreak central}, \texttt{curv\_\allowbreak gradient\_\allowbreak eq\_\allowbreak zero}, \texttt{curv\_\allowbreak split}, \texttt{curv\_\allowbreak eq\_\allowbreak rotation\_\allowbreak part}, \texttt{curv\_\allowbreak pure\_\allowbreak scale\_\allowbreak eq\_\allowbreak zero}\par\smallskip
\noindent\textbf{\texttt{Coupling.lean}} (8)\quad \texttt{bracket\_\allowbreak scaling\_\allowbreak scaling}, \texttt{bracket\_\allowbreak rotation\_\allowbreak scaling}, \texttt{bracket\_\allowbreak rotation\_\allowbreak rotation}, \texttt{coupling\_\allowbreak structure}, \texttt{bracket\_\allowbreak self}, \texttt{no\_\allowbreak rotation\_\allowbreak of\_\allowbreak commuting}, \texttt{boost\_\allowbreak bracket\_\allowbreak eq\_\allowbreak rotation}, \texttt{scalings\_\allowbreak do\_\allowbreak not\_\allowbreak commute}\par\smallskip
\noindent\textbf{\texttt{Algebra.lean}} (18)\quad \texttt{scaling\_\allowbreak rotation\_\allowbreak orthogonal}, \texttt{boost\_\allowbreak zero\_\allowbreak zero}, \texttt{boost\_\allowbreak zero\_\allowbreak succ}, \texttt{boost\_\allowbreak succ\_\allowbreak zero}, \texttt{boost\_\allowbreak succ\_\allowbreak succ}, \texttt{boost\_\allowbreak isScaling}, \texttt{boost\_\allowbreak add}, \texttt{boost\_\allowbreak smul}, \texttt{boost\_\allowbreak mul\_\allowbreak succ\_\allowbreak succ}, \texttt{boost\_\allowbreak bracket\_\allowbreak spatial}, \texttt{boosts\_\allowbreak commute\_\allowbreak iff\_\allowbreak parallel}, \texttt{commuting\_\allowbreak boosts\_\allowbreak lie\_\allowbreak on\_\allowbreak a\_\allowbreak line}, \texttt{commuting\_\allowbreak family\_\allowbreak one\_\allowbreak parameter}, \texttt{no\_\allowbreak two\_\allowbreak dimensional\_\allowbreak commuting\_\allowbreak family}, \texttt{boost\_\allowbreak drift}, \texttt{commuting\_\allowbreak drifts\_\allowbreak proportional}, \texttt{boost\_\allowbreak spatial\_\allowbreak block\_\allowbreak zero}, \texttt{boost\_\allowbreak eq\_\allowbreak zero\_\allowbreak iff}\par\smallskip
\noindent\textbf{\texttt{Dimension.lean}} (10)\quad \texttt{algebra\_\allowbreak bigger\_\allowbreak than\_\allowbreak directions}, \texttt{rep\_\allowbreak dim\_\allowbreak determined}, \texttt{rot\_\allowbreak matches\_\allowbreak scale\_\allowbreak iff}, \texttt{alg\_\allowbreak matches\_\allowbreak rep\_\allowbreak iff}, \texttt{two\_\allowbreak matching\_\allowbreak conditions\_\allowbreak differ}, \texttt{wedge\_\allowbreak smul\_\allowbreak both}, \texttt{wedge\_\allowbreak ratio\_\allowbreak invariant}, \texttt{only\_\allowbreak ratios\_\allowbreak are\_\allowbreak fixed}, \texttt{density\_\allowbreak normalization\_\allowbreak relation}, \texttt{normalization\_\allowbreak returns\_\allowbreak input}\par\smallskip
\noindent\textbf{\texttt{Slice.lean}} (16)\quad \texttt{mem\_\allowbreak newContent}, \texttt{newContent\_\allowbreak nonempty\_\allowbreak of\_\allowbreak threshold}, \texttt{newContent\_\allowbreak mono}, \texttt{newContent\_\allowbreak empty\_\allowbreak of\_\allowbreak no\_\allowbreak threshold}, \texttt{no\_\allowbreak finite\_\allowbreak universal\_\allowbreak labelling}, \texttt{label\_\allowbreak count\_\allowbreak measures\_\allowbreak range}, \texttt{wedge\_\allowbreak eq\_\allowbreak bracket}, \texttt{wedge\_\allowbreak antisymm}, \texttt{wedge\_\allowbreak bilinear\_\allowbreak left}, \texttt{wedge\_\allowbreak eq\_\allowbreak zero\_\allowbreak iff\_\allowbreak parallel}, \texttt{homogeneous\_\allowbreak carries\_\allowbreak no\_\allowbreak rotation}, \texttt{rotation\_\allowbreak requires\_\allowbreak biaxial}, \texttt{wedge\_\allowbreak dim\_\allowbreak eq\_\allowbreak iff}, \texttt{three\_\allowbreak is\_\allowbreak unique}, \texttt{more\_\allowbreak wedges\_\allowbreak than\_\allowbreak directions}, \texttt{fewer\_\allowbreak wedges\_\allowbreak than\_\allowbreak directions}\par\smallskip
\noindent\textbf{\texttt{Axes.lean}} (9)\quad \texttt{rotation\_\allowbreak needs\_\allowbreak variation}, \texttt{combable\_\allowbreak no\_\allowbreak rotation}, \texttt{wound\_\allowbreak carries\_\allowbreak rotation}, \texttt{rotationless\_\allowbreak is\_\allowbreak trivial}, \texttt{rotation\_\allowbreak iff\_\allowbreak variation}, \texttt{no\_\allowbreak second\_\allowbreak axis}, \texttt{axis\_\allowbreak or\_\allowbreak variation}, \texttt{axis\_\allowbreak stabilizer\_\allowbreak nontrivial}, \texttt{multiplet\_\allowbreak flatness\_\allowbreak needs\_\allowbreak per\_\allowbreak direction\_\allowbreak scalings}\par\smallskip

\subsection*{第三部分\quad 分叉点}
\noindent\textbf{\texttt{Locus.lean}} (11)\quad \texttt{observability\_\allowbreak eq\_\allowbreak postulate}, \texttt{three\_\allowbreak from\_\allowbreak observability}, \texttt{observability\_\allowbreak fails\_\allowbreak elsewhere}, \texttt{observability\_\allowbreak iff\_\allowbreak three}, \texttt{no\_\allowbreak scale\_\allowbreak invariant\_\allowbreak interaction}, \texttt{constant\_\allowbreak coupling\_\allowbreak iff\_\allowbreak no\_\allowbreak content}, \texttt{isotropic\_\allowbreak pair\_\allowbreak parallel}, \texttt{isotropic\_\allowbreak no\_\allowbreak rotation}, \texttt{labels\_\allowbreak require\_\allowbreak anisotropy}, \texttt{everything\_\allowbreak switches\_\allowbreak on\_\allowbreak together}, \texttt{particle\_\allowbreak locus\_\allowbreak is\_\allowbreak anisotropic}\par\smallskip
\noindent\textbf{\texttt{Dynamics.lean}} (13)\quad \texttt{no\_\allowbreak flatness\_\allowbreak from\_\allowbreak commuting\_\allowbreak derivations}, \texttt{nonflat\_\allowbreak witness}, \texttt{commutative\_\allowbreak sector\_\allowbreak transports\_\allowbreak commute}, \texttt{commutative\_\allowbreak curvature\_\allowbreak undetectable}, \texttt{source\_\allowbreak conserved\_\allowbreak of\_\allowbreak field\_\allowbreak equation}, \texttt{no\_\allowbreak field\_\allowbreak equation\_\allowbreak of\_\allowbreak nonconserved}, \texttt{source\_\allowbreak antisymm}, \texttt{field\_\allowbreak equation\_\allowbreak content}, \texttt{same\_\allowbreak multiplet\_\allowbreak no\_\allowbreak rotation}, \texttt{rotation\_\allowbreak implies\_\allowbreak different\_\allowbreak multiplet}, \texttt{larger\_\allowbreak multiplet\_\allowbreak fewer\_\allowbreak generators}, \texttt{multiplets\_\allowbreak fixed\_\allowbreak by\_\allowbreak pattern}, \texttt{multiplet\_\allowbreak change\_\allowbreak needs\_\allowbreak pattern\_\allowbreak change}\par\smallskip

\subsection*{第四部分\quad 各向同性扇区}
\noindent\textbf{\texttt{Conformal.lean}} (20)\quad \texttt{hess\_\allowbreak symm}, \texttt{d\_\allowbreak sig\_\allowbreak comm}, \texttt{sum\_\allowbreak mul\_\allowbreak gradsq}, \texttt{sum\_\allowbreak const\_\allowbreak fin}, \texttt{scaleCurv\_\allowbreak eq\_\allowbreak zero}, \texttt{Chr\_\allowbreak symm}, \texttt{sum\_\allowbreak d\_\allowbreak Chr}, \texttt{sum\_\allowbreak Chr\_\allowbreak diag}, \texttt{sum\_\allowbreak d\_\allowbreak Chr\_\allowbreak trace}, \texttt{sum\_\allowbreak ChrTrace\_\allowbreak Chr}, \texttt{sum\_\allowbreak Chr\_\allowbreak Chr\_\allowbreak full}, \texttt{sum\_\allowbreak Chr\_\allowbreak Chr}, \texttt{Ric\_\allowbreak eq}, \texttt{RscBare\_\allowbreak eq}, \texttt{Defm\_\allowbreak symm}, \texttt{d\_\allowbreak Chr\_\allowbreak diff}, \texttt{two\_\allowbreak Rm\_\allowbreak eq}, \texttt{Rm\_\allowbreak eq\_\allowbreak zero\_\allowbreak of\_\allowbreak Defm\_\allowbreak eq\_\allowbreak zero}, \texttt{Ric\_\allowbreak eq\_\allowbreak zero\_\allowbreak of\_\allowbreak Defm\_\allowbreak eq\_\allowbreak zero}, \texttt{two\_\allowbreak Ric\_\allowbreak eq}\par\smallskip
\noindent\textbf{\texttt{Frame.lean}} (5)\quad \texttt{isotropic\_\allowbreak metric}, \texttt{reciprocal\_\allowbreak metric}, \texttt{isotropic\_\allowbreak reciprocal\_\allowbreak forces\_\allowbreak flat}, \texttt{isotropic\_\allowbreak reciprocal\_\allowbreak metric\_\allowbreak bare}, \texttt{exists\_\allowbreak reciprocal\_\allowbreak nonisotropic}\par\smallskip
\noindent\textbf{\texttt{Unify.lean}} (9)\quad \texttt{toMetric\_\allowbreak of\_\allowbreak isotropic}, \texttt{ofScalar\_\allowbreak isotropic}, \texttt{ofScalar\_\allowbreak toMetric}, \texttt{isotropic\_\allowbreak iff\_\allowbreak ofScalar}, \texttt{toDirScale\_\allowbreak isotropic}, \texttt{toDirScale\_\allowbreak toMetric}, \texttt{sector\_\allowbreak coherence}, \texttt{isotropic\_\allowbreak ratio\_\allowbreak const}, \texttt{not\_\allowbreak isotropic\_\allowbreak of\_\allowbreak ne}\par\smallskip
\noindent\textbf{\texttt{Newton.lean}} (7)\quad \texttt{lap\_\allowbreak newtPot}, \texttt{RscBare\_\allowbreak linear}, \texttt{poisson}, \texttt{poisson\_\allowbreak three}, \texttt{sig\_\allowbreak inr}, \texttt{gradsq\_\allowbreak inr}, \texttt{poisson\_\allowbreak exact}\par\smallskip
\noindent\textbf{\texttt{Schwarzschild.lean}} (5)\quad \texttt{ricci\_\allowbreak combination}, \texttt{vacuum\_\allowbreak log\_\allowbreak derivative}, \texttt{scale\_\allowbreak sum\_\allowbreak derivative\_\allowbreak zero}, \texttt{vacuum\_\allowbreak iff\_\allowbreak scale\_\allowbreak sum\_\allowbreak stationary}, \texttt{area\_\allowbreak element\_\allowbreak stationary}\par\smallskip
\noindent\textbf{\texttt{RicciDiag.lean}} (4)\quad \texttt{ricciTT\_\allowbreak eq}, \texttt{ricciRR\_\allowbreak eq}, \texttt{gravity\_\allowbreak chain}, \texttt{area\_\allowbreak stationary\_\allowbreak derived}\par\smallskip
\noindent\textbf{\texttt{Vacuum.lean}} (5)\quad \texttt{const\_\allowbreak of\_\allowbreak deriv\_\allowbreak zero\_\allowbreak and\_\allowbreak vanishing}, \texttt{reciprocity\_\allowbreak of\_\allowbreak asymptotic\_\allowbreak flatness}, \texttt{scale\_\allowbreak product\_\allowbreak eq\_\allowbreak one}, \texttt{product\_\allowbreak constant\_\allowbreak without\_\allowbreak flatness}, \texttt{tr\_\allowbreak area\_\allowbreak element\_\allowbreak constant}\par\smallskip
\noindent\textbf{\texttt{Deflection.lean}} (7)\quad \texttt{iso\_\allowbreak is\_\allowbreak exactly\_\allowbreak half}, \texttt{iso\_\allowbreak ratio\_\allowbreak half}, \texttt{missing\_\allowbreak half}, \texttt{solarObs\_\allowbreak value}, \texttt{solarIso\_\allowbreak value}, \texttt{solar\_\allowbreak ratio}, \texttt{iso\_\allowbreak fails\_\allowbreak by\_\allowbreak far}\par\smallskip
\noindent\textbf{\texttt{PPN.lean}} (9)\quad \texttt{gamma\_\allowbreak zero\_\allowbreak is\_\allowbreak isotropic}, \texttt{gamma\_\allowbreak one\_\allowbreak is\_\allowbreak observed}, \texttt{reciprocal\_\allowbreak expansion}, \texttt{gamma\_\allowbreak eq\_\allowbreak one\_\allowbreak of\_\allowbreak reciprocal}, \texttt{gamma\_\allowbreak eq\_\allowbreak one\_\allowbreak first\_\allowbreak order}, \texttt{deflection\_\allowbreak reciprocal\_\allowbreak eq\_\allowbreak observed}, \texttt{halves\_\allowbreak equal}, \texttt{solar\_\allowbreak reciprocal\_\allowbreak value}, \texttt{gamma\_\allowbreak prediction\_\allowbreak is\_\allowbreak sharp}\par\smallskip
\noindent\textbf{\texttt{Precession.lean}} (8)\quad \texttt{coeff\_\allowbreak reciprocal}, \texttt{coeff\_\allowbreak isotropic}, \texttt{isotropic\_\allowbreak is\_\allowbreak third}, \texttt{coeff\_\allowbreak affine\_\allowbreak in\_\allowbreak gamma}, \texttt{mercuryFactor\_\allowbreak bounds}, \texttt{mercury\_\allowbreak value}, \texttt{mercury\_\allowbreak isotropic\_\allowbreak value}, \texttt{two\_\allowbreak tests\_\allowbreak different\_\allowbreak factors}\par\smallskip
\noindent\textbf{\texttt{Covariance.lean}} (7)\quad \texttt{sig\_\allowbreak eq}, \texttt{Chr\_\allowbreak eq}, \texttt{Rm\_\allowbreak eq}, \texttt{Ric\_\allowbreak eq\_\allowbreak of\_\allowbreak drift}, \texttt{Defm\_\allowbreak eq}, \texttt{field\_\allowbreak equation\_\allowbreak form\_\allowbreak invariant}, \texttt{no\_\allowbreak local\_\allowbreak expansion}\par\smallskip
\noindent\textbf{\texttt{Singularity.lean}} (6)\quad \texttt{scale\_\allowbreak ne\_\allowbreak zero}, \texttt{energy\_\allowbreak ne\_\allowbreak zero}, \texttt{conformal\_\allowbreak factor\_\allowbreak isUnit}, \texttt{conformal\_\allowbreak factor\_\allowbreak invertible}, \texttt{curvature\_\allowbreak indep\_\allowbreak of\_\allowbreak scale\_\allowbreak value}, \texttt{curvature\_\allowbreak zero\_\allowbreak of\_\allowbreak uniform\_\allowbreak scale}\par\smallskip
\noindent\textbf{\texttt{Horizon.lean}} (15)\quad \texttt{gtt\_\allowbreak ne\_\allowbreak zero}, \texttt{reciprocal\_\allowbreak never\_\allowbreak degenerate}, \texttt{redshift\_\allowbreak never\_\allowbreak infinite}, \texttt{no\_\allowbreak degeneracy\_\allowbreak anywhere}, \texttt{null\_\allowbreak cone\_\allowbreak scale\_\allowbreak independent}, \texttt{waves\_\allowbreak travel\_\allowbreak on\_\allowbreak the\_\allowbreak light\_\allowbreak cone}, \texttt{no\_\allowbreak energy\_\allowbreak dependent\_\allowbreak speed}, \texttt{two\_\allowbreak polarizations}, \texttt{no\_\allowbreak breathing\_\allowbreak mode}, \texttt{polarization\_\allowbreak count}, \texttt{entropy\_\allowbreak of\_\allowbreak scale\_\allowbreak ratio}, \texttt{entropy\_\allowbreak doubling\_\allowbreak is\_\allowbreak additive}, \texttt{entropy\_\allowbreak increment\_\allowbreak independent\_\allowbreak of\_\allowbreak size}, \texttt{no\_\allowbreak minimum\_\allowbreak length}, \texttt{no\_\allowbreak final\_\allowbreak description}\par\smallskip
\noindent\textbf{\texttt{Waves.lean}} (32)\quad \texttt{conserved\_\allowbreak has\_\allowbreak no\_\allowbreak dynamics}, \texttt{monopole\_\allowbreak does\_\allowbreak not\_\allowbreak radiate}, \texttt{dipole\_\allowbreak does\_\allowbreak not\_\allowbreak radiate}, \texttt{quadrupole\_\allowbreak is\_\allowbreak leading}, \texttt{power\_\allowbreak pos}, \texttt{power\_\allowbreak ratio\_\allowbreak halving}, \texttt{inspiral\_\allowbreak shrinks\_\allowbreak orbit}, \texttt{separation\_\allowbreak rate}, \texttt{chirpMass\_\allowbreak symm}, \texttt{chirpMass\_\allowbreak equal}, \texttt{chirp\_\allowbreak exponent}, \texttt{chirp\_\allowbreak mass\_\allowbreak exponent}, \texttt{hulse\_\allowbreak taylor\_\allowbreak agreement}, \texttt{dipole\_\allowbreak bounded\_\allowbreak by\_\allowbreak hulse\_\allowbreak taylor}, \texttt{gw250114\_\allowbreak radiated\_\allowbreak fraction}, \texttt{area\_\allowbreak law\_\allowbreak confirmed}, \texttt{gw150914\_\allowbreak equal\_\allowbreak mass\_\allowbreak component}, \texttt{dipole\_\allowbreak bounded\_\allowbreak by\_\allowbreak double\_\allowbreak pulsar}, \texttt{double\_\allowbreak pulsar\_\allowbreak tighter}, \texttt{reflectivity\_\allowbreak nonzero\_\allowbreak but\_\allowbreak unbounded}, \texttt{echo\_\allowbreak delay\_\allowbreak logarithmic}, \texttt{delay\_\allowbreak insensitive\_\allowbreak to\_\allowbreak reflectivity}, \texttt{echo\_\allowbreak train\_\allowbreak is\_\allowbreak a\_\allowbreak comb}, \texttt{echo\_\allowbreak amplitudes\_\allowbreak geometric}, \texttt{null\_\allowbreak search\_\allowbreak bounds\_\allowbreak reflectivity}, \texttt{damping\_\allowbreak is\_\allowbreak the\_\allowbreak looser\_\allowbreak constraint}, \texttt{scale\_\allowbreak ratio\_\allowbreak never\_\allowbreak decreases}, \texttt{log\_\allowbreak law\_\allowbreak is\_\allowbreak not\_\allowbreak area\_\allowbreak law}, \texttt{gw241011\_\allowbreak spin\_\allowbreak is\_\allowbreak rapid}, \texttt{spin\_\allowbreak bound\_\allowbreak beats\_\allowbreak echo\_\allowbreak bound}, \texttt{no\_\allowbreak ultralight\_\allowbreak scalar\_\allowbreak available}, \texttt{same\_\allowbreak constant\_\allowbreak both\_\allowbreak sectors}\par\smallskip

\subsection*{第四部分附\quad 结构的见证}
\noindent\textbf{\texttt{Witness.lean}} (6)\quad \texttt{shift\_\allowbreak is\_\allowbreak exact}, \texttt{shift\_\allowbreak nontrivial}, \texttt{shift\_\allowbreak commutes\_\allowbreak with\_\allowbreak derivation}, \texttt{linearFlow\_\allowbreak moves}, \texttt{constMoment\_\allowbreak conserved}, \texttt{scaleShift\_\allowbreak is\_\allowbreak witnessed}\par\smallskip

\subsection*{第五部分\quad 各向异性扇区}
\noindent\textbf{\texttt{Gauge.lean}} (5)\quad \texttt{mem\_\allowbreak stabilizer\_\allowbreak iff}, \texttt{stabilizer\_\allowbreak eq\_\allowbreak bot\_\allowbreak of\_\allowbreak injective}, \texttt{stabilizer\_\allowbreak eq\_\allowbreak top\_\allowbreak of\_\allowbreak isotropic}, \texttt{sameOrbit\_\allowbreak iff\_\allowbreak eq\_\allowbreak scale}, \texttt{internal\_\allowbreak symmetry\_\allowbreak determined}\par\smallskip
\noindent\textbf{\texttt{Transport.lean}} (6)\quad \texttt{transports\_\allowbreak self\_\allowbreak iff}, \texttt{transport\_\allowbreak unique\_\allowbreak of\_\allowbreak nondegenerate}, \texttt{transport\_\allowbreak coset}, \texttt{transports\_\allowbreak of\_\allowbreak stabilizer}, \texttt{transport\_\allowbreak unique\_\allowbreak iff\_\allowbreak stabilizer\_\allowbreak trivial}, \texttt{nondegenerate\_\allowbreak no\_\allowbreak gauge\_\allowbreak freedom}\par\smallskip
\noindent\textbf{\texttt{Axis.lean}} (9)\quad \texttt{bareForm\_\allowbreak neg}, \texttt{physForm\_\allowbreak neg}, \texttt{neg\_\allowbreak indistinguishable}, \texttt{physForm\_\allowbreak indistinguishable}, \texttt{isotropic\_\allowbreak no\_\allowbreak order\_\allowbreak parameter}, \texttt{nondegenerate\_\allowbreak full\_\allowbreak frame}, \texttt{uniaxial\_\allowbreak stabilizer}, \texttt{uniaxial\_\allowbreak axis\_\allowbreak distinguished}, \texttt{scalar\_\allowbreak picture\_\allowbreak is\_\allowbreak the\_\allowbreak isotropic\_\allowbreak locus}\par\smallskip
\noindent\textbf{\texttt{Defect.lean}} (13)\quad \texttt{winding\_\allowbreak unique}, \texttt{vacuum\_\allowbreak winding}, \texttt{combine\_\allowbreak winding}, \texttt{anti\_\allowbreak winding}, \texttt{pair\_\allowbreak winding\_\allowbreak zero}, \texttt{vacuum\_\allowbreak isTrivial}, \texttt{stable\_\allowbreak of\_\allowbreak winding\_\allowbreak ne\_\allowbreak zero}, \texttt{winding\_\allowbreak eq\_\allowbreak zero\_\allowbreak of\_\allowbreak trivial}, \texttt{mass\_\allowbreak pos}, \texttt{antiparticle\_\allowbreak same\_\allowbreak mass}, \texttt{mass\_\allowbreak zero}, \texttt{mass\_\allowbreak ratio\_\allowbreak geometric}, \texttt{mass\_\allowbreak ratio\_\allowbreak const}\par\smallskip
\noindent\textbf{\texttt{Charges.lean}} (15)\quad \texttt{comp\_\allowbreak block}, \texttt{comp\_\allowbreak hedgehog}, \texttt{anti\_\allowbreak block}, \texttt{anti\_\allowbreak hedgehog}, \texttt{triv\_\allowbreak observable}, \texttt{block\_\allowbreak confined}, \texttt{hedgehog\_\allowbreak free}, \texttt{hedgehog\_\allowbreak surjective\_\allowbreak on\_\allowbreak observables}, \texttt{exactly\_\allowbreak one\_\allowbreak confined}, \texttt{hedgehog\_\allowbreak not\_\allowbreak determined\_\allowbreak by\_\allowbreak block}, \texttt{block\_\allowbreak not\_\allowbreak determined\_\allowbreak by\_\allowbreak hedgehog}, \texttt{charges\_\allowbreak independent}, \texttt{pair\_\allowbreak observable}, \texttt{pair\_\allowbreak hedgehog\_\allowbreak zero}, \texttt{observable\_\allowbreak pair\_\allowbreak any\_\allowbreak hedgehog}\par\smallskip
\noindent\textbf{\texttt{Particle.lean}} (16)\quad \texttt{boost\_\allowbreak injective}, \texttt{pattern\_\allowbreak is\_\allowbreak one\_\allowbreak vector}, \texttt{mixture\_\allowbreak second\_\allowbreak moment}, \texttt{mixture\_\allowbreak cvSq\_\allowbreak ge\_\allowbreak one}, \texttt{mixture\_\allowbreak cvSq\_\allowbreak eq\_\allowbreak one\_\allowbreak iff}, \texttt{observed\_\allowbreak disfavours\_\allowbreak higher\_\allowbreak rank}, \texttt{rate\_\allowbreak value\_\allowbreak not\_\allowbreak fixed}, \texttt{block\_\allowbreak two\_\allowbreak valued}, \texttt{block\_\allowbreak cannot\_\allowbreak be\_\allowbreak three\_\allowbreak valued}, \texttt{block\_\allowbreak self\_\allowbreak inverse}, \texttt{bind\_\allowbreak hedgehog}, \texttt{bind\_\allowbreak block}, \texttt{two\_\allowbreak odd\_\allowbreak make\_\allowbreak even}, \texttt{hedgehog\_\allowbreak conserved\_\allowbreak in\_\allowbreak splitting}, \texttt{no\_\allowbreak further\_\allowbreak label}, \texttt{particle\_\allowbreak label\_\allowbreak structure}\par\smallskip
\noindent\textbf{\texttt{Emergence.lean}} (9)\quad \texttt{mem\_\allowbreak resolved}, \texttt{resolved\_\allowbreak mono}, \texttt{content\_\allowbreak never\_\allowbreak complete}, \texttt{always\_\allowbreak more\_\allowbreak above}, \texttt{no\_\allowbreak fundamental\_\allowbreak list}, \texttt{mass\_\allowbreak iff\_\allowbreak threshold}, \texttt{mass\_\allowbreak mono}, \texttt{content\_\allowbreak eq\_\allowbreak of\_\allowbreak no\_\allowbreak threshold}, \texttt{content\_\allowbreak changes\_\allowbreak only\_\allowbreak at\_\allowbreak threshold}\par\smallskip

\subsection*{第六部分\quad 量子扇区}
\noindent\textbf{\texttt{Deformation.lean}} (12)\quad \texttt{poisson\_\allowbreak antisymm}, \texttt{poisson\_\allowbreak self}, \texttt{poisson\_\allowbreak leibniz\_\allowbreak left}, \texttt{poisson\_\allowbreak leibniz\_\allowbreak right}, \texttt{star\_\allowbreak assoc}, \texttt{star\_\allowbreak classical}, \texttt{star\_\allowbreak comm\_\allowbreak classical}, \texttt{star\_\allowbreak commutator}, \texttt{star\_\allowbreak comm\_\allowbreak iff\_\allowbreak poisson\_\allowbreak zero}, \texttt{canonical\_\allowbreak poisson}, \texttt{ccr\_\allowbreak of\_\allowbreak canonical}, \texttt{canonical\_\allowbreak poisson\_\allowbreak self}\par\smallskip
\noindent\textbf{\texttt{Crossed.lean}} (11)\quad \texttt{map\_\allowbreak zero}, \texttt{shiftOp\_\allowbreak apply}, \texttt{mulOp\_\allowbreak apply}, \texttt{shift\_\allowbreak mul\_\allowbreak eq}, \texttt{commutator\_\allowbreak eq}, \texttt{comm\_\allowbreak iff\_\allowbreak shift\_\allowbreak trivial}, \texttt{delta\_\allowbreak twisted\_\allowbreak leibniz}, \texttt{delta\_\allowbreak leibniz\_\allowbreak of\_\allowbreak untwisted}, \texttt{delta\_\allowbreak eq\_\allowbreak zero\_\allowbreak iff}, \texttt{ScaleShift}, \texttt{shift\_\allowbreak mul\_\allowbreak eq\_\allowbreak comp}\par\smallskip
\noindent\textbf{\texttt{NCConformal.lean}} (18)\quad \texttt{kr\_\allowbreak symm}, \texttt{kr\_\allowbreak comm}, \texttt{kr\_\allowbreak self}, \texttt{sum\_\allowbreak kr\_\allowbreak left}, \texttt{sum\_\allowbreak kr\_\allowbreak right}, \texttt{kr\_\allowbreak swap}, \texttt{scalar\_\allowbreak mul\_\allowbreak scalar}, \texttt{hess\_\allowbreak symm}, \texttt{Chr\_\allowbreak symm}, \texttt{Defm\_\allowbreak antisymm\_\allowbreak part}, \texttt{Defm\_\allowbreak symm\_\allowbreak iff\_\allowbreak commutator}, \texttt{Defm\_\allowbreak symm\_\allowbreak iff\_\allowbreak commute}, \texttt{symmetric\_\allowbreak part\_\allowbreak uncorrected}, \texttt{correction\_\allowbreak is\_\allowbreak pure\_\allowbreak antisymmetric}, \texttt{sum\_\allowbreak Chr\_\allowbreak Chr\_\allowbreak full\_\allowbreak nc}, \texttt{sum\_\allowbreak Chr\_\allowbreak Chr\_\allowbreak full\_\allowbreak classical}, \texttt{quantum\_\allowbreak correction\_\allowbreak couples\_\allowbreak to\_\allowbreak rotation}, \texttt{correction\_\allowbreak vanishes\_\allowbreak on\_\allowbreak commuting}\par\smallskip

\subsection*{第七部分\quad 标度流与内容}
\noindent\textbf{\texttt{RG.lean}} (13)\quad \texttt{observable\_\allowbreak const\_\allowbreak at\_\allowbreak fixedPoint}, \texttt{fixedPoint\_\allowbreak invariant}, \texttt{callan\_\allowbreak symanzik}, \texttt{scale\_\allowbreak invariant\_\allowbreak of\_\allowbreak fixedPoint}, \texttt{mem\_\allowbreak active}, \texttt{active\_\allowbreak mono}, \texttt{active\_\allowbreak eq\_\allowbreak of\_\allowbreak no\_\allowbreak threshold}, \texttt{active\_\allowbreak finite\_\allowbreak range}, \texttt{active\_\allowbreak eq\_\allowbreak univ\_\allowbreak of\_\allowbreak ge}, \texttt{active\_\allowbreak eq\_\allowbreak empty\_\allowbreak of\_\allowbreak lt}, \texttt{beta\_\allowbreak const\_\allowbreak between\_\allowbreak thresholds}, \texttt{beta\_\allowbreak finite\_\allowbreak range}, \texttt{beta\_\allowbreak eq\_\allowbreak of\_\allowbreak all\_\allowbreak active}\par\smallskip
\noindent\textbf{\texttt{Running.lean}} (6)\quad \texttt{resolvedCount\_\allowbreak self}, \texttt{invSqCoupling\_\allowbreak at\_\allowbreak base}, \texttt{hasDerivAt\_\allowbreak invSqCoupling}, \texttt{slope\_\allowbreak eq\_\allowbreak two\_\allowbreak b}, \texttt{invSqCoupling\_\allowbreak eq\_\allowbreak one\_\allowbreak loop}, \texttt{asympt\_\allowbreak free\_\allowbreak iff\_\allowbreak pos\_\allowbreak density}\par\smallskip
\noindent\textbf{\texttt{QCD.lean}} (10)\quad \texttt{running\_\allowbreak inv\_\allowbreak sq}, \texttt{asymptotic\_\allowbreak freedom}, \texttt{lambda\_\allowbreak invariant}, \texttt{lambda\_\allowbreak pos}, \texttt{hadron\_\allowbreak mass\_\allowbreak ratio}, \texttt{spectrum\_\allowbreak rigid}, \texttt{geometric\_\allowbreak tower\_\allowbreak fails}, \texttt{ratio₁\_\allowbreak value}, \texttt{ratio₂\_\allowbreak value}, \texttt{lightestGlueballInLambda\_\allowbreak value}\par\smallskip
\noindent\textbf{\texttt{Spectrum.lean}} (9)\quad \texttt{mass\_\allowbreak threshold\_\allowbreak inverse}, \texttt{equal\_\allowbreak gaps\_\allowbreak variance\_\allowbreak zero}, \texttt{observed\_\allowbreak cvSq\_\allowbreak value}, \texttt{meanGap\_\allowbreak value}, \texttt{observed\_\allowbreak not\_\allowbreak geometric}, \texttt{observed\_\allowbreak within\_\allowbreak scale\_\allowbreak invariant\_\allowbreak range}, \texttt{densityFromGaps\_\allowbreak value}, \texttt{deltaN\_\allowbreak inverse}, \texttt{density\_\allowbreak gap\_\allowbreak inverse}\par\smallskip
\noindent\textbf{\texttt{Entropy.lean}} (9)\quad \texttt{mem\_\allowbreak sat}, \texttt{subset\_\allowbreak sat}, \texttt{card\_\allowbreak le\_\allowbreak card\_\allowbreak sat}, \texttt{card\_\allowbreak image\_\allowbreak perm}, \texttt{entropy\_\allowbreak nondecreasing}, \texttt{no\_\allowbreak reverse\_\allowbreak law}, \texttt{evolve\_\allowbreak no\_\allowbreak reverse\_\allowbreak law}, \texttt{charge\_\allowbreak change\_\allowbreak implies\_\allowbreak nonequilibrium}, \texttt{entropy\_\allowbreak production\_\allowbreak implies\_\allowbreak nonequilibrium}\par\smallskip
\noindent\textbf{\texttt{Light.lean}} (6)\quad \texttt{unit\_\allowbreak sq\_\allowbreak mul\_\allowbreak eq\_\allowbreak zero\_\allowbreak iff}, \texttt{isNull\_\allowbreak iff\_\allowbreak bare}, \texttt{isNull\_\allowbreak scale\_\allowbreak invariant}, \texttt{scale\_\allowbreak determined\_\allowbreak of\_\allowbreak nonnull}, \texttt{physForm\_\allowbreak ne\_\allowbreak bareForm}, \texttt{null\_\allowbreak weight\_\allowbreak zero}\par\smallskip
\noindent\textbf{\texttt{Observation.lean}} (6)\quad \texttt{null\_\allowbreak has\_\allowbreak energy}, \texttt{redshift\_\allowbreak is\_\allowbreak scale\_\allowbreak ratio}, \texttt{scale\_\allowbreak differs\_\allowbreak of\_\allowbreak energy\_\allowbreak differs}, \texttt{obsEnergy\_\allowbreak rescale}, \texttt{redshift\_\allowbreak fiducial\_\allowbreak invariant}, \texttt{causal\_\allowbreak structure\_\allowbreak still\_\allowbreak blind}\par\smallskip
\noindent\textbf{\texttt{Particles.lean}} (9)\quad \texttt{energy\_\allowbreak pos}, \texttt{energy\_\allowbreak ge}, \texttt{one\_\allowbreak le\_\allowbreak scaleRatio}, \texttt{scaleRatio\_\allowbreak mono}, \texttt{lifetime\_\allowbreak ge}, \texttt{lifetime\_\allowbreak mono}, \texttt{lifetime\_\allowbreak at\_\allowbreak rest}, \texttt{available\_\allowbreak mono}, \texttt{not\_\allowbreak available\_\allowbreak of\_\allowbreak lt}\par\smallskip
\noindent\textbf{\texttt{DarkMatter.lean}} (4)\quad \texttt{detResp\_\allowbreak pos}, \texttt{detection\_\allowbreak failure}, \texttt{dark\_\allowbreak matter\_\allowbreak exists}, \texttt{response\_\allowbreak vanishes}\par\smallskip
\noindent\textbf{\texttt{DarkEnergy.lean}} (8)\quad \texttt{hasDerivAt\_\allowbreak exp\_\allowbreak mul}, \texttt{dissipation\_\allowbreak solution}, \texttt{energy\_\allowbreak strictly\_\allowbreak decreasing}, \texttt{scaleFactor\_\allowbreak eq}, \texttt{hasDerivAt\_\allowbreak scaleFactor}, \texttt{hasDerivAt\_\allowbreak hubble}, \texttt{acceleration\_\allowbreak pos}, \texttt{open\_\allowbreak universe\_\allowbreak accelerates}\par\smallskip

\subsection*{第八部分\quad 预言与数据}
\noindent\textbf{\texttt{Predictions.lean}} (17)\quad \texttt{accel\_\allowbreak numerator}, \texttt{powerlaw\_\allowbreak accel\_\allowbreak numerator}, \texttt{accelerates\_\allowbreak iff}, \texttt{eos\_\allowbreak dissipationExponent}, \texttt{dissipationExponent\_\allowbreak eos}, \texttt{eos\_\allowbreak one}, \texttt{eos\_\allowbreak eq\_\allowbreak neg\_\allowbreak one\_\allowbreak iff}, \texttt{eos\_\allowbreak lt\_\allowbreak third\_\allowbreak iff}, \texttt{eos\_\allowbreak lt\_\allowbreak neg\_\allowbreak one\_\allowbreak iff}, \texttt{worked\_\allowbreak inversion}, \texttt{muon\_\allowbreak lifetime\_\allowbreak value}, \texttt{muon\_\allowbreak matches\_\allowbreak experiment}, \texttt{poundRebka\_\allowbreak value}, \texttt{darkEnergyDensity\_\allowbreak value}, \texttt{logScale\_\allowbreak step}, \texttt{logScale\_\allowbreak ratio\_\allowbreak const}, \texttt{logScale\_\allowbreak closed}\par\smallskip
\noindent\textbf{\texttt{Cosmos.lean}} (13)\quad \texttt{shift\_\allowbreak invariant\_\allowbreak is\_\allowbreak constant}, \texttt{rate\_\allowbreak constant\_\allowbreak of\_\allowbreak fiducial\_\allowbreak invariance}, \texttt{w\_\allowbreak does\_\allowbreak not\_\allowbreak evolve}, \texttt{scale\_\allowbreak ratio\_\allowbreak depends\_\allowbreak only\_\allowbreak on\_\allowbreak difference}, \texttt{no\_\allowbreak first\_\allowbreak moment}, \texttt{finite\_\allowbreak lookback\_\allowbreak no\_\allowbreak origin}, \texttt{isotropic\_\allowbreak gravitates\_\allowbreak without\_\allowbreak labels}, \texttt{dark\_\allowbreak matter\_\allowbreak has\_\allowbreak no\_\allowbreak species}, \texttt{two\_\allowbreak dark\_\allowbreak routes\_\allowbreak differ}, \texttt{resolvable\_\allowbreak needs\_\allowbreak width\_\allowbreak below\_\allowbreak gap}, \texttt{mean\_\allowbreak gap\_\allowbreak bound}, \texttt{broad\_\allowbreak structures\_\allowbreak unresolvable}, \texttt{resolvability\_\allowbreak boundary}\par\smallskip
\noindent\textbf{\texttt{Native.lean}} (15)\quad \texttt{not\_\allowbreak two\_\allowbreak places\_\allowbreak of\_\allowbreak parallel}, \texttt{separation\_\allowbreak is\_\allowbreak oriented}, \texttt{single\_\allowbreak anisotropy\_\allowbreak boundary}, \texttt{label\_\allowbreak forces\_\allowbreak anisotropy}, \texttt{zero\_\allowbreak window\_\allowbreak resolves\_\allowbreak nothing}, \texttt{content\_\allowbreak is\_\allowbreak interval\_\allowbreak valued}, \texttt{no\_\allowbreak global\_\allowbreak field\_\allowbreak equation}, \texttt{ordering\_\allowbreak discrepancy\_\allowbreak is\_\allowbreak the\_\allowbreak observable}, \texttt{parity\_\allowbreak conserved\_\allowbreak under\_\allowbreak binding}, \texttt{both\_\allowbreak labels\_\allowbreak conserved}, \texttt{entropy\_\allowbreak order\_\allowbreak is\_\allowbreak scale\_\allowbreak order}, \texttt{entropy\_\allowbreak strict\_\allowbreak iff\_\allowbreak scale\_\allowbreak strict}, \texttt{arrows\_\allowbreak are\_\allowbreak one\_\allowbreak order}, \texttt{no\_\allowbreak maximum\_\allowbreak anisotropy}, \texttt{no\_\allowbreak cutoff\_\allowbreak either\_\allowbreak end}\par\smallskip
\noindent\textbf{\texttt{Data.lean}} (19)\quad \texttt{desi\_\allowbreak tension\_\allowbreak exceeds\_\allowbreak three\_\allowbreak sigma}, \texttt{desi\_\allowbreak significance\_\allowbreak dataset\_\allowbreak dependent}, \texttt{desi\_\allowbreak combinations\_\allowbreak disagree\_\allowbreak internally}, \texttt{falsification\_\allowbreak threshold}, \texttt{ccl\_\allowbreak above\_\allowbreak bound}, \texttt{gkpy\_\allowbreak below\_\allowbreak bound}, \texttt{determinations\_\allowbreak straddle\_\allowbreak bound}, \texttt{gkpy\_\allowbreak updated\_\allowbreak inside}, \texttt{kappa\_\allowbreak below\_\allowbreak bound}, \texttt{rho\_\allowbreak below\_\allowbreak bound}, \texttt{delta\_\allowbreak below\_\allowbreak bound}, \texttt{z\_\allowbreak boson\_\allowbreak far\_\allowbreak inside}, \texttt{bound\_\allowbreak separates\_\allowbreak kappa\_\allowbreak from\_\allowbreak sigma}, \texttt{breit\_\allowbreak wigner\_\allowbreak spans\_\allowbreak the\_\allowbreak bound}, \texttt{dispersion\_\allowbreak limit\_\allowbreak above\_\allowbreak planck}, \texttt{framework\_\allowbreak predicts\_\allowbreak no\_\allowbreak finite\_\allowbreak scale}, \texttt{one\_\allowbreak charge\_\allowbreak prediction}, \texttt{zero\_\allowbreak nu\_\allowbreak beta\_\allowbreak beta\_\allowbreak decides}, \texttt{scorecard}\par\smallskip

\subsection*{第九部分\quad 登记册}
\noindent\textbf{\texttt{Color.lean}} (17)\quad \texttt{comp\_\allowbreak perm}, \texttt{anti\_\allowbreak perm}, \texttt{triv\_\allowbreak perm}, \texttt{charge\_\allowbreak comp}, \texttt{charge\_\allowbreak anti}, \texttt{triv\_\allowbreak observable}, \texttt{confinement}, \texttt{pair\_\allowbreak observable}, \texttt{pair\_\allowbreak charge}, \texttt{rep\_\allowbreak perm}, \texttt{rep\_\allowbreak charge}, \texttt{bound\_\allowbreak state\_\allowbreak observable}, \texttt{bound\_\allowbreak state\_\allowbreak size}, \texttt{no\_\allowbreak partial\_\allowbreak state}, \texttt{constituent\_\allowbreak charge\_\allowbreak fraction}, \texttt{constituent\_\allowbreak charge\_\allowbreak rat}, \texttt{bound\_\allowbreak charge\_\allowbreak divisible}\par\smallskip
\noindent\textbf{\texttt{ColorAudit.lean}} (8)\quad \texttt{identification\_\allowbreak predicts\_\allowbreak two}, \texttt{identification\_\allowbreak predicts\_\allowbreak halves}, \texttt{identification\_\allowbreak fails\_\allowbreak baryon}, \texttt{identification\_\allowbreak fails\_\allowbreak charge}, \texttt{spatial\_\allowbreak block\_\allowbreak is\_\allowbreak not\_\allowbreak colour\_\allowbreak block}, \texttt{three\_\allowbreak is\_\allowbreak forced\_\allowbreak by\_\allowbreak data}, \texttt{three\_\allowbreak matches\_\allowbreak charge}, \texttt{geometry\_\allowbreak offers\_\allowbreak two\_\allowbreak data\_\allowbreak needs\_\allowbreak three}\par\smallskip
\noindent\textbf{\texttt{Codimension.lean}} (12)\quad \texttt{defectDim\_\allowbreak two}, \texttt{defectDim\_\allowbreak three}, \texttt{defectDim\_\allowbreak four}, \texttt{pointlike\_\allowbreak iff\_\allowbreak two\_\allowbreak dim}, \texttt{three\_\allowbreak dim\_\allowbreak gives\_\allowbreak strings}, \texttt{pointlike\_\allowbreak forces\_\allowbreak three}, \texttt{our\_\allowbreak world\_\allowbreak gives\_\allowbreak strings}, \texttt{no\_\allowbreak pointlike\_\allowbreak in\_\allowbreak four}, \texttt{string\_\allowbreak mass\_\allowbreak not\_\allowbreak topological}, \texttt{stringMass\_\allowbreak strictMono}, \texttt{stringMass\_\allowbreak zero}, \texttt{extended\_\allowbreak or\_\allowbreak low\_\allowbreak dimensional}\par\smallskip
\noindent\textbf{\texttt{Expressive.lean}} (10)\quad \texttt{diffPattern\_\allowbreak shift}, \texttt{invariant\_\allowbreak factors}, \texttt{diffPattern\_\allowbreak invariant}, \texttt{invariant\_\allowbreak iff\_\allowbreak diffPattern}, \texttt{eval\_\allowbreak not\_\allowbreak invariant}, \texttt{scale\_\allowbreak says\_\allowbreak nothing\_\allowbreak about\_\allowbreak rotation}, \texttt{rotation\_\allowbreak says\_\allowbreak nothing\_\allowbreak about\_\allowbreak scale}, \texttt{channels\_\allowbreak independent}, \texttt{neither\_\allowbreak channel\_\allowbreak complete}, \texttt{no\_\allowbreak channel\_\allowbreak coupling}\par\smallskip
\noindent\textbf{\texttt{CrossCheck.lean}} (9)\quad \texttt{rhoAll\_\allowbreak value}, \texttt{rhoQuark\_\allowbreak value}, \texttt{densities\_\allowbreak differ}, \texttt{density\_\allowbreak ratio\_\allowbreak near\_\allowbreak two}, \texttt{sector\_\allowbreak relation\_\allowbreak is\_\allowbreak identity}, \texttt{qcd\_\allowbreak extraction\_\allowbreak returns\_\allowbreak input}, \texttt{no\_\allowbreak second\_\allowbreak determination}, \texttt{quarkCvSq\_\allowbreak value}, \texttt{quark\_\allowbreak not\_\allowbreak geometric}\par\smallskip
\noindent\textbf{\texttt{MassAudit.lean}} (9)\quad \texttt{massExp\_\allowbreak neg}, \texttt{massQuad\_\allowbreak neg}, \texttt{massExp\_\allowbreak pos}, \texttt{massQuad\_\allowbreak pos}, \texttt{massExp\_\allowbreak zero}, \texttt{massQuad\_\allowbreak zero}, \texttt{mass\_\allowbreak law\_\allowbreak underdetermined}, \texttt{both\_\allowbreak laws\_\allowbreak antiparticle\_\allowbreak degenerate}, \texttt{massQuad\_\allowbreak ratio\_\allowbreak not\_\allowbreak constant}\par\smallskip
\noindent\textbf{\texttt{Audit.lean}} (5)\quad \texttt{scalar\_\allowbreak for\_\allowbreak directional\_\allowbreak count}, \texttt{register\_\allowbreak nonempty}, \texttt{more\_\allowbreak corrected\_\allowbreak than\_\allowbreak retracted}, \texttt{register\_\allowbreak grew\_\allowbreak under\_\allowbreak audit}, \texttt{clean\_\allowbreak count}\par\smallskip

\normalsize

\end{document}
